% Leçon 27 — Vocabulaire et compréhension de texte : corruption — Chinois 1ère A — Cours signé OnBuch+
% Programme officiel MINESEC. Compiler avec : tectonic ZHA27-vocabulaire-corruption.tex
\newcommand{\DOCMATIERE}{Chinois}
\newcommand{\DOCNIVEAU}{1ère A}
\newcommand{\DOCMODULE}{Module 3 : Citoyenneté et ouverture au monde}
\newcommand{\DOCLECON}{ZHA27}
\newcommand{\DOCTITRE}{Leçon 27 — Vocabulaire et compréhension de texte : corruption}
% OnBuch+ — préambule « cours signés » ENRICHI (compilé avec tectonic / XeLaTeX)
% Système visuel « Pop » d'OnBuch+ : fond crème (jamais blanc), bordures encre
% épaisses, ombres « dures » décalées, titres Archivo Black en capitales.
% Version MATHÉMATIQUES : graphes (pgfplots/tikz), boîtes pédagogiques
% (définition, propriété, méthode, attention, le savais-tu, exercice résolu,
% exercice, TP, bilan), page de garde et signature OnBuch+ ancrée en bas.
%
% Macros attendues dans main.tex AVANT \input{preamble} :
%   \DOCMATIERE  \DOCNIVEAU  \DOCMODULE  \DOCLECON (numéro)  \DOCTITRE
\documentclass[11pt]{article}

\usepackage{fontspec}
\usepackage[french]{babel} % français = langue principale ; le chinois via \zh{..} / textebox (xeCJK)
\usepackage{xeCJK}
\usepackage{pifont} % \ding{51} \ding{55} utilisés par les modèles
\usepackage[a4paper,margin=2cm,top=3.4cm,bottom=2.8cm,headheight=1.4cm,footskip=1.3cm]{geometry}
\usepackage{amsmath,amssymb,amsfonts}
\usepackage{mathtools} % \xmapsto, \coloneqq... souvent utilisés par les modèles
\usepackage{cancel} % simplifications barrées (bilans de forces)
% \abs{z} (module, valeur absolue) et \norm{u} (norme), utilisés par les modèles
\providecommand{\abs}{}\let\abs\relax\DeclarePairedDelimiter{\abs}{\lvert}{\rvert}
\providecommand{\norm}{}\let\norm\relax\DeclarePairedDelimiter{\norm}{\lVert}{\rVert}
\usepackage[table]{xcolor} % [table] : fournit aussi \rowcolors (alternance de lignes)
\usepackage{pagecolor}
\usepackage{tikz}
\usetikzlibrary{arrows.meta,positioning,calc,shapes.geometric,shapes.misc,shapes.arrows,decorations.pathmorphing,decorations.pathreplacing,decorations.markings,patterns,shadows,mindmap,trees,fit,backgrounds,matrix,intersections,angles,quotes}
\usepackage{pgfplots}
\usepackage[version=4]{mhchem} % équations nucléaires \ce{^{235}_{92}U}
\usepackage[european,siunitx]{circuitikz} % schémas électriques
\pgfplotsset{compat=1.18}
\usepgfplotslibrary{fillbetween,groupplots} % aires entre courbes (calcul intégral)
\usepackage{enumitem}
\usepackage{fancyhdr}
% [most] charge toutes les bibliothèques tcolorbox (listings, minted, raster,
% documentation...) et gonfle inutilement la mémoire TeX sur un document long
% (~300 boîtes) : on ne charge que ce qui est réellement utilisé.
\usepackage[skins,breakable]{tcolorbox}
\usepackage{graphicx}
\usepackage{colortbl}
\usepackage{booktabs}
\usepackage{tabularx}
\usepackage{diagbox} % case d'en-tête diagonale des tableaux à double entrée
% Colonnes extensibles centrée / gauche / droite (C, L, R), souvent utilisées par les modèles
\newcolumntype{C}{>{\centering\arraybackslash}X}
\newcolumntype{L}{>{\raggedright\arraybackslash}X}
\newcolumntype{R}{>{\raggedleft\arraybackslash}X}
\usepackage{array}
\usepackage{multicol}
\usepackage{multirow}
\usepackage{siunitx}
% parse-numbers=false : accepte aussi \SI{2,5 \times 10^{-3}}{...}
\sisetup{locale=FR,output-decimal-marker={,},group-minimum-digits=4,parse-numbers=false}
\DeclareSIUnit\ohm{\ensuremath{\symup{\Omega}}} % U+2126 absent de la police
\DeclareSIUnit\division{div} % réglages d'oscilloscope (ms/div, V/div)
\DeclareSIUnit\tr{tr} % tours (vitesse de rotation, ex. \SI{1500}{\tr\per\minute})
\DeclareSIUnit\molar{M} % concentration molaire (ex. \SI{3}{\molar}, \SI{1}{\micro\molar})
\DeclareSIUnit\an{an} % année (le modèle confond parfois avec \year, un compteur TeX) — ex. \SI{5}{\kilo\meter\per\an}
\DeclareSIUnit\FCFA{FCFA} % franc CFA (ex. \SI{500}{\FCFA\per\kilo\gram})
\DeclareSIUnit\degreeBrix{\textdegree Brix} % teneur en sucre d'un sirop/jus (réfractométrie)
\DeclareSIUnit\month{mois} % le modèle utilise parfois \month, un compteur TeX, comme unité de durée
\usepackage{qrcode}
% \degree : commande fréquemment utilisée par les modèles pour un angle ou une
% température en dehors de \SI{}{} (non fournie par les paquets chargés).
\providecommand{\degree}{^{\circ}}
% \qed (fin de démonstration), utilisé par les modèles sans amsthm
\providecommand{\qed}{\hfill$\square$}
% \barre : double barre des valeurs interdites dans un tableau de variations (tkz-tab)
\providecommand{\barre}{\ensuremath{\|}}
% environnement proof (amsthm non chargé) utilisé par les modèles
\ifdefined\proof\else\newenvironment{proof}[1][Démonstration]{\par\noindent\textit{#1.}\ }{\qed\par}\fi
\usepackage{lastpage}
\usepackage{hyperref}
\hypersetup{hidelinks}

% ============================================================
% Palette « Pop » OnBuch+ — mêmes hex que la classe P (plus_ui.dart)
% ============================================================
\definecolor{poporange}{HTML}{F59321}
\definecolor{popdark}{HTML}{E07A0C}
\definecolor{popcream}{HTML}{FFF6E8}
\definecolor{popink}{HTML}{1C1714}
\definecolor{poppurple}{HTML}{7A5AE0}
\definecolor{popgreen}{HTML}{2E9E6A}
\definecolor{poppink}{HTML}{E0457B}
\definecolor{popblue}{HTML}{2F7DE1}
\definecolor{popgold}{HTML}{C98A12}
\definecolor{popmuted}{HTML}{8A7F76}
\definecolor{popline}{HTML}{EADFD2}
\definecolor{popgray}{HTML}{8A7F76}
\colorlet{popgrey}{popgray}
% Alias tolérants (le modèle nomme parfois les couleurs différemment)
\colorlet{popwhite}{white}
\colorlet{popblack}{popink}
\colorlet{popred}{poppink}
\colorlet{popyellow}{popgold}
% Teintes claires (fonds de tableaux/encadrés) — le modèle nomme parfois
% les couleurs avec un suffixe L (ex. popblueL) sans les avoir définies.
\colorlet{poporangeL}{poporange!20}
\colorlet{popdarkL}{popdark!20}
\colorlet{poppurpleL}{poppurple!20}
\colorlet{popgreenL}{popgreen!20}
\colorlet{poppinkL}{poppink!20}
\colorlet{popblueL}{popblue!20}
\colorlet{popgoldL}{popgold!20}
\colorlet{popredL}{popred!20}
\colorlet{popgrayL}{popgray!20}
% Teintes claires pour fonds de tableaux / zones
\colorlet{popblueL}{popblue!10!white}
\colorlet{poporangeL}{poporange!14!white}
\colorlet{popgreenL}{popgreen!12!white}
\colorlet{poppurpleL}{poppurple!12!white}
\colorlet{poppinkL}{poppink!12!white}
\colorlet{popgoldL}{popgold!14!white}

\pagecolor{popcream}

% ============================================================
% Typographie
% ============================================================
\defaultfontfeatures{Ligatures=TeX}
\setmainfont{PlusJakartaSans-Medium.ttf}[Path=../../../fonts/,BoldFont=PlusJakartaSans-Bold.ttf]
% Symboles, API et emojis absents de la police principale : repli sur DejaVu Sans (caractères actifs)
\newfontface\symfont{DejaVuSans.ttf}[Path=../../../fonts/]
\makeatletter
\newcommand\symactive[1]{\catcode#1=13 \begingroup\lccode`\~=#1 \lowercase{\endgroup\def~}{{\symfont\char#1}}}
\newcommand\symblank[1]{\catcode#1=13 \begingroup\lccode`\~=#1 \lowercase{\endgroup\def~}{}}
\newcommand\symhyph[1]{\catcode#1=13 \begingroup\lccode`\~=#1 \lowercase{\endgroup\def~}{\mbox{-}}}
\newcount\symc
\newcommand\symrange[2]{\symc=#1 \@whilenum\symc<#2 \do{\expandafter\symactive\expandafter{\the\symc}\advance\symc by 1 }}
\newcommand\symblankrange[2]{\symc=#1 \@whilenum\symc<#2 \do{\expandafter\symblank\expandafter{\the\symc}\advance\symc by 1 }}
\makeatother
\symrange{"0250}{"02FF}\symactive{"2713}\symactive{"2714}\symactive{"2717}\symactive{"2718}\symactive{"2610}\symactive{"2611}\symactive{"2612}
\symrange{"2600}{"27BF}
\catcode"2705=13 \begingroup\lccode`\~="2705 \lowercase{\endgroup\def~}{{\symfont\char"2713}}\symblank{"26BD}\symblank{"26A1}
\symrange{"2460}{"257F}\symactive{"223C}\symactive{"2248}\symactive{"2260}
\symactive{"01F8}\symactive{"01F9}\symactive{"1E3E}\symactive{"1E3F}
\symhyph{"2011}\symhyph{"2E1A}\symblankrange{"1F300}{"1FAFF}\symblank{"FE0F}
% Caractères chinois : WenQuanYi Zen Hei (sous-ensemble GB2312 + pinyin), gras/italique simulés
\setCJKmainfont{WenQuanYiZenHei-subset.ttf}[Path=../../../fonts/,AutoFakeBold=3,AutoFakeSlant=0.2]
\xeCJKsetup{CJKspace=false}
% Pinyin : voyelles à caron (ǎ ǐ ǒ ǔ ǖ ǘ ǚ ǜ) absentes de la police principale -> caractères actifs
\newfontface\pyfont{WenQuanYiZenHei-subset.ttf}[Path=../../../fonts/]
\newcommand\pyactive[1]{\catcode#1=13 \begingroup\lccode`\~=#1 \lowercase{\endgroup\def~}{{\pyfont\char#1}}}
\pyactive{"01CD}\pyactive{"01CE}\pyactive{"01CF}\pyactive{"01D0}\pyactive{"01D1}\pyactive{"01D2}\pyactive{"01D3}\pyactive{"01D4}\pyactive{"01D5}\pyactive{"01D6}\pyactive{"01D7}\pyactive{"01D8}\pyactive{"01D9}\pyactive{"01DA}\pyactive{"01DB}\pyactive{"01DC}
% Maths en sans-serif (Fira Math) avec chiffres et lettres droites de Plus
% Jakarta, pour que les formules se fondent dans le texte.
\usepackage[math-style=ISO,bold-style=ISO]{unicode-math}
\setmathfont{FiraMath-Regular.otf}[Path=../../../fonts/]
\setmathfont{PlusJakartaSans-Medium.ttf}[Path=../../../fonts/,range={up/num,up/{Latin,latin},\mathup}]
\usepackage{icomma} % virgule décimale sans espace en mode math
% Caractères mathématiques Unicode absents des polices de texte (Plus Jakarta,
% Archivo Black) : rendus par leur équivalent mathématique, en texte comme en
% formule — sinon ils s'affichent en carré vide (ex. « Arithmétique dans ℤ »).
\usepackage{newunicodechar}
\newunicodechar{ℝ}{\ensuremath{\mathbb{R}}}
\newunicodechar{ℤ}{\ensuremath{\mathbb{Z}}}
\newunicodechar{ℂ}{\ensuremath{\mathbb{C}}}
\newunicodechar{ℕ}{\ensuremath{\mathbb{N}}}
\newunicodechar{ℚ}{\ensuremath{\mathbb{Q}}}
\newunicodechar{𝔻}{\ensuremath{\mathbb{D}}}
\newunicodechar{α}{\ensuremath{\alpha}}
\newunicodechar{β}{\ensuremath{\beta}}
\newunicodechar{γ}{\ensuremath{\gamma}}
\newunicodechar{δ}{\ensuremath{\delta}}
\newunicodechar{ε}{\ensuremath{\varepsilon}}
\newunicodechar{θ}{\ensuremath{\theta}}
\newunicodechar{λ}{\ensuremath{\lambda}}
\newunicodechar{μ}{\ensuremath{\mu}}
\newunicodechar{σ}{\ensuremath{\sigma}}
\newunicodechar{τ}{\ensuremath{\tau}}
\newunicodechar{φ}{\ensuremath{\varphi}}
\newunicodechar{ψ}{\ensuremath{\psi}}
\newunicodechar{ω}{\ensuremath{\omega}}
\newunicodechar{Σ}{\ensuremath{\Sigma}}
% majuscules produites par \MakeUppercase dans les titres (α-aminés -> Α)
\newunicodechar{Α}{\ensuremath{\alpha}}
\newunicodechar{Β}{\ensuremath{\beta}}
\newunicodechar{Γ}{\ensuremath{\Gamma}}
\newunicodechar{Δ}{\ensuremath{\Delta}}
\newunicodechar{Ε}{\ensuremath{\varepsilon}}
\newunicodechar{Θ}{\ensuremath{\Theta}}
\newunicodechar{Λ}{\ensuremath{\Lambda}}
\newunicodechar{Μ}{\ensuremath{\mu}}
\newunicodechar{Τ}{\ensuremath{\tau}}
\newunicodechar{Φ}{\ensuremath{\Phi}}
\newunicodechar{Ψ}{\ensuremath{\Psi}}
\newunicodechar{Ω}{\ensuremath{\Omega}}
\newunicodechar{↦}{\ensuremath{\mapsto}}
\newunicodechar{∈}{\ensuremath{\in}}
\newunicodechar{∉}{\ensuremath{\notin}}
\newunicodechar{∧}{\ensuremath{\wedge}}
\newunicodechar{⇔}{\ensuremath{\Leftrightarrow}}
\newunicodechar{⟺}{\ensuremath{\Longleftrightarrow}}
\newunicodechar{⇒}{\ensuremath{\Rightarrow}}
\newunicodechar{⟹}{\ensuremath{\Longrightarrow}}
\newunicodechar{∘}{\ensuremath{\circ}}
\newunicodechar{∪}{\ensuremath{\cup}}
\newunicodechar{∩}{\ensuremath{\cap}}
\newunicodechar{≡}{\ensuremath{\equiv}}
\newunicodechar{⊂}{\ensuremath{\subset}}
\newunicodechar{⊥}{\ensuremath{\perp}}
\newunicodechar{∃}{\ensuremath{\exists}}
\newunicodechar{∀}{\ensuremath{\forall}}
\newunicodechar{∛}{\ensuremath{\sqrt[3]{\,}}}
\newunicodechar{∜}{\ensuremath{\sqrt[4]{\,}}}
\newunicodechar{⃗}{\ensuremath{{}^{\to}}}
\newunicodechar{ⁿ}{\ifmmode^{n}\else\textsuperscript{\ensuremath{n}}\fi}
\newunicodechar{ˣ}{\ifmmode^{x}\else\textsuperscript{\ensuremath{x}}\fi}
\newunicodechar{ᵇ}{\ifmmode^{b}\else\textsuperscript{\ensuremath{b}}\fi}
\newunicodechar{ʳ}{\ifmmode^{r}\else\textsuperscript{\ensuremath{r}}\fi}
\newunicodechar{ᵘ}{\ifmmode^{u}\else\textsuperscript{\ensuremath{u}}\fi}
\newunicodechar{ʸ}{\ifmmode^{y}\else\textsuperscript{\ensuremath{y}}\fi}
\newunicodechar{ᵃ}{\ifmmode^{a}\else\textsuperscript{\ensuremath{a}}\fi}
\newunicodechar{ᵉ}{\ifmmode^{e}\else\textsuperscript{\ensuremath{e}}\fi}
\newunicodechar{ᵏ}{\ifmmode^{k}\else\textsuperscript{\ensuremath{k}}\fi}
\newunicodechar{ᵖ}{\ifmmode^{p}\else\textsuperscript{\ensuremath{p}}\fi}
\newunicodechar{ᴺ}{\ifmmode^{N}\else\textsuperscript{\ensuremath{N}}\fi}
\newunicodechar{ᶜ}{\ifmmode^{c}\else\textsuperscript{\ensuremath{c}}\fi}
\newunicodechar{ʰ}{\ifmmode^{h}\else\textsuperscript{\ensuremath{h}}\fi}
\newunicodechar{ᵠ}{\ifmmode^{q}\else\textsuperscript{\ensuremath{q}}\fi}
\newunicodechar{ₙ}{\ifmmode_{n}\else\textsubscript{\ensuremath{n}}\fi}
\newunicodechar{ₐ}{\ifmmode_{a}\else\textsubscript{\ensuremath{a}}\fi}
\newunicodechar{ᵢ}{\ifmmode_{i}\else\textsubscript{\ensuremath{i}}\fi}
\newunicodechar{ᵣ}{\ifmmode_{r}\else\textsubscript{\ensuremath{r}}\fi}
\newunicodechar{ₖ}{\ifmmode_{k}\else\textsubscript{\ensuremath{k}}\fi}
\newunicodechar{ᵦ}{\ifmmode_{\beta}\else\textsubscript{\ensuremath{\beta}}\fi}
\newunicodechar{ₓ}{\ifmmode_{x}\else\textsubscript{\ensuremath{x}}\fi}
\newfontfamily\popdisplay{ArchivoBlack-Regular.ttf}[Path=../../../fonts/]
\newfontfamily\poplogo{SpaceGrotesk-Bold.ttf}[Path=../../../fonts/]
\newfontfamily\popsemi{PlusJakartaSans-SemiBold.ttf}[Path=../../../fonts/]
\newfontfamily\popxbold{PlusJakartaSans-ExtraBold.ttf}[Path=../../../fonts/]
\newfontfamily\popxboldls{PlusJakartaSans-ExtraBold.ttf}[Path=../../../fonts/,LetterSpace=3.0]

\setlist[enumerate]{leftmargin=1.7em,itemsep=5pt,topsep=4pt}
\setlist[itemize]{leftmargin=1.7em,itemsep=4pt,topsep=4pt,label={\color{poporange}\textbullet}}
\setlength{\parindent}{0pt}
\setlength{\parskip}{5pt}
\renewcommand{\arraystretch}{1.25}

% pgfplots : style Pop par défaut
\pgfplotsset{
  every axis/.append style={axis line style={popink,line width=1pt},tick style={popink},
    grid style={popline,line width=0.5pt},label style={font=\small},tick label style={font=\footnotesize}},
  popcurve/.style={poporange,line width=1.8pt,no markers,smooth},
}
% Même style disponible dans un \draw TikZ simple (hors pgfplots)
\tikzset{popcurve/.style={poporange,line width=1.8pt}}
\tikzset{popgrid/.style={popline,very thin}}
\colorlet{popcurve}{poporange} % le nom du style est parfois utilisé comme couleur

% ============================================================
% Pill / squiggle
% ============================================================
\newcommand{\poppill}[3]{%
  \tikz[baseline=-0.55ex]\node[draw=popink,line width=1.2pt,rounded corners=2.6pt,%
    fill=#2,inner xsep=6.5pt,inner ysep=3pt]{%
    \popxboldls\fontsize{7}{7}\selectfont\color{#3}\MakeUppercase{#1}};%
}
\newcommand{\popsquiggle}[1]{%
  \tikz[baseline=-0.4ex]{
    \draw[#1,line width=2.6pt,line cap=round]
      plot [smooth] coordinates {(0,0.09) (0.34,0.01) (0.68,0.21) (1.02,0.01) (1.36,0.10)};
  }%
}
% Mot-clé surligné (marqueur orange sous le texte)
\newcommand{\cle}[1]{{\popxbold\color{popdark}#1}}

% ============================================================
% En-tête / pied
% ============================================================
\pagestyle{fancy}
\fancyhf{}
\renewcommand{\headrulewidth}{0pt}
\renewcommand{\footrulewidth}{0pt}
\lhead{\raisebox{-0.15ex}{\poplogo\fontsize{13}{13}\selectfont\color{popink}OnBuch\color{poporange}+}}
\rhead{\poppill{\DOCMATIERE\ \textbullet\ \DOCNIVEAU\ \textbullet\ Leçon \DOCLECON}{white}{popink}}
\lfoot{\popxbold\fontsize{7.6}{7.6}\selectfont\color{popmuted}\MakeUppercase{Cours signé OnBuch+}\ \textbullet\ {\popsemi\DOCTITRE}}
\rfoot{\tikz[baseline=-0.6ex]\node[rounded corners=8.5pt,fill=popink,minimum height=17pt,minimum width=17pt,inner xsep=5pt,inner sep=0]{\popxbold\fontsize{8}{8}\selectfont\color{popcream}\thepage\,/\,\pageref*{LastPage}};}

% ============================================================
% Titres
% ============================================================
\newcommand{\chaptitre}[2]{%
  \begin{tcolorbox}[enhanced,colback=poporange,colframe=popink,boxrule=2.6pt,arc=11pt,
      drop shadow=popink,left=15pt,right=15pt,top=12pt,bottom=11pt,before skip=3pt,after skip=18pt]
    \poppill{#2}{popink}{popcream}\par\vspace{9pt}
    {\raggedright\hyphenpenalty=10000\exhyphenpenalty=10000\popdisplay\fontsize{22}{24}\selectfont\color{popcream}\MakeUppercase{#1}\par}
  \end{tcolorbox}
}
% Section numérotée : pastille orange avec le numéro + titre Archivo
\newcounter{popsec}
\newcommand{\coursec}[1]{%
  \stepcounter{popsec}\setcounter{popsub}{0}%
  \par\vspace{16pt}\noindent
  \begin{minipage}{\linewidth}
  \tikz[baseline=-0.7ex]\node[draw=popink,line width=1.6pt,fill=poporange,rounded corners=4pt,minimum size=22pt,inner sep=0]{\popdisplay\fontsize{12}{12}\selectfont\color{popcream}\thepopsec};%
  \hspace{8pt}{\popdisplay\fontsize{15}{17}\selectfont\color{popink}\MakeUppercase{#1}}\par
  \vspace{4pt}\noindent{\color{poporange}\rule{36pt}{3.6pt}}
  \end{minipage}\par\nopagebreak\vspace{8pt}
}
\newcounter{popsub}
\newcommand{\courssub}[1]{%
  \stepcounter{popsub}%
  \par\vspace{9pt}\noindent{\popxbold\fontsize{11.5}{13}\selectfont\color{popink}{\color{poporange}\thepopsec.\thepopsub}\hspace{6pt}#1}\par\nopagebreak\vspace{4pt}
}

% ============================================================
% Boîtes pédagogiques — même recette : fond blanc, bordure épaisse colorée,
% ombre dure encre, badge plein. Titre optionnel [#1].
% ============================================================
\tcbset{popbase/.style={breakable,enhanced,colback=white,boxrule=2.3pt,arc=9pt,
  shadow={3pt}{-3pt}{0pt}{popink},left=14pt,right=14pt,top=11pt,bottom=11pt,before skip=11pt,after skip=15pt,
  fontupper=\popsemi\fontsize{10.6}{14.5}\selectfont\color{popink},
  fontlower=\popsemi\fontsize{10.6}{14.5}\selectfont\color{popink}}}
% Coche dessinée (le glyphe ✓ n'existe pas dans les polices du document)
\newcommand{\popcheck}[1]{\tikz[baseline=-0.5ex]\draw[#1,line width=1.6pt,line cap=round,line join=round](-0.13,0.0)--(-0.04,-0.09)--(0.13,0.1);}
\providecommand{\coche}{\popcheck{popgreen}}
\providecommand{\resultat}[1]{\par\smallskip\noindent\fcolorbox{popgreen}{popgreenL}{\parbox{\dimexpr\linewidth-2\fboxsep-2\fboxrule\relax}{\textbf{Résultat.} #1}}\par\smallskip}
\newcommand{\popboxhead}[3]{\poppill{#1}{#2}{white}\ifx\relax#3\relax\else\hspace{6pt}{\popxbold\fontsize{10.6}{12}\selectfont\color{popink}#3}\fi\par\vspace{7pt}}

\newtcolorbox{aretenir}[1][]{popbase,colframe=popgreen,before upper={\popboxhead{À retenir}{popgreen}{#1}}}
% Texte en chinois (document de lecture, dialogue, extrait) : cadre
% d'étude. Un titre optionnel (source, auteur) s'affiche dans l'en-tête.
\newtcolorbox{textebox}[1][]{popbase,colframe=popblue,before upper={\popboxhead{文本 · Texte}{popblue}{#1}},after upper={}}
% Marques ✓ / ✗ (forme correcte / fautive) souvent utilisées par les modèles
\providecommand{\cross}{\textcolor{poppink}{\ensuremath{\times}}}
\providecommand{\xmark}{\cross}\providecommand{\crossmark}{\cross}\providecommand{\cmark}{\ensuremath{\checkmark}}
% Ligne de réponse / justification à compléter (inventée par les modèles)
\providecommand{\justification}[1]{\par\noindent\textit{Justification~:} \dotfill\par}
% \textipa (package tipa non chargé) : la police ne couvre pas l'API -> simple passage
\providecommand{\textipa}[1]{#1}
% Soulignement (ulem non chargé) : \uline, \ul
\providecommand{\uline}[1]{\underline{#1}}\providecommand{\ul}[1]{\underline{#1}}
% \glossaire{\item ...} inventée par les modèles : liste de glossaire
\providecommand{\glossaire}[1]{\begin{itemize}#1\end{itemize}}
% \alert (beamer) inventée par les modèles : texte mis en évidence
\providecommand{\alert}[1]{\textbf{\textcolor{poppink}{#1}}}
% variantes ASCII d'ordinaux écrites en macro
\providecommand{\iere}{\textsuperscript{re}}\providecommand{\ieme}{\textsuperscript{e}}\providecommand{\ere}{\textsuperscript{re}}\providecommand{\eme}{\textsuperscript{e}}\providecommand{\er}{\textsuperscript{er}}\providecommand{\ie}{\textsuperscript{e}}
% Ordinaux français écrits en macro par les modèles : 1\ière, 3\ième
\newcommand{\ière}{\textsuperscript{re}}\newcommand{\ième}{\textsuperscript{e}}
% \exemple (inventée par les modèles) : étiquette « Exemple : »
\providecommand{\exemple}{\textbf{Exemple~:}\ }
% \square utilisable aussi hors mode math (cases à cocher)
\AtBeginDocument{\let\squareMath\square\renewcommand{\square}{\ensuremath{\squareMath}}}
% Mot ou phrase en chinois à l'intérieur d'un paragraphe français
\newcommand{\zh}[1]{\ifmmode\text{#1}\else{#1}\fi}
\let\eng\zh \let\esp\zh \let\all\zh % alias tolérés
% flèches d'intonation inventées par les modèles
\providecommand{\textsearrow}{}\renewcommand{\textsearrow}{\ensuremath{\searrow}}\providecommand{\textnearrow}{}\renewcommand{\textnearrow}{\ensuremath{\nearrow}}
% \fr{..} : mot français (inventé par les modèles)
\providecommand{\fr}[1]{\textit{#1}}
% zhbox : encadré de texte chinois inventé par les modèles
\newenvironment{zhbox}{\begin{quote}}{\end{quote}}
% \courssubsub : titre de niveau 3 inventé par les modèles
\providecommand{\courssubsub}[1]{\par\medskip\noindent\textbf{#1}\par\smallskip}
% \zhbf : texte chinois en gras inventé par les modèles
\providecommand{\zhbf}[1]{\textbf{#1}}
% \vs : « versus » inventé par les modèles
\providecommand{\vs}{\textit{vs}\ }
% \blank : case à compléter inventée par les modèles
\providecommand{\blank}[1][]{\underline{\hspace{1.6em}}}
\newtcolorbox{exemplebox}[1][]{popbase,colframe=poporange,before upper={\popboxhead{Exemple}{poporange}{#1}}}
\newtcolorbox{definition}[1][]{popbase,colframe=popblue,before upper={\popboxhead{Définition}{popblue}{#1}}}
\newtcolorbox{propriete}[1][]{popbase,colframe=poppurple,before upper={\popboxhead{Propriété}{poppurple}{#1}}}
\newtcolorbox{methode}[1][]{popbase,colframe=popgold,before upper={\popboxhead{Méthode}{popgold}{#1}}}
\newtcolorbox{attention}[1][]{popbase,colframe=poppink,before upper={\popboxhead{Attention}{poppink}{#1}}}
\newtcolorbox{experience}[1][]{popbase,colframe=popblue,colback=popblueL,before upper={\popboxhead{Expérience}{popblue}{#1}}}
% « Le savais-tu ? » : carte violette pleine (contexte, culture, Cameroun)
\newtcolorbox{savaistu}[1][]{popbase,colback=poppurple,colframe=popink,
  fontupper=\popsemi\fontsize{10.4}{14.2}\selectfont\color{white},
  before upper={\poppill{Le savais-tu\,?}{white}{poppurple}\ifx\relax#1\relax\else\hspace{6pt}{\popxbold\color{white}#1}\fi\par\vspace{7pt}}}
% Exercice résolu : énoncé en haut, \tcblower puis la solution sur fond crème
\newcounter{popexo}\newcounter{popexr}
\newtcolorbox{exoresolu}[1][]{popbase,colframe=poporange,colbacklower=popcream,
  segmentation style={popink,line width=1.2pt,dashed},
  before upper={\stepcounter{popexr}\popboxhead{Exercice résolu \thepopexr}{poporange}{#1}},
  before lower={\poppill{Solution}{popink}{popcream}\par\vspace{6pt}}}
% Exercice (non résolu en place) — la correction est dans la partie Corrigés
\newtcolorbox{exercice}[1][]{popbase,colframe=popink,
  before upper={\stepcounter{popexo}\popboxhead{Exercice \thepopexo}{popink}{#1}}}
% Corrigé
\newtcolorbox{corrige}[1][]{popbase,colframe=popgreen,colback=popgreenL,
  before upper={\popboxhead{Corrigé}{popgreen}{#1}}}
% Objectifs / prérequis / bilan
\newtcolorbox{objectifs}{popbase,colframe=popink,colback=poporangeL,
  before upper={\popboxhead{Objectifs de la leçon}{popink}{}}}
\newtcolorbox{prerequis}{popbase,colframe=popink,colback=popblueL,
  before upper={\popboxhead{Prérequis}{popblue}{}}}
\newtcolorbox{bilan}{popbase,colframe=popink,colback=white,boxrule=2.8pt,
  before upper={\popboxhead{Fiche bilan}{poporange}{L'essentiel à connaître pour l'examen}}}
% Figure encadrée avec légende
\newtcolorbox{popfigure}[1][]{enhanced,breakable=false,colback=white,colframe=popline,boxrule=1.4pt,arc=8pt,
  left=10pt,right=10pt,top=10pt,bottom=8pt,before skip=10pt,after skip=12pt,halign=center,
  lower separated=false,halign lower=center,
  fontlower=\popsemi\fontsize{9}{11.5}\selectfont\color{popmuted},
  #1}
\newcommand{\legende}[1]{\par\vspace{6pt}{\popsemi\fontsize{9}{11.5}\selectfont\color{popmuted}#1}}

% ============================================================
% Signature OnBuch+
% ============================================================
\newcommand{\poplogoinline}[1]{{\poplogo\fontsize{#1}{#1}\selectfont\color{popink}OnBuch\color{poporange}+}}
\newcommand{\signatureonbuch}{%
  \begin{tcolorbox}[enhanced,colback=popink,colframe=popink,boxrule=2.2pt,arc=10pt,
      drop shadow=poporange,left=14pt,right=14pt,top=10pt,bottom=10pt,before skip=0pt,after skip=0pt]
    \tikz[baseline=-0.7ex]\node[circle,fill=popgreen,draw=popcream,line width=1.3pt,minimum size=22pt,inner sep=0]{\popcheck{white}};%
    \hspace{9pt}\begin{minipage}[c]{0.62\linewidth}
      {\popxbold\fontsize{12}{13}\selectfont\color{popcream}Signé\ \textbullet\ {\poplogo OnBuch\color{poporange}+}}\par\vspace{2pt}
      {\raggedright\popsemi\fontsize{8}{10.5}\selectfont\color{popline}\MakeUppercase{Équipe pédagogique OnBuch+}\\\MakeUppercase{Conforme au programme officiel MINESEC}\par}
    \end{minipage}\hfill
    {\poplogo\fontsize{22}{22}\selectfont\color{popcream}OnBuch\color{poporange}+}
  \end{tcolorbox}%
}

% ============================================================
% Page de garde
% #1 titre · #2 matière · #3 niveau · #4 module · #5 n° leçon · #6 durée ·
% #7 description / objectifs (texte ou itemize)
% ============================================================
\newcommand{\pagegarde}[7]{%
  \thispagestyle{empty}
  \newgeometry{a4paper,margin=1.7cm,top=1.6cm,bottom=1.4cm}%
  \begin{tikzpicture}[remember picture,overlay]
    \fill[poporange!18] ([xshift=-3.2cm,yshift=-3.6cm]current page.north east) circle (4.6cm);
    \fill[poppurple!14] ([xshift=2.4cm,yshift=7.6cm]current page.south west) circle (3.8cm);
  \end{tikzpicture}%
  {\poplogo\fontsize{30}{30}\selectfont\color{popink}OnBuch\color{poporange}+}\hfill
  \raisebox{0.6ex}{\poppill{Cours signé \textbullet\ Premium}{popink}{popcream}}\par
  \vspace{1pt}\popsquiggle{poppurple}\par
  \vspace{8pt}
  \begin{tcolorbox}[enhanced,colback=poporange,colframe=popink,boxrule=2.8pt,arc=13pt,
      drop shadow=popink,left=18pt,right=18pt,top=12pt,bottom=13pt,after skip=10pt]
    \poppill{#2 \textbullet\ #3}{popink}{popcream}\hspace{5pt}\poppill{#4}{white}{popink}\par\vspace{8pt}
    {\popdisplay\fontsize{14}{14}\selectfont\color{popink}LEÇON #5}\par\vspace{4pt}
    {\raggedright\hyphenpenalty=10000\exhyphenpenalty=10000\popdisplay\fontsize{25}{27}\selectfont\color{popcream}\MakeUppercase{#1}\par}
    \vspace{8pt}
    {\popxbold\fontsize{9.5}{9.5}\selectfont\color{popink}\MakeUppercase{Durée conseillée : #6}}
  \end{tcolorbox}
  \begin{tcolorbox}[enhanced,colback=white,colframe=popink,boxrule=2.2pt,arc=10pt,
      drop shadow=popink,left=16pt,right=16pt,top=10pt,bottom=10pt,after skip=8pt]
    \poppill{Dans cette leçon}{poppurple}{white}\par\vspace{5pt}
    \popsemi\fontsize{9.3}{12.6}\selectfont\color{popink}#7
  \end{tcolorbox}
  \noindent\begin{minipage}[t]{0.487\linewidth}
    \begin{tcolorbox}[enhanced,colback=popgreen,colframe=popink,boxrule=2.2pt,arc=10pt,
        drop shadow=popink,left=11pt,right=11pt,top=10pt,bottom=10pt,height=3.1cm,valign=center]
      \tcbox[colback=white,colframe=popink,boxrule=1.3pt,arc=5pt,boxsep=0pt,nobeforeafter,
        left=3pt,right=3pt,top=3pt,bottom=3pt,valign=center]{\qrcode[height=1.9cm]{https://play.google.com/store/apps/details?id=com.luvvix.onbuch&hl=fr}}%
      \hspace{8pt}\begin{minipage}[c]{0.5\linewidth}
        {\popdisplay\fontsize{11}{12.5}\selectfont\color{white}\MakeUppercase{Télécharge OnBuch}}\par\vspace{4pt}
        {\raggedright\popsemi\fontsize{8.4}{11}\selectfont\color{white}Gratuit \textbullet\ Google Play\\Tuteur IA, annales, concours, résultats\par}
      \end{minipage}
    \end{tcolorbox}
  \end{minipage}\hfill
  \begin{minipage}[t]{0.487\linewidth}
    \begin{tcolorbox}[enhanced,colback=poppurple,colframe=popink,boxrule=2.2pt,arc=10pt,
        drop shadow=popink,left=11pt,right=11pt,top=10pt,bottom=10pt,height=3.1cm,valign=center]
      \tikz[baseline=-0.7ex]\node[circle,fill=white,draw=popink,line width=1.3pt,minimum size=1.6cm,inner sep=0]{\popdisplay\fontsize{24}{24}\selectfont\color{poppurple}+};%
      \hspace{8pt}\begin{minipage}[c]{0.6\linewidth}
        {\popdisplay\fontsize{11}{12.5}\selectfont\color{white}\MakeUppercase{Passe à OnBuch+}}\par\vspace{4pt}
        {\raggedright\popsemi\fontsize{8.4}{11}\selectfont\color{white}Tous les cours signés, Tuteur IA illimité, fiches PDF — onglet OnBuch+ de l'app\par}
      \end{minipage}
    \end{tcolorbox}
  \end{minipage}
  \vspace{8pt}
  \popcontenu
  \vfill
  \signatureonbuch
  \restoregeometry
  \newpage
}

% Rangée « Ce cours contient » (page de garde)
\newcommand{\poptile}[3]{% #1 couleur #2 gros texte #3 légende
  \begin{tcolorbox}[enhanced,colback=#1,colframe=popink,boxrule=2pt,arc=9pt,shadow={2.5pt}{-2.5pt}{0pt}{popink},
      left=6pt,right=6pt,top=9pt,bottom=9pt,halign=center,valign=center,height=1.75cm,width=\linewidth,nobeforeafter]
    {\popdisplay\fontsize{12.5}{13}\selectfont\color{white}\MakeUppercase{#2}}\par\vspace{3pt}
    {\centering\popsemi\fontsize{7.8}{9.5}\selectfont\color{white}#3\par}
  \end{tcolorbox}}
\newcommand{\popcontenu}{%
  \par\noindent{\popxboldls\fontsize{7.5}{7.5}\selectfont\color{popmuted}CE COURS SIGNÉ CONTIENT}\par\vspace{7pt}
  \noindent\begin{minipage}[t]{0.235\linewidth}\poptile{popblue}{Cours}{complet, illustré, pas à pas}\end{minipage}\hfill
  \begin{minipage}[t]{0.235\linewidth}\poptile{poporange}{Méthodes}{et exercices résolus}\end{minipage}\hfill
  \begin{minipage}[t]{0.235\linewidth}\poptile{poppink}{Exercices}{type examen + corrigés}\end{minipage}\hfill
  \begin{minipage}[t]{0.235\linewidth}\poptile{popgreen}{Fiche bilan}{l'essentiel à retenir}\end{minipage}\par}

% Sommaire
\newtcolorbox{sommaire}{enhanced,colback=white,colframe=popink,boxrule=2.2pt,arc=10pt,
  drop shadow=popink,left=18pt,right=18pt,top=13pt,bottom=13pt,before skip=6pt,after skip=20pt,
  before upper={\poppill{Sommaire}{poppurple}{white}\par\vspace{9pt}\popsemi\fontsize{10.6}{14.5}\selectfont\color{popink}}}

% Signature de fin de cours — poussée tout en bas de la dernière page
\newcommand{\findecours}{%
  \par\vfill\nopagebreak
  \begin{center}\popsquiggle{poporange}\end{center}\vspace{4pt}
  \signatureonbuch
}
\AtBeginDocument{\def\llbracket{\mathopen{[\mkern-3.5mu[}}\def\rrbracket{\mathclose{]\mkern-3.5mu]}}} % intervalles d'entiers (glyphe ⟦ absent de la police)
% Glyphes absents de Fira Math : redéfinis à partir de glyphes disponibles
\AtBeginDocument{%
  \def\setminus{\mathbin{\backslash}}\let\smallsetminus\setminus
  \def\vdots{\mathord{\vbox{\baselineskip3.2pt\lineskiplimit0pt\kern4pt\hbox{.}\hbox{.}\hbox{.}}}}%
  \def\ddots{\mathinner{\mkern1mu\raise6.5pt\hbox{.}\mkern2mu\raise3.6pt\hbox{.}\mkern2mu\raise0.7pt\hbox{.}\mkern1mu}}%
  \def\triangle{\mathord{\scalebox{0.9}{$\symup{\Delta}$}}}%
  \def\nabla{\mathord{\raisebox{\depth}{\scalebox{1}[-1]{$\symup{\Delta}$}}}}%
  \def\checkmark{\popcheck{popdark}}%
  \def\star{\mathbin{\ast}}\def\bigstar{\mathord{\ast}}%
  \def\diamond{\mathbin{\scalebox{0.6}{$\Diamond$}}}%
  \def\bigcup{\mathop{\mathchoice{\raisebox{-0.3em}{\scalebox{1.7}{$\cup$}}}{\scalebox{1.25}{$\cup$}}{\cup}{\cup}}\displaylimits}%
  \def\bigcap{\mathop{\mathchoice{\raisebox{-0.3em}{\scalebox{1.7}{$\cap$}}}{\scalebox{1.25}{$\cap$}}{\cap}{\cap}}\displaylimits}%
  \def\lesseqqgtr{\mathrel{\lessgtr}}\def\bigoplus{\mathop{\oplus}\displaylimits}%
}
\providecommand{\ssi}{si et seulement si} % abréviation fréquente
\AtBeginDocument{\providecommand{\wideparen}{\overparen}} % arc de cercle

\begin{document}

\pagegarde{Leçon 27 — Vocabulaire et compréhension de texte : corruption}{Chinois}{1ère A}{Module 3 : Citoyenneté et ouverture au monde}{ZHA27}{2 h}{Cette leçon de vocabulaire et compréhension de texte porte sur le thème de la corruption. À partir d'un texte original en chinois simplifié, les élèves acquièrent le lexique essentiel (腐败, 贪污, 贿赂, 反腐败...), étudient les clés et l'ordre des traits de caractères clés, puis répondent à des questions de compréhension avant une production guidée.
\vspace{4pt}
\begin{multicols}{2}\small
\begin{itemize}
\item À la fin de cette leçon, je sais lire à voix haute un court texte chinois sur la corruption avec une prononciation correcte (pinyin et tons).
\item À la fin de cette leçon, je sais reconnaître et traduire au moins 15 mots du champ lexical de la corruption.
\item À la fin de cette leçon, je sais identifier la clé (radical) et tracer l'ordre des traits des caractères 腐, 败, 贪, 贿, 廉.
\item À la fin de cette leçon, je sais répondre par écrit à des questions de compréhension sur le texte support.
\item À la fin de cette leçon, je sais employer les collocations fréquentes : 反腐败, 贪污受贿, 廉洁政府, 打击腐败.
\item À la fin de cette leçon, je sais comparer brièvement la lutte contre la corruption au Cameroun et en Chine.
\end{itemize}\end{multicols}}

\begin{sommaire}
\begin{multicols}{2}
\begin{enumerate}
\item Mise en route et découverte du thème
\item Texte support : lecture et compréhension globale
\item Glossaire du thème et prononciation
\item Clés (radicaux) et ordre des traits
\item Questions de compréhension et approfondissement
\item Production guidée et réinvestissement
\item Activité d'intégration
\item Méthodes et exercices résolus
\item Exercices
\item Corrigés des exercices
\item Fiche bilan
\end{enumerate}
\end{multicols}
\end{sommaire}

\begin{objectifs}
\begin{itemize}
\item À la fin de cette leçon, je sais lire à voix haute un court texte chinois sur la corruption avec une prononciation correcte (pinyin et tons).
\item À la fin de cette leçon, je sais reconnaître et traduire au moins 15 mots du champ lexical de la corruption.
\item À la fin de cette leçon, je sais identifier la clé (radical) et tracer l'ordre des traits des caractères 腐, 败, 贪, 贿, 廉.
\item À la fin de cette leçon, je sais répondre par écrit à des questions de compréhension sur le texte support.
\item À la fin de cette leçon, je sais employer les collocations fréquentes : 反腐败, 贪污受贿, 廉洁政府, 打击腐败.
\item À la fin de cette leçon, je sais comparer brièvement la lutte contre la corruption au Cameroun et en Chine.
\item À la fin de cette leçon, je sais rédiger 4 à 5 phrases guidées donnant mon avis sur la corruption.
\end{itemize}
\end{objectifs}

\begin{prerequis}
\begin{itemize}
\item Lecture du pinyin et des quatre tons (Seconde)
\item Structures de base : sujet + verbe + complément ; 是 / 有 (Seconde)
\item Vocabulaire de la vie en société : 政府, 法律, 社会, 人民 (début de 1ère)
\item Notions sur les radicaux et l'ordre des traits (Seconde)
\item Expression de l'opinion simple : 我认为... / 我觉得... (leçons antérieures du module)
\end{itemize}
\end{prerequis}

\newpage

\coursec{Mise en route et découverte du thème}

\courssub{Situation d'accroche et problématique}

Au marché Mokolo de Yaoundé, un commerçant raconte qu'on lui a demandé un « pot-de-vin » pour obtenir un permis de construire. Au Nigeria, on parle de « \emph{dash} », au Ghana de « \emph{brown envelope} ». La corruption est un fléau que la Chine combat aussi avec des campagnes très médiatisées. Comment en parler en chinois ? Cette leçon nous donne les mots pour lire, comprendre et débattre de ce sujet de société.

\begin{textebox}[Amorce lexicale]
腐败是一个严重的问题。\\
Fǔbài shì yí ge yánzhòng de wèntí.\\
La corruption est un problème grave.
\end{textebox}

\courssub{Brainstorming : qu'est-ce que la corruption ?}

En classe entière, l'enseignant écrit au tableau le mot \textbf{corruption} et note les propositions des élèves. On guide la discussion vers des exemples concrets du quotidien camerounais :

\begin{itemize}
    \item \textbf{Pot-de-vin} (donner de l'argent pour obtenir un service dû) : payer un policier pour éviter une contravention, « graisser la patte » d'un agent administratif pour accélérer un dossier.
    \item \textbf{Détournement de fonds publics} : un gestionnaire qui s'approprie l'argent de la cantine scolaire ou du budget d'un hôpital.
    \item \textbf{Favoritisme / Népotisme} : embaucher son cousin ou son ami au lieu du candidat le plus compétent ; attribuer un marché public à une entreprise amie sans appel d'offres.
    \item \textbf{Abus de pouvoir} : un chef de service qui utilise sa fonction pour des intérêts personnels.
\end{itemize}

\begin{attention}[Piège lexical : « corruption » vs « pourriture »]
En français, \emph{corruption} vient du latin \emph{corrumpere} (briser, gâter). En chinois, \zh{腐败} (fǔbài) signifie littéralement « pourri (腐) + gâté/défait (败) ». Le sens est identique : quelque chose qui pourrit de l'intérieur. Attention à ne pas confondre \zh{腐} (fǔ, pourri, comme dans \zh{豆腐} dòufu, tofu « haricot pourri/fermenté ») et \zh{府} (fǔ, mansion, gouvernement, comme dans \zh{政府} zhèngfǔ). Ce sont des homophones de ton différent (3e ton vs 3e ton mais caractères et sens radicalement différents).
\end{attention}

\courssub{Présentation du mot-clé : 腐败 fǔbài}

\begin{definition}[Mot-clé : 腐败 fǔbài]
\textbf{Caractères :} \zh{腐败} \\
\textbf{Pinyin :} fǔbài (3e ton + 4e ton) \\
\textbf{Nature :} Verbe d'état / Nom (verbe adjectival) \\
\textbf{Sens :} Être corrompu ; corruption, pourriture, dépravation. \\
\textbf{Structure interne :} \zh{腐} (radical \zh{月} « chair/lune » + \zh{府} indicateur phonétique) « pourrir, putréfier » + \zh{败} (radical \zh{攵} « frapper/action » + \zh{贝} « coquillage/monnaie ») « être vaincu, gâter, faire faillite ».
\end{definition}

\paragraph{Écriture au tableau et prononciation.}
L'enseignant écrit le caractère en grand, trait par trait, en verbalisant l'ordre (voir section 4). La classe répète en chœur : \textbf{fǔ-bài} (voix qui descend puis chute brutale). On insiste sur le 3e ton de \zh{腐} (descendant-remontant, souvent réalisé « demi-troisième ton » bas en milieu de phrase) et le 4e ton de \zh{败} (chutant net).

\paragraph{Collocations immédiates (à mémoriser).}

\begin{center}
\begin{tabularx}{\linewidth}{|X|X|X|X|}
\hline
\rowcolor{popblueL} \textbf{Caractères} & \textbf{Pinyin} & \textbf{Structure} & \textbf{Traduction} \\
\hline
反腐败 & fǎn fǔbài & Verbe + Objet & Lutter contre la corruption \\
\hline
腐败分子 & fǔbài fènzǐ & Nom + Suffixe & Corrompu, élément corrompu \\
\hline
贪污腐败 & tānwū fǔbài & Verbe + Verbe & Détournement et corruption (binôme fixe) \\
\hline
预防腐败 & yùfáng fǔbài & Verbe + Objet & Prévenir la corruption \\
\hline
\end{tabularx}
\end{center}

\courssub{Prédiction à partir du titre du texte}

Le titre du texte support est : \zh{《反腐败，人人有责》} (Fǎn fǔbài, rénrén yǒu zé).

\begin{methode}[Comment prédire le contenu d'un texte à partir de son titre]
\begin{enumerate}
    \item \textbf{Segmenter le titre :} \zh{反腐败} (Lutter contre la corruption) + \zh{，} + \zh{人人有责} (Chacun a sa part de responsabilité / C'est l'affaire de tous).
    \item \textbf{Identifier les mots connus :} \zh{反} (contre, anti-), \zh{腐败} (corruption), \zh{人} (personne), \zh{有} (avoir), \zh{责} (responsabilité, blâme).
    \item \textbf{Formuler des hypothèses :} Le texte parlera probablement de la lutte anti-corruption, de lois, de sanctions, et insistera sur le rôle du citoyen ordinaire (dénonciation, intégrité personnelle).
    \item \textbf{Noter les questions :} Quelles mesures concrètes ? Quel rôle pour les jeunes / les élèves ?
\end{enumerate}
\end{methode}

\courssub{Première écoute / lecture globale par l'enseignant}

L'enseignant lit le texte à voix normale (ou diffuse l'audio). Les élèves ne cherchent pas à tout comprendre, mais \textbf{repèrent les mots déjà connus} (mots « amis » du HSK 1-2). Ils cochent sur leur fiche ou lèvent la main quand ils entendent :

\begin{center}
\begin{tabularx}{\linewidth}{|X|X|X|}
\hline
\rowcolor{popgreenL} \textbf{Mot repéré} & \textbf{Pinyin} & \textbf{Sens} \\
\hline
政府 & zhèngfǔ & Gouvernement \\
\hline
人民 & rénmín & Peuple \\
\hline
法律 & fǜlǜ & Loi \\
\hline
社会 & shèhuì & Société \\
\hline
问题 & wèntí & Problème \\
\hline
重要 & zhòngyào & Important \\
\hline
\end{tabularx}
\end{center}

\begin{exemplebox}[Modélisation par l'enseignant]
L'enseignant : « J'entends \zh{政府} (zhèngfǔ) au début de la deuxième phrase. J'entends \zh{法律} (fǜlǜ) vers la fin. Cela confirme que le texte parle de l'action de l'État et du cadre légal. »
\end{exemplebox}

\courssub{Carte mentale lexicale : le champ de la corruption}

\begin{popfigure}
\begin{tikzpicture}[
    mindmap,
    grow cyclic,
    every node/.style={concept, execute at begin node=\hskip0pt},
    concept color=popblue!80,
    level 1/.append style={level distance=4.5cm, sibling angle=72, concept color=popblue!60},
    level 2/.append style={level distance=3cm, sibling angle=45, concept color=poporange!70},
    text=white,
    font=\small\bfseries
]
\node[align=center]{腐败\\fǔbài}
    child[concept color=popred!70] {node[align=center]{贪污\\tānwū\\détournement}
        child {node[align=center]{贪污犯\\tānwūfàn}}
        child {node[align=center]{巨款\\jùkuǎn}}
    }
    child[concept color=poporange!80] {node[align=center]{贿赂\\huìlù\\pot-de-vin}
        child {node[align=center]{行贿\\xínghuì}}
        child {node[align=center]{受贿\\shòuhuì}}
    }
    child[concept color=popgreen!70] {node[align=center]{法律\\fǜlǜ\\loi}
        child {node[align=center]{刑法\\xíngfǎ}}
        child {node[align=center]{监察\\jiānchá}}
    }
    child[concept color=poppurple!70] {node[align=center]{廉洁\\liánjié\\intégrité}
        child {node[align=center]{廉政\\liánzhèng}}
        child {node[align=center]{自律\\zìlǜ}}
    }
    child[concept color=poppink!70] {node[align=center]{惩罚\\chéngfá\\punition}
        child {node[align=center]{判刑\\pànxíng}}
        child {node[align=center]{双规\\shuāngguī}} % Note: Shuanggui is specific CCP discipline, maybe keep as "discipline interne" but okay for vocab context
    };
\end{tikzpicture}
\legende{Carte mentale du champ lexical « corruption » (vocabulaire HSK 3-4 visé).}
\end{popfigure}

\courssub{Le saviez-vous ? : La campagne anti-corruption en Chine}

\begin{savaistu}[打虎拍蝇 : « Frapper les tigres, écraser les mouches »]
Depuis 2012, la Chine mène une campagne anti-corruption sans précédent sous l'impulsion de Xi Jinping. Elle vise deux cibles symboliques :
\begin{itemize}
    \item \textbf{打虎 (dǎ hǔ) :} « Frapper les tigres » = sanctionner les \textbf{hauts cadres} (ministres, généraux, patrons d'entreprises d'État, secrétaires de province). Ex : Bo Xilai, Zhou Yongkang.
    \item \textbf{拍蝇 (pāi ying) :} « Écraser les mouches » = sanctionner les \textbf{petits fonctionnaires} de base (guichets d'administration, police locale, directeurs d'école, médecins) dont la corruption pourrit le quotidien des citoyens.
\end{itemize}
L'expression officielle est \zh{打虎、拍蝇、猎狐} (dǎ hǔ, pāi ying, liè hú) : « Frapper les tigres, écraser les mouches, chasser les renards » (ceux qui ont fui à l'étranger avec l'argent). Des milliers de fonctionnaires ont été sanctionnés. La Chine a aussi créé la \zh{国家监察委员会} (Commission nationale de supervision) en 2018, organe indépendant du Parti et de l'État, pour renforcer la \zh{监察} (jiānchá, supervision/contrôle).
\end{savaistu}

\courssub{Méthode : Aborder un texte chinois inconnu}

\begin{methode}[4 étapes pour une première lecture efficace]
\begin{enumerate}
    \item \textbf{Lire le titre} (et regarder l'image s'il y en a une) : activer ses connaissances du thème, prédire le vocabulaire.
    \item \textbf{Repérer les mots connus} (mots « amis », noms propres, chiffres, connecteurs logiques comme \zh{因为}、\zh{所以}、\zh{但是}) : se rassurer, poser des jalons.
    \item \textbf{Lire une première fois sans dictionnaire} : viser la compréhension globale (qui ? quoi ? où ? quand ? idée principale). Accepter de ne pas tout comprendre.
    \item \textbf{Vérifier avec le glossaire / le professeur} : lever les obstacles lexicaux précis, relire pour la compréhension détaillée.
\end{enumerate}
\end{methode}

\courssub{Exercice résolu : Identifier le vocabulaire du thème}

\begin{exoresolu}[Repérage lexical]
\textbf{Énoncé :} Parmi les phrases suivantes, surlignez (ou citez) les mots qui appartiennent au champ lexical de la corruption et de sa lutte. Donnez le pinyin et le sens.

1. \zh{他因为贪污被判刑了。} \\
2. \zh{我们要保持廉洁，拒绝贿赂。} \\
3. \zh{政府加强了监察力度。} \\

\textbf{Solution détaillée :}
\begin{enumerate}
    \item \zh{\textbf{贪污} (tānwū) : détournement de fonds ; \textbf{判刑} (pànxíng) : condamner à une peine de prison.} Phrase : « Il a été condamné à la prison pour détournement de fonds. »
    \item \zh{\textbf{廉洁} (liánjié) : intégrité, honnête ; \textbf{贿赂} (huìlù) : pot-de-vin.} Phrase : « Nous devons garder notre intégrité et refuser les pots-de-vin. »
    \item \zh{\textbf{监察} (jiānchá) : supervision, inspection, contrôle (administratif).} Phrase : « Le gouvernement a renforcé l'intensité du contrôle/supervision. »
\end{enumerate}
\end{exoresolu}

\courssub{Synthèse de la mise en route}

\begin{aretenir}[Bilan de la découverte]
\begin{itemize}
    \item Le mot-clé est \cle{腐败 fǔbài} (corruption). Il s'écrit avec les radicaux \zh{月} (chair) et \zh{攵} (action).
    \item Le titre \zh{《反腐败，人人有责》} annonce un texte sur la lutte collective.
    \item La stratégie de lecture : Titre $\rightarrow$ Mots connus $\rightarrow$ Lecture globale $\rightarrow$ Glossaire.
    \item Contexte culturel : En Chine, la campagne \zh{打虎拍蝇} cible à la fois les hauts cadres (« tigres ») et la petite corruption (« mouches »).
\end{itemize}
\end{aretenir}

\begin{exercice}[Entraînement personnel - Devoirs / Travail en autonomie]
\begin{enumerate}
    \item Écrire 3 fois le mot \zh{腐败} en respectant l'ordre des traits (voir section 4). Lire à voix haute à chaque fois.
    \item Compléter le tableau ci-dessous avec les mots de la carte mentale (pinyin + traduction).
    \item Préparer une phrase personnelle avec \zh{廉洁} ou \zh{贿赂} pour la prochaine séance.
\end{enumerate}
\end{exercice}

\coursec{Mise en route et découverte du thème}

Le thème de cette leçon est la \cle{corruption} (\zh{腐败} \emph{fǔbài}), un fléau qui touche toutes les sociétés, y compris le Cameroun et la Chine. Avant de lire le texte, activons nos connaissances.

\begin{experience}[Brainstorming : mots et idées associés]
En classe entière, au tableau, construisez une carte mentale autour du mot \zh{腐败}. L'enseignant note en chinois (caractères + pinyin) les propositions des élèves : mots du vocabulaire connu (\zh{坏} \emph{huài}, \zh{偷} \emph{tōu}, \zh{钱} \emph{qián}...), noms d'institutions (\zh{政府} \emph{zhèngfǔ}, \zh{警察} \emph{jǐngchá}), verbes d'action (\zh{抓} \emph{zhuā}, \zh{判} \emph{pàn}). Cela permet de vérifier les prérequis lexicaux (niveau HSK 1-2) et d'introduire naturellement le vocabulaire cible de la leçon (HSK 3-4).
\end{experience}

\begin{savaistu}[Contexte Cameroun -- Chine]
Au Cameroun, la lutte contre la corruption est pilotée par la \textbf{CONAC} (Commission Nationale Anti-Corruption, créée en 2006) et la \textbf{CONSUPE} (Commission Spéciale de la Surveillance des Marchés Publics). L'opération \emph{Épervier} (depuis 2006) a permis de sanctionner de nombreux détournements de fonds publics. En Chine, la \textbf{CCDI} (Commission Centrale pour la Discipline et l'Inspection du PCC) et la \textbf{Commission Nationale de Supervision} (depuis 2018) mènent la campagne \zh{反腐败} (\emph{fǎnfǔbài}) avec les opérations \zh{猎狐} (\emph{Lièhú}, « Chasse au renard », contre la fuite de corrompus à l'étranger) et \zh{天网} (\emph{Tiānwǎng}, « Filet céleste »). Les deux pays partagent le même objectif : \zh{零容忍} (\emph{líng róngrěn}, « tolérance zéro »).
\end{savaistu}

\coursec{Texte support : lecture et compréhension globale}

Nous allons maintenant étudier le texte support, un court article d'opinion intitulé \zh{反腐败，人人有责} (\emph{Fǎn fǔbài, rénrén yǒu zé} — « Lutter contre la corruption, c'est l'affaire de tous »). Objectifs : comprendre le sens global, repérer la structure argumentative, identifier le vocabulaire clé et les connecteurs logiques.

\courssub{Lecture du texte support}

\textbf{Étape 1 — Écoute active.} L'enseignant lit le texte à voix haute, lentement, en respectant scrupuleusement les tons et les pauses (virgules \zh{，}, points \zh{。}). Les élèves écoutent \emph{sans lire le texte}, puis suivent du doigt pendant une deuxième lecture.

\textbf{Étape 2 — Lecture en chaîne et correction phonétique.} Chaque élève lit une phrase. Le professeur corrige immédiatement les tons, points d'attention particuliers :
\begin{itemize}
\item \zh{腐败} \emph{fǔbài} (3ᵉ + 4ᵉ ton) : attention au 3ᵉ ton bas sur \zh{腐} (pas montant).
\item \zh{受贿} \emph{shòuhuì} (4ᵉ + 4ᵉ ton) : deux 4ᵉ tons, bien distincts.
\item \zh{廉洁} \emph{liánjié} (2ᵉ + 2ᵉ ton) : deux tons montants, ne pas « chanter » le second.
\item \zh{贿赂} \emph{huìlù} (4ᵉ + 4ᵉ ton) : \zh{赂} se lit \emph{lù} (4ᵉ), pas \emph{lǔ}.
\item \zh{制定} \emph{zhìdìng} (4ᵉ + 4ᵉ ton) : ne pas confondre avec \zh{制订} (synonyme, même prononciation).
\end{itemize}

\begin{textebox}[反腐败，人人有责 — Lutter contre la corruption, l'affaire de tous]
腐败是社会的一个大问题。\\
\emph{Fǔbài shì shèhuì de yí ge dà wèntí.}\\
La corruption est un grand problème de la société.\\[4pt]
有些官员贪污受贿，损害了人民的利益。\\
\emph{Yǒuxiē guānyuán tānwū shòuhuì, sǔnhài le rénmín de lìyì.}\\
Certains fonctionnaires se livrent à la corruption et aux pots-de-vin, ce qui nuit aux intérêts du peuple.\\[4pt]
中国政府非常重视反腐败工作，制定了很多法律，惩罚腐败分子。\\
\emph{Zhōngguó zhèngfǔ fēicháng zhòngshì fǎnfǔbài gōngzuò, zhìdìng le hěn duō fǎlǜ, chéngfá fǔbài fènzǐ.}\\
Le gouvernement chinois attache une grande importance à la lutte contre la corruption : il a élaboré de nombreuses lois et punit les corrompus.\\[4pt]
近年来，反腐败取得了很大的成绩。\\
\emph{Jìnnián lái, fǎnfǔbài qǔdé le hěn dà de chéngjì.}\\
Ces dernières années, la lutte contre la corruption a obtenu de grands résultats.\\[4pt]
老百姓希望政府更加廉洁，社会更加公平。\\
\emph{Lǎobǎixìng xīwàng zhèngfǔ gèngjiā liánjié, shèhuì gèngjiā gōngpíng.}\\
Les gens du peuple espèrent un gouvernement plus intègre et une société plus juste.\\[4pt]
反腐败不只是政府的事情，也是我们每个人的责任。\\
\emph{Fǎnfǔbài bù zhǐ shì zhèngfǔ de shìqing, yě shì wǒmen měi ge rén de zérèn.}\\
Lutter contre la corruption n'est pas seulement l'affaire du gouvernement, c'est aussi la responsabilité de chacun de nous.\\[4pt]
我们应该从小事做起，诚实做人，拒绝贿赂。\\
\emph{Wǒmen yīnggāi cóng xiǎoshì zuòqǐ, chéngshí zuòrén, jùjué huìlù.}\\
Nous devons commencer par les petites choses, être honnêtes et refuser les pots-de-vin.\\[4pt]
只有大家共同努力，社会才会更加美好。\\
\emph{Zhǐyǒu dàjiā gòngtóng nǔlì, shèhuì cái huì gèngjiā měihǎo.}\\
Seulement si tout le monde fait des efforts ensemble, la société deviendra meilleure.
\end{textebox}

\courssub{Compréhension globale du texte}

Avant l'analyse détaillée, répondons à trois questions pour saisir la macro-structure :

\begin{enumerate}
\item \textbf{De quoi parle le texte ?} Du phénomène de la \cle{corruption} (\zh{腐败}) comme problème social majeur, et de la \cle{lutte contre la corruption} (\zh{反腐败}) comme responsabilité partagée.
\item \textbf{Quels sont les acteurs ?} Les \cle{fonctionnaires} (\zh{官员}) fautifs, le \cle{gouvernement} (\zh{政府}) qui légifère et sanctionne, le \cle{peuple} (\zh{人民} / \zh{老百姓}) qui attend des résultats, et \cle{chacun de nous} (\zh{我们每个人}).
\item \textbf{Quelle est la thèse finale (dernière phrase) ?} \zh{只有大家共同努力，社会才会更加美好} : la condition nécessaire (\zh{只有\ldots 才\ldots}) d'une société meilleure est l'effort collectif.
\end{enumerate}

\begin{popfigure}
\begin{tikzpicture}[node distance=0.55cm and 0.9cm,
  bloc/.style={rectangle, rounded corners=4pt, draw=popline, thick, align=center, minimum width=3.1cm, minimum height=1.15cm, font=\small},
  fleche/.style={-{Stealth[length=3mm]}, thick, popdark}]
\node[bloc, fill=poppinkL, align=center] (pb){\textbf{Constat : Problème}\\腐败严重\\\emph{fǔbài yánzhòng}};
\node[bloc, fill=popblueL, right=of pb, align=center] (ac){\textbf{Action : Gouvernement}\\法律、惩罚\\\emph{fǎlǜ, chéngfá}};
\node[bloc, fill=popgreenL, right=of ac, align=center] (re){\textbf{Bilan : Résultats}\\取得成绩\\\emph{qǔdé chéngjì}};
\node[bloc, fill=popgoldL, right=of re, align=center] (ap){\textbf{Appel : Responsabilité partagée}\\人人有责\\\emph{rénrén yǒu zé}};
\draw[fleche] (pb) -- (ac);
\draw[fleche] (ac) -- (re);
\draw[fleche] (re) -- (ap);
\end{tikzpicture}
\legende{Structure argumentative du texte : Problème $\rightarrow$ Mesures $\rightarrow$ Bilan $\rightarrow$ Appel à l'action collective.}
\end{popfigure}

Cette progression \textbf{Constat $\rightarrow$ Mesures $\rightarrow$ Bilan $\rightarrow$ Appel} est classique dans l'écriture argumentative chinoise (éditoriaux, discours officiels, dissertations). La repérer facilite grandement la compréhension de textes plus longs.

\courssub{Glossaire du thème et prononciation}

Voici le vocabulaire essentiel du thème « Corruption et lutte anticorruption » (niveau HSK 3-4). Apprenez les caractères, le pinyin \textbf{avec les tons}, la nature grammaticale et le sens.

\begin{center}
\begin{tabularx}{\linewidth}{|X|X|X|X|}
\hline
\rowcolor{popblueL} Caractères & Pinyin & Nature & Traduction \\
\hline
腐败 & fǔbài & n. / adj. & corruption ; corrompu \\
\hline
反腐败 & fǎnfǔbài & n. & lutte contre la corruption \\
\hline
贪污 & tānwū & v. & détourner (fonds publics), se corrompre \\
\hline
受贿 & shòuhuì & v. & recevoir un pot-de-vin (corruption passive) \\
\hline
行贿 & xínghuì & v. & donner un pot-de-vin (corruption active) \\
\hline
贿赂 & huìlù & n. / v. & pot-de-vin ; soudoyer \\
\hline
廉洁 & liánjié & adj. & intègre, honnête (qualifie un responsable) \\
\hline
公平 & gōngpíng & adj. & juste, équitable \\
\hline
惩罚 & chéngfá & v. / n. & punir ; punition, sanction \\
\hline
责任 & zérèn & n. & responsabilité \\
\hline
诚实 & chéngshí & adj. & honnête, sincère (qualifie une personne) \\
\hline
拒绝 & jùjué & v. & refuser, rejeter \\
\hline
制定 & zhìdìng & v. & élaborer, formuler (loi, plan, règle) \\
\hline
分子 & fènzǐ & n. & élément, individu (souvent péjoratif : \zh{腐败分子}) \\
\hline
近年来 & jìnnián lái & t. & ces dernières années, depuis quelques années \\
\hline
更加 & gèngjiā & adv. & encore plus, d'avantage \\
\hline
共同 & gòngtóng & adj. / adv. & commun, ensemble, conjointement \\
\hline
努力 & nǔlì & v. / adv. & faire des efforts ; avec effort \\
\hline
美好 & měihǎo & adj. & beau, meilleur, heureux (avenir, vie) \\
\hline
从小事做起 & cóng xiǎoshì zuòqǐ & loc. & commencer par les petites choses \\
\hline
\end{tabularx}
\end{center}

\textbf{Point de prononciation :} Les mots composées \zh{贪污} (\emph{tānwū}), \zh{受贿} (\emph{shòuhuì}), \zh{行贿} (\emph{xínghuì}), \zh{贿赂} (\emph{huìlù}), \zh{廉洁} (\emph{liánjié}), \zh{惩罚} (\emph{chéngfá}), \zh{制定} (\emph{zhìdìng}) sont tous des mots bisyllabiques. Respectez la neutralisation tonale possible sur le second syllabe en parole rapide (ex: \emph{fǔbai} léger), mais en diction soignée (lecture, examen oral), prononcez les deux tons pleins.

\coursec{Clés (radicaux) et ordre des traits}

L'analyse des radicaux (\zh{部首} \emph{bùshǒu}) aide à mémoriser le sens et à classer les caractères dans le dictionnaire. Voici 7 caractères clés du thème.

\begin{definition}[Analyse de caractères clés : radicaux et ordre des traits]
\begin{itemize}
\item \zh{\textbf{腐} (fǔ)} : Radical \zh{月} (viande, \emph{yuè/zì}, 4 traits) -- partie gauche. Phonétique \zh{府} (\emph{fǔ}). \textbf{Ordre} : haut \rightarrow bas (月), puis gauche \rightarrow droite (府). Total : 13 traits.
\item \zh{\textbf{败} (bài)} : Radical \zh{攵} (frapper, \emph{bù}, 4 traits) -- partie droite. \textbf{Ordre} : écrire la partie gauche \zh{啇} (haut \rightarrow bas), puis le radical \zh{攵} (point, horizontal, crochet, point). Total : 11 traits.
\item \zh{\textbf{贿} (huì)} : Radical \zh{贝} (coquillage/argent, \emph{bèi}, 4 traits) -- partie basse. Phonétique \zh{穴} (\emph{xué}) en haut. \textbf{Ordre} : haut \rightarrow bas (穴 puis 贝). Total : 13 traits. \textit{Note : 贝 en bas s'écrit différemment de 贝 seul (dernier trait point au lieu de horizontal).}
\item \zh{\textbf{廉} (lián)} : Radical \zh{广} (toit, \emph{guǎng}, 3 traits) -- haut gauche. \textbf{Ordre} : radical \zh{广} (point, horizontal, crochet), puis l'intérieur \zh{兼} (haut \rightarrow bas). Total : 13 traits.
\item \zh{\textbf{洁} (jié)} : Radical \zh{氵} (eau, \emph{shuǐ}, 3 traits) -- gauche. Phonétique \zh{齐} (\emph{qí}). \textbf{Ordre} : radical \zh{氵} (deux points, trait montant), puis \zh{齐} (haut \rightarrow bas). Total : 9 traits.
\item \zh{\textbf{惩} (chéng)} : Radical \zh{忄} (cœur, \emph{xīn}, 3 traits) -- gauche. Phonétique \zh{徵} (\emph{zhēng}). \textbf{Ordre} : radical \zh{忄} (deux points, trait vertical), puis \zh{徵} (haut \rightarrow bas). Total : 11 traits.
\item \zh{\textbf{责} (zé)} : Radical \zh{贝} (coquillage/argent, \emph{bèi}, 4 traits) -- bas. Haut : \zh{朿} (\emph{cī}). \textbf{Ordre} : haut \rightarrow bas. Total : 8 traits.
\end{itemize}
\end{definition}

\begin{methode}[Écrire correctement les caractères composés]
\begin{enumerate}
\item Identifier le radical (clé de dictionnaire) et le composant phonétique/sémantique.
\item Respecter l'ordre universel : \textbf{haut $\rightarrow$ bas}, \textbf{gauche $\rightarrow$ droite}, \textbf{horizontal $\rightarrow$ vertical}, \textbf{extérieur $\rightarrow$ intérieur $\rightarrow$ fermeture}.
\item Pour les radicaux \zh{忄}, \zh{氵}, \zh{月}, \zh{广} : écrire le radical \textbf{en premier} (sauf \zh{贝} en bas qui s'écrit \textbf{en dernier}).
\item S'entraîner sur papier quadrillé (carreaux de 1.5~cm) : un caractère par case, centré.
\end{enumerate}
\end{methode}

\coursec{Questions de compréhension et approfondissement}

\courssub{Repérage des connecteurs logiques}

Le texte utilise des connecteurs pour articuler l'argumentation. Repérez-les dans le texte (encadrez-les) et observez leur fonction.

\begin{center}
\begin{tabularx}{\linewidth}{|X|X|X|}
\hline
\rowcolor{popblueL} Connecteur & Exemple / Position & Fonction logique \\
\hline
\zh{因为\ldots 所以\ldots} (\emph{yīnwèi\ldots suǒyǐ\ldots}) & Implicite : \zh{损害了利益，所以制定法律} & Cause $\rightarrow$ Conséquence \\
\hline
\zh{不但\ldots 而且\ldots} (\emph{búdàn\ldots érqiě\ldots}) & \zh{不只是\ldots 也是\ldots} (phrase 6) & Accumulation, progression \\
\hline
\zh{只有\ldots 才\ldots} (\emph{zhǐyǒu\ldots cái\ldots}) & \zh{只有大家共同努力，社会才会更加美好} (phrase 8) & Condition \textbf{nécessaire} \\
\hline
\zh{应该} (\emph{yīnggāi}) & \zh{我们应该从小事做起} (phrase 7) & Obligation morale, conseil \\
\hline
\end{tabularx}
\end{center}

\begin{propriete}[La structure 只有\ldots 才\ldots (\emph{zhǐyǒu\ldots cái\ldots})]
Exprime une \cle{condition nécessaire} : « Seulement si [condition] \ldots , [alors seulement] [résultat] \ldots ».

\textbf{Schéma syntaxique :} \zh{只有} + Condition (sujet + verbe) \zh{，} + Sujet + \zh{才} + Verbe / Adjectif.

\begin{center}
\begin{tabularx}{\linewidth}{|X|X|X|}
\hline
\rowcolor{popblueL} Structure & Exemple (Caractères + Pinyin) & Traduction \\
\hline
\zh{只有\ldots 才\ldots} & \zh{只有大家共同努力，社会才会更加美好。} \emph{Zhǐyǒu dàjiā gòngtóng nǔlì, shèhuì cái huì gèngjiā měihǎo.} & Seulement si tout le monde s'unit, la société s'améliorera. \\
\hline
\zh{只有\ldots 才\ldots} & \zh{只有拒绝贿赂，官员才会廉洁。} \emph{Zhǐyǒu jùjué huìlù, guānyuán cái huì liánjié.} & Seulement en refusant les pots-de-vin les fonctionnaires seront intègres. \\
\hline
\zh{只有\ldots 才\ldots} & \zh{只有努力学习中文，你才能通过HSK4。} \emph{Zhǐyǒu nǔlì xuéxí Zhōngwén, nǐ cái néng tōngguò HSK4.} & Ce n'est qu'en travaillant le chinois que tu passeras le HSK4. \\
\hline
\end{tabularx}
\end{center}

\textbf{Points de vigilance (erreurs fréquentes) :}
\begin{enumerate}
\item \textbf{Ne pas omettre \zh{才}.} *\zh{只有努力，你能成功} $\rightarrow$ \textbf{Faux}. Correct : \zh{只有努力，你\underline{才}能成功}.
\item \textbf{Ordre des mots :} La condition suit \zh{只有}, le résultat suit \zh{才}. Ne pas inverser.
\item \textbf{Distinction avec \zh{只要\ldots 就\ldots} (\emph{zhǐyào\ldots jiù\ldots}) :} \zh{只要} = « du moment que », condition \emph{suffisante} (facile à réaliser). \zh{只有} = « seulement si », condition \emph{nécessaire} (difficile, exclusive). Ex: \zh{只要你来，我就去} (Si tu viens, j'irai) vs \zh{只有你来，我才去} (Seulement si TU viens [et personne d'autre], j'irai).
\end{enumerate}
\end{propriete}

\begin{attention}[受贿 (\emph{shòuhuì}) vs 行贿 (\emph{xínghuì}) : le piège classique !]
C'est la confusion n°1 en examen de vocabulaire et de compréhension.
\begin{itemize}
\item \zh{\textbf{受贿} (\emph{shòuhuì})} : \zh{受} (\emph{shòu}) = recevoir, subir. \textbf{Sujet = corrompu (celui qui prend l'argent).} \textit{Ex: 官员受贿。}
\item \zh{\textbf{行贿} (\emph{xínghuì})} : \zh{行} (\emph{xíng}) = faire, pratiquer, commettre. \textbf{Sujet = corrupteur (celui qui donne l'argent).} \textit{Ex: 商人行贿。}
\end{itemize}
\textbf{Astuce mnémotechnique :} \zh{受} (toit + main qui reçoit) $\rightarrow$ \textbf{Recevoir}. \zh{行} (croix + pas) $\rightarrow$ \textbf{Agir/Donner}.
Dans le texte : \zh{官员\ldots 受贿} (Les fonctionnaires \textbf{reçoivent}). Celui qui donne l'argent au fonctionnaire commet \zh{行贿}. Les deux sont punis par la loi (\zh{惩罚}).
\end{attention}

\courssub{Méthode : répondre aux questions de compréhension écrite}

\begin{methode}[Trois étapes pour une réponse réussie]
\begin{enumerate}
\item \textbf{Repérer le mot-clé} de la question dans le texte (qui, quoi, comment, pourquoi $\rightarrow$ \zh{谁}, \zh{什么}, \zh{怎么}, \zh{为什么}).
\item \textbf{Localiser et recopier} la phrase exacte du texte qui contient l'information.
\item \textbf{Reformuler} en phrase complète, en reprenant le sujet de la question et en adaptant le pronom/personne si nécessaire. Ne pas copier tout le paragraphe.
\end{enumerate}
\end{methode}

\begin{exoresolu}[Application guidée : Question \zh{怎么} (Comment)]
\textbf{Énoncé.} 中国政府怎么反腐败？\emph{Zhōngguó zhèngfǔ zěnme fǎnfǔbài ?} — Comment le gouvernement chinois lutte-t-il contre la corruption ?
\tcblower
\textbf{Solution détaillée.}\\
\textbf{Étape 1 :} Mots-clés : \zh{中国政府} (sujet), \zh{怎么} (manière), \zh{反腐败} (action).\\
\textbf{Étape 2 :} Localisation : Phrase 3 : \zh{中国政府非常重视反腐败工作，制定了很多法律，惩罚腐败分子。}\\
\textbf{Étape 3 :} Reformulation (réponse attendue en examen) :\\
\zh{中国政府非常重视反腐败工作，制定了很多法律，惩罚腐败分子。}\\
\emph{Zhōngguó zhèngfǔ fēicháng zhòngshì fǎnfǔbài gōngzuò, zhìdìng le hěn duō fǎlǜ, chéngfá fǔbài fènzǐ.}\\
\textit{Traduction :} Le gouvernement chinois attache une grande importance au travail de lutte contre la corruption, a formulé beaucoup de lois, et punit les éléments corrompus.
\end{exoresolu}

\begin{exercice}[Entraînement : Questions de compréhension]
Répondez en chinois, phrases complètes, caractères + pinyin.
\begin{enumerate}
\item 腐败损害了谁的利益？\emph{Fǔbài sǔnhài le shéi de lìyì ?}
\item 老百姓希望政府和社会怎么样？\emph{Lǎobǎixìng xīwàng zhèngfǔ hé shèhuì zěnmeyàng ?}
\item 根据课文，我们每个人应该做什么？\emph{Gēnjù kèwén, wǒmen měi ge rén yīnggāi zuò shénme ?}
\item 为什么说「反腐败，人人有责」？\emph{Wèishénme shuō « Fǎnfǔbài, rénrén yǒu zé » ?} (Réponse inférée à partir de la dernière phrase).
\end{enumerate}
\end{exercice}

\begin{corrige}[Exercice « Entraînement : Questions de compréhension »]
\begin{enumerate}
\item \zh{腐败损害了人民的利益。} \emph{Fǔbài sǔnhài le rénmín de lìyì.} (La corruption nuit aux intérêts du peuple.)
\item \zh{老百姓希望政府更加廉洁，社会更加公平。} \emph{Lǎobǎixìng xīwàng zhèngfǔ gèngjiā liánjié, shèhuì gèngjiā gōngpíng.} (Le peuple espère un gouvernement plus intègre et une société plus juste.)
\item \zh{我们应该从小事做起，诚实做人，拒绝贿赂。} \emph{Wǒmen yīnggāi cóng xiǎoshì zuòqǐ, chéngshí zuòrén, jùjué huìlù.} (Nous devons commencer par les petites choses, être honnêtes, refuser les pots-de-vin.)
\item \zh{因为只有大家共同努力，社会才会更加美好。} \emph{Yīnwèi zhǐyǒu dàjiā gòngtóng nǔlì, shèhuì cái huì gèngjiā měihǎo.} (Parce que seulement si tout le monde fait des efforts ensemble, la société deviendra meilleure.)
\end{enumerate}
\end{corrige}

\coursec{Production guidée et réinvestissement}

\courssub{Expression orale : Exposé court (1-2 min)}

\begin{experience}[Présentation : « Mon engagement contre la corruption »]
\textbf{Consigne :} Préparez un exposé de 6-8 phrases en vous appuyant sur le vocabulaire et la structure \zh{只有\ldots 才\ldots}. Utilisez le canevas ci-dessous.
\begin{enumerate}
\item \textbf{Constat :} \zh{腐败是\ldots 问题。} (La corruption est un problème \ldots)
\item \textbf{Exemple local :} \zh{在喀麦隆，\ldots} (Au Cameroun, \ldots -- ex: \zh{贪污}, \zh{受贿} dans les marchés, administrations).
\item \textbf{Action gouvernementale :} \zh{政府\ldots{}，制定了\ldots{}，惩罚\ldots} (Gouvernement \ldots, lois \ldots, punit \ldots).
\item \textbf{Appel personnel :} \zh{只有\ldots 才\ldots} (Seulement si \ldots alors \ldots).
\item \textbf{Engagement concret :} \zh{我应该\ldots{}，拒绝\ldots} (Je dois \ldots, refuser \ldots).
\end{enumerate}
\textbf{Évaluation :} Prononciation (tons), fluidité, utilisation correcte de \zh{只有\ldots 才\ldots}, richesse lexicale (min. 5 mots du glossaire).
\end{experience}

\courssub{Expression écrite : Message citoyen (SMS / Réseau social)}

\begin{methode}[Rédiger un message court d'intérêt général]
\begin{enumerate}
\item \textbf{Accroche :} Titre clair (\zh{拒绝腐败，从我做起}).
\item \textbf{Argument :} Une phrase avec \zh{因为\ldots 所以\ldots} ou \zh{只有\ldots 才\ldots}.
\item \textbf{Appel à l'action :} Verbes d'action (\zh{举报} \emph{jǔbào} signaler, \zh{监督} \emph{jiāndū} surveiller, \zh{拒绝} \emph{jùjué} refuser).
\item \textbf{Signature/Hashtag :} \zh{—— 一名中学生} / \zh{\#廉洁喀麦隆\#}.
\end{enumerate}
\end{methode}

\begin{exemplebox}[Modèle de message citoyen (WeChat / WhatsApp / SMS)]
\zh{【拒绝腐败，从我做起】}\\
\zh{腐败损害大家的利益。只有人人拒绝贿赂，社会才会公平。}\\
\zh{发现腐败，请举报：12388 (中国) / CONAC 1517 (喀麦隆)。}\\
\zh{—— 一名喀麦隆中学生 \#廉洁社会\# \#反腐倡廉\#}\\
\emph{【Jùjué fǔbài, cóng wǒ zuòqǐ】 Fǔbài sǔnhài dàjiā de lìyì. Zhǐyǒu rénrén jùjué huìlù, shèhuì cái huì gōngpíng. Fāxiàn fǔbài, qǐng jǔbào: 12388 (Zhōngguó) / CONAC 1517 (Kāmàilóng). — Yī míng Kāmàilóng zhōngxuéshēng \#Liánjié shèhuì\# \#Fǎnfūlián\#}\\
\textit{« [Refuser la corruption, ça commence par moi] La corruption nuit aux intérêts de tous. Seulement si chacun refuse les pots-de-vin, la société sera juste. Si vous constatez de la corruption, signalez-le : 12388 (Chine) / CONAC 1517 (Cameroun). — Un lycéen camerounais \#SociétéIntègre\# \#LutteAntiCorruption\# »}
\end{exemplebox}

\begin{exercice}[Production écrite : À vous de jouer]
Rédigez un message similaire (60-80 caractères chinois) adapté au contexte de votre lycée (ex: corruption aux examens, pots-de-vin pour notes, détournement fonds association). Utilisez \zh{只有\ldots 才\ldots} et au moins 4 mots du glossaire.
\end{exercice}

\coursec{Élément socioculturel approfondi : Cameroun -- Chine}

\begin{savaistu}[Les institutions de la lutte anticorruption : parallèles institutionnels]
\begin{center}
\begin{tabularx}{\linewidth}{|X|X|X|}
\hline
\rowcolor{popblueL} Aspect & Cameroun & Chine \\
\hline
Institution suprême & \textbf{CONAC} (Commission Nationale Anti-Corruption, 2006) & \textbf{CCDI} (Commission Centrale Discipline/Inspection) + \textbf{Commission Nationale de Supervision} (2018) \\
\hline
Opérations emblématiques & \emph{Opération Épervier} (depuis 2006, récupération fonds publics) & \zh{猎狐行动} (\emph{Lièhú Xíngdòng}, « Chasse au renard », fuite à l'étranger) ; \zh{天网行动} (\emph{Tiānwǎng Xíngdòng}, « Filet céleste ») \\
\hline
Numéro vert / Signalement & \textbf{1517} (Appel gratuit CONAC) & \textbf{12388} (Site web \& hotline CCDI) \\
\hline
Vocabulaire institutionnel & \zh{贪污} \emph{tānwū}, \zh{受贿} \emph{shòuhuì}, \zh{挪用公款} \emph{nuóyòng gōngkuǎn} & \zh{违纪} \emph{wéijì} (violation discipline), \zh{违法} \emph{wéifǎ} (illégal), \zh{双规} \emph{shuāngguī} (ancien: double règlement), \zh{留置} \emph{liúzhì} (rétention/détention supervisée) \\
\hline
\end{tabularx}
\end{center}
\textbf{Point de culture générale :} La Chine a inscrit la lutte anticorruption dans sa Constitution (2018, création Commission de Supervision). Le Cameroun a ratifié la Convention de l'Union Africaine sur la prévention et la lutte contre la corruption (2003). Les deux pays coopèrent via le Forum sur la Coopération Sino-Africaine (FOCAC) sur la gouvernance.
\end{savaistu}

\begin{aretenir}[Synthèse de la leçon 27 — Vocabulaire \& Compréhension : Corruption]
\begin{itemize}
\item \textbf{Texte :} Structure argumentative \textbf{Problème $\rightarrow$ Action Gov $\rightarrow$ Résultat $\rightarrow$ Responsabilité partagée}. Titre : \zh{反腐败，人人有责}.
\item \textbf{Vocabulaire clé (HSK 3-4) :} \zh{腐败、贪污、受贿、行贿、贿赂、廉洁、惩罚、制定、分子、近年来、更加、共同、努力、美好、从小事做起}. Distinction vitale : \zh{受贿} (recevoir) $\neq$ \zh{行贿} (donner).
\item \textbf{Grammaire :} \zh{只有\ldots 才\ldots} = condition nécessaire (obligation de \zh{才}). Différence avec \zh{只要\ldots 就\ldots} (condition suffisante). Connecteurs : \zh{不但\ldots 而且\ldots} / \zh{不只是\ldots 也是\ldots}.
\item \textbf{Écriture :} 7 caractères analysés (radicaux \zh{月、攵、贝、广、氵、忄} + ordre traits).
\item \textbf{Compétences :} Lire/Comprendre texte d'opinion $\rightarrow$ Répondre questions (Méthode 3 étapes) $\rightarrow$ S'exprimer oralement (exposé) $\rightarrow$ Écrire message citoyen.
\item \textbf{Culture :} CONAC / 1517 (Cameroun) vs CCDI-Supervision / 12388 (Chine). Objectif commun : \zh{零容忍} (tolérance zéro).
\end{itemize}
\end{aretenir}


\coursec{Glossaire du thème et prononciation}

Après avoir découvert le texte support dans la section précédente et en avoir saisi le sens global, nous allons maintenant nous concentrer sur le \textbf{vocabulaire spécialisé} du thème de la corruption. Maîtriser ces termes, leurs tons précis, leurs collocations et leur structure interne (radicaux) est indispensable pour comprendre des articles de presse, des rapports officiels ou pour participer à un débat citoyen en chinois.

\courssub{Vocabulaire thématique : tableau récapitulatif}

Le tableau ci-dessous regroupe les \textbf{16 mots-clés} du thème. Observez la nature grammaticale (certains mots fonctionnent comme nom \emph{et} adjectif, d'autres comme verbe \emph{et} nom) et mémorisez les tons. La prononciation exacte est votre passeport pour l'oral.

\begin{popfigure}
\legende{Glossaire du thème « Corruption » — 1ère A}
\begin{center}
\begin{tabularx}{\linewidth}{|>{\centering\arraybackslash}X|>{\centering\arraybackslash}X|>{\centering\arraybackslash}X|X|}
\hline
\rowcolor{popblueL} \textbf{Caractères} & \textbf{Pinyin} & \textbf{Nature} & \textbf{Traduction} \\ \hline
腐败 & fǔbài & n. / adj. & corruption ; corrompu \\ \hline
贪污 & tānwū & v. & détourner (des fonds publics), piller \\ \hline
贿赂 & huìlù & v. / n. & corrompre / pot-de-vin, argent sale \\ \hline
受贿 & shòuhuì & v. & accepter un pot-de-vin (être corrompu) \\ \hline
行贿 & xínghuì & v. & offrir un pot-de-vin (corrompre qqn) \\ \hline
官员 & guānyuán & n. & fonctionnaire, officiel \\ \hline
廉洁 & liánjié & adj. & intègre, honnête, incorruptible \\ \hline
公平 & gōngpíng & adj. & équitable, juste \\ \hline
法律 & fǎlǜ & n. & loi, législation \\ \hline
惩罚 & chéngfá & v. / n. & punir / punition, sanction \\ \hline
打击 & dǎjī & v. & combattre, réprimer, frapper \\ \hline
责任 & zérèn & n. & responsabilité, devoir \\ \hline
诚实 & chéngshí & adj. & honnête, sincère \\ \hline
拒绝 & jùjué & v. & refuser, rejeter \\ \hline
利益 & lìyì & n. & intérêt, profit, avantage \\ \hline
举报 & jǔbào & v. & signaler, dénoncer \\ \hline
\end{tabularx}
\end{center}
\end{popfigure}

\begin{propriete}[Nature grammaticale variable]
Plusieurs mots de cette liste appartiennent à deux catégories grammaticales sans changer de forme. Le contexte (place dans la phrase, présence de \zh{的} \emph{de}, aspect verbal \zh{了/过/着}) détermine leur fonction.
\begin{itemize}
    \item \zh{腐败 fǔbài} : \textbf{Nom} \zh{反腐败} (lutte contre la corruption) / \textbf{Adjectif} \zh{腐败官员} (fonctionnaire corrompu).
    \item \zh{贿赂 huìlù} : \textbf{Verbe} \zh{贿赂官员} (corrompre un fonctionnaire) / \textbf{Nom} \zh{收受贿赂} (accepter des pots-de-vin).
    \item \zh{惩罚 chéngfá} : \textbf{Verbe} \zh{惩罚罪犯} (punir le criminel) / \textbf{Nom} \zh{受到惩罚} (subir une sanction).
\end{itemize}
\end{propriete}

\courssub{Travail sur les tons difficiles : discrimination auditive}

Le chinois est une langue tonale : un changement de ton change le sens du mot. Pour le vocabulaire de la corruption, quatre paires de tons posent souvent problème aux francophones. Écoutez attentivement la différence de \textbf{contour mélodique} (montant, descendant, haut, bas).

\begin{definition}[Les 4 paires critiques de la leçon]
\begin{enumerate}
    \item \zh{\textbf{腐败} fǔbài} \textbf{(3-4)} : Ton 3 (bas-descendant-remontant \emph{mais souvent demi-ton 3 bas en milieu de phrase}) + Ton 4 (descendant sec). \emph{Attention : pas « fúbài » (2-4) ni « fǔbǎi » (3-3).}
    \item \zh{\textbf{贪污} tānwū} \textbf{(1-1)} : Deux tons 1 (haut plat). Stabilité absolue. \emph{Piège : ne pas faire descendre le 2e ton (wū ≠ wǔ).}
    \item \zh{\textbf{贿赂} huìlù} \textbf{(4-4)} : Deux tons 4 (descendant sec, comme un « Non ! » ferme). Énergie identique sur les deux syllabes.
    \item \zh{\textbf{廉洁} liánjié} \textbf{(2-2)} : Deux tons 2 (montant, comme une question « Hein ? »). Montée continue.
\end{enumerate}
\end{definition}

\begin{methode}[Exercice de discrimination auditive — « Levez la carte »]
\textbf{Matériel :} 4 cartes par élève (ou 4 doigts levés) étiquetées A, B, C, D.
\textbf{Déroulement :}
\begin{enumerate}
    \item L'enseignant prononce \textbf{un seul mot} de la liste (ex: \zh{huìlù}).
    \item Les élèves identifient la \textbf{paire de tons} correspondante et lèvent la bonne carte :
        \begin{itemize}
            \item Carte A : \zh{fǔbài} (3-4)
            \item Carte B : \zh{tānwū} (1-1)
            \item Carte C : \zh{huìlù} (4-4)
            \item Carte D : \zh{liánjié} (2-2)
        \end{itemize}
    \item L'enseignant vérifie visuellement, corrige si besoin, fait répéter en chœur.
    \item \textbf{Variante :} L'enseignant donne la définition en français (« Pot-de-vin »), les élèves donnent le mot chinois \emph{avec les tons justes}.
\end{enumerate}
\textbf{Objectif :} Automatiser la perception et la production des contours tonaux pour ces mots à haute fréquence.
\end{methode}

\begin{exemplebox}[Exemples contextualisés pour l'entraînement tonal]
\begin{itemize}
    \item \zh{这个官员很腐败。} \emph{Zhège guānyuán hěn fǔbài.} (Cet officiel est très corrompu.) — \textbf{Adj.}
    \item \zh{我们要反腐败。} \emph{Wǒmen yào fǎn fǔbài.} (Nous devons lutter contre la corruption.) — \textbf{Nom.}
    \item \zh{他因贪污被捕。} \emph{Tā yīn tānwū bèi bǔ.} (Il a été arrêté pour détournement de fonds.) — \textbf{Verbe.}
    \item \zh{拒绝贿赂，人人有责。} \emph{Jùjué huìlù, rénrén yǒu zé.} (Refuser les pots-de-vin, c'est le devoir de tous.) — \textbf{Nom.}
    \item \zh{他贿赂了法官。} \emph{Tā huìlù le fǎguān.} (Il a corrompu le juge.) — \textbf{Verbe.}
    \item \zh{我们要选廉洁的领导。} \emph{Wǒmen yào xuǎn liánjié de lǐngdǎo.} (Nous devons élire des dirigeants intègres.) — \textbf{Adj.}
\end{itemize}
\end{exemplebox}

\courssub{Collocations indispensables (chengyu / suites lexicales)}

En chinois, les mots ne vivent pas isolés. Pour parler de corruption avec justesse et naturel, vous devez mémoriser ces \textbf{associations figées} (collocations). Elles fonctionnent comme des « blocs » prêts à l'emploi.

\begin{aretenir}[5 collocations à connaître par cœur]
\begin{center}
\begin{tabularx}{\linewidth}{|>{\centering\arraybackslash}X|>{\centering\arraybackslash}X|X|}
\hline
\rowcolor{poporangeL} \textbf{Chinois (Pinyin)} & \textbf{Structure} & \textbf{Traduction / Usage} \\ \hline
\zh{反腐败} (fǎn fǔbài) & Verbe \zh{反} + Nom \zh{腐败} & Lutter contre la corruption (terme officiel, slogan). \\ \hline
\zh{贪污受贿} (tānwū shòuhuì) & Verbe \zh{贪污} + Verbe \zh{受贿} & Détournements de fonds \textbf{et} acceptation de pots-de-vin (chef d'accusation standard). \\ \hline
\zh{打击腐败分子} (dǎjī fǔbài fènzǐ) & Verbe \zh{打击} + Nom composé & Réprimer / frapper les éléments corrompus (langage policier/judiciaire). \\ \hline
\zh{廉洁政府} (liánjié zhèngfǔ) & Adj \zh{廉洁} + Nom \zh{政府} & Gouvernement intègre / propre (objectif politique). \\ \hline
\zh{拒绝贿赂} (jùjué huìlù) & Verbe \zh{拒绝} + Nom \zh{贿赂} & Refuser les pots-de-vin (engagement moral, serment). \\ \hline
\end{tabularx}
\end{center}
\end{aretenir}

\begin{exoresolu}[Emploi des collocations dans une phrase complète]
\textbf{Énoncé :} Traduisez en chinois : « Le gouvernement camerounais lutte fermement contre la corruption et veut construire une administration intègre. »
\\
\textbf{Solution détaillée :}
\begin{enumerate}
    \item \textbf{Analyse :} « Lutte fermement contre la corruption » $\rightarrow$ \zh{坚决反腐败} (jiānjué fǎn fǔbài). « Construire » $\rightarrow$ \zh{建设} (jiànshè) ou \zh{建立} (jiànlì). « Administration intègre » $\rightarrow$ \zh{廉洁政府} ou \zh{廉洁的行政机关} (liánjié de xíngzhèng jīguān).
    \item \textbf{Assemblage :} \zh{喀麦隆政府坚决反腐败，致力于建设廉洁政府。}
    \item \textbf{Pinyin :} Kāmàilóng zhèngfǔ jiānjué fǎn fǔbài, zhìlì yú jiànshè liánjié zhèngfǔ.
    \item \textbf{Vérification tons :} jiān-jué (1-2), fǎn fǔ-bài (3-3-4 $\rightarrow$ sandhi : fán fǔbài), jiàn-shè (4-4), lián-jié (2-2).
\end{enumerate}
\end{exoresolu}

\courssub{Analyse morphologique : la famille du caractère \zh{贿} (huì)}

Comprendre la construction des caractères (clé phonétique + clé sémantique) permet de deviner le sens de mots inconnus et de mémoriser le vocabulaire par \textbf{réseaux sémantiques} plutôt que par listes isolées.

\begin{propriete}[Le caractère \zh{贿} (huì) : « bien offert pour obtenir une faveur »]
\textbf{Clé (radical) :} \zh{贝} (bèi) — \emph{coquillage, monnaie, richesse, bien de valeur}. Tous les caractères avec \zh{贝} ont un lien avec l'argent, l'échange, la valeur (ex: \zh{财} cái richesse, \zh{买} mǎi acheter, \zh{卖} mài vendre, \zh{贫} pín pauvre).
\\
\textbf{Partie phonétique :} \zh{穴} (xué) — \emph{trou, caverne} (indique la prononciation \emph{huì/xué}).
\\
\textbf{Sens originel :} Biens précieux offerts en secret (dans une « caverne ») pour acheter un pouvoir.
\\
\textbf{Famille lexicale (mots de la leçon) :}
\begin{itemize}
    \item \zh{\textbf{贿}赂 huìlù} : Pot-de-vin (Nom/Verbe). \zh{赂} (lù) = cadeau, don.
    \item \zh{受\textbf{贿} shòuhuì} : Accepter un pot-de-vin. \zh{受} (shòu) = recevoir.
    \item \zh{行\textbf{贿} xínghuì} : Offrir un pot-de-vin. \zh{行} (xíng) = accomplir, mettre en œuvre.
\end{itemize}
\textbf{Autres mots utiles (HSK 4+) :} \zh{贿选} (huìxuǎn) achat de votes, \zh{行贿罪} (xínghuì zuì) délit de corruption active.
\end{propriete}

\begin{experience}[Atelier « Arbre lexical » — Travail en binôme]
\textbf{Consigne :} Sur une feuille A3, dessinez un arbre.
\begin{enumerate}
    \item \textbf{Racines (Clé) :} Écrivez \zh{贝} (bèi) et son sens « argent/valeur ».
    \item \textbf{Tronc (Caractère pivot) :} Écrivez \zh{贿} (huì) au centre.
    \item \textbf{Grosses branches (Mots connus) :} \zh{贿赂}, \zh{受贿}, \zh{行贿}. Reliez-les au tronc. Écrivez le pinyin et la traduction.
    \item \textbf{Petites branches (Mots apparentés avec \zh{贝}) :} Ajoutez \zh{财} (richesse), \zh{贫} (pauvre), \zh{账} (compte/facture), \zh{赚} (gagner de l'argent), \zh{赔} (perdre/indemniser).
    \item \textbf{Fruits (Phrases) :} Accrochez une phrase exemple par mot connu.
\end{enumerate}
\textbf{Restitution :} Chaque binôme présente une branche à la classe en expliquant le lien sens/clé.
\end{experience}

\courssub{Focus grammaire : distinction \zh{贪污} vs \zh{腐败} vs \zh{贿赂}}

C'est le \textbf{point de grammaire lexicale} le plus important de cette leçon. Les francophones ont un seul mot « corruption » pour couvrir trois réalités juridiques et sémantiques distinctes en chinois. Ne les confondez plus.

\begin{attention}[Piège majeur : « Corruption » $\neq$ \zh{腐败} unique]
\begin{center}
\begin{tabularx}{\linewidth}{|>{\centering\arraybackslash}X|>{\centering\arraybackslash}X|>{\centering\arraybackslash}X|X|}
\hline
\rowcolor{poporangeL} \textbf{Mot chinois} & \textbf{Nature} & \textbf{Sens précis (Droit / Usage)} & \textbf{Sujet typique} \\ \hline
\zh{贪污 tānwū} & \textbf{Verbe} & \textbf{Détournement de fonds publics} : un fonctionnaire s'approprie l'argent de l'État dont il a la garde. & Fonctionnaire \zh{官员}, Comptable \zh{会计} \\ \hline
\zh{受贿 shòuhuì} & \textbf{Verbe} & \textbf{Corruption passive} : accepter de l'argent/avantages \textbf{pour} user de son pouvoir au profit du donneur. & Fonctionnaire \zh{官员}, Juge \zh{法官} \\ \hline
\zh{行贿 xínghuì} & \textbf{Verbe} & \textbf{Corruption active} : offrir de l'argent/avantages \textbf{à} un fonctionnaire pour obtenir un traitement de faveur. & Entrepreneur \zh{老板}, Citoyen \zh{市民} \\ \hline
\zh{贿赂 huìlù} & Verbe / Nom & Terme général : « corrompre » (v.) ou « pot-de-vin » (n.). Englober les deux côtés. & Les deux parties \\ \hline
\zh{腐败 fǔbài} & Nom / Adj. & \textbf{État, phénomène, qualificatif} : « la corruption » (fléau social), « corrompu » (personne/système). & Système, phénomène, personne \\ \hline
\end{tabularx}
\end{center}
\textbf{Règle d'or :} \zh{贪污} = \emph{vol dans la caisse}. \zh{受贿/行贿} = \emph{échange argent $\leftrightarrow$ pouvoir}. \zh{腐败} = \emph{nom du fléau} ou \emph{adjectif}.
\end{attention}

\begin{exemplebox}[Cas pratiques : choisissez le bon verbe]
\begin{enumerate}
    \item Un directeur d'école utilise l'argent de la cantine pour s'acheter une voiture.
    \zh{校长\textbf{贪污}了食堂经费。} (Xiàozhǎng \textbf{tānwū} le shítáng jīngfèi.) $\rightarrow$ \textbf{Détournement.}
    \item Un entrepreneur donne 50 000 FCFA à un agent des douanes pour passer sans contrôle.
    \zh{老板\textbf{行贿}海关官员。} (Lǎobǎn \textbf{xínghuì} hǎiguān guānyuán.) $\rightarrow$ \textbf{Corruption active.}
    \item L'agent des douanes prend l'argent et laisse passer.
    \zh{海关官员\textbf{受贿}放行。} (Hǎiguān guānyuán \textbf{shòuhuì} fàngxíng.) $\rightarrow$ \textbf{Corruption passive.}
    \item Ce phénomène pourrit la société.
    \zh{\textbf{腐败}现象很严重。} (\textbf{Fǔbài} xiànxiàng hěn yánzhòng.) $\rightarrow$ \textbf{Nom : le fléau.}
    \item Ce policier est pourri.
    \zh{这个警察很\textbf{腐败}。} (Zhège jǐngchá hěn \textbf{fǔbài}.) $\rightarrow$ \textbf{Adjectif : corrompu.}
\end{enumerate}
\end{exemplebox}

\courssub{Élément socioculturel : Cameroun — Chine, même combat}

\begin{savaistu}[Les institutions anti-corruption : CONAC vs 中央纪委]
La lutte contre la corruption est une priorité affichée dans les deux pays, avec des institutions dédiées puissantes.

\begin{itemize}
    \item \textbf{Cameroun :} La \textbf{CONAC} (\emph{Commission Nationale Anti-Corruption}), créée en 2006, rattachée à la Présidence. Elle a pour mission la prévention, l'éducation et la répression. Elle publie des rapports annuels et gère le numéro vert \textbf{1517}.
    \item \textbf{Chine :} La \textbf{中央纪委} (Zhōngyāng Jìwěi) — \emph{Commission centrale d'inspection disciplinaire} du PCC. C'est l'organe suprême de contrôle interne du Parti. Depuis 2012, la campagne \zh{打虎拍蝇} (dǎhǔ pāiying) — « frapper les tigres (hauts cadres) et écraser les mouches (petits fonctionnaires) » — a sanctionné des millions de cadres.
    \item \textbf{Point commun :} Les deux pays ont criminalisé la corruption transnationale et coopèrent dans le cadre de l'ONU (Convention de Mérida) et du Forum sur la Coopération Sino-Africaine (FOCAC).
    \item \textbf{Vocabulaire lié :} \zh{举报} (jǔbào) signaler/dénoncer ; \zh{零容忍} (líng róngrěn) tolérance zéro ; \zh{双规} (shuāngguī) « double règlement » (mesure d'enquête disciplinaire interne au Parti, spécifique à la Chine).
\end{itemize}
\end{savaistu}

\courssub{Méthode de mémorisation : regrouper par familles (radicaux / phonétiques)}

\begin{methode}[Comment enrichir son lexique durablement ? — 4 étapes]
\textbf{Étape 1 — Identifier le « caractère-pivot » (clé phonétique).}
Ex: \zh{贿} (huì) dans \zh{贿赂, 受贿, 行贿}. \zh{污} (wū) dans \zh{贪污, 污染, 污水}.
\textbf{Étape 2 — Construire l'arbre (cf. Atelier ci-dessus).}
Relier sens du radical (\zh{贝}=argent, \zh{氵}=eau/salissure) $\rightarrow$ sens du caractère $\rightarrow$ sens du mot.
\textbf{Étape 3 — Mémoriser les \textbf{couples opposés} (antonymes).}
\begin{itemize}
    \item \zh{腐败} (corrompu) $\leftrightarrow$ \zh{廉洁} (intègre)
    \item \zh{贪污} (détourner) $\leftrightarrow$ \zh{奉公守法} (fènggōng shǒufǎ — servir le public et respecter la loi)
    \item \zh{行贿/受贿} (corrompre/être corrompu) $\leftrightarrow$ \zh{拒绝贿赂} (refuser les pots-de-vin)
\end{itemize}
\textbf{Étape 4 — Contextualiser par l'écriture manuscrite.}
Écrire 3 fois chaque mot \textbf{en respectant l'ordre des traits} (voir section 4) tout en le prononçant à voix haute. La mémoire motrice + auditive + visuelle = ancrage long terme.
\end{methode}

\courssub{Texte d'entraînement complémentaire : « Le serment de Xiao Li »}

Ce texte original réinvestit \textbf{15 des 16 mots} du glossaire (seul \zh{举报} n'y figure pas). Lisez-le à haute voix, soulignez les mots du vocabulaire, puis répondez aux questions.

\begin{textebox}[Texte original — Niveau HSK 3/4]
小李是一名年轻的公务员，在喀麦隆首都雅温得工作。他深知\zh{腐败} (fǔbài) 是社会的毒药，\zh{贪污} (tānwū) 和 \zh{受贿} (shòuhuì) 会毁掉一个人的前途。去年，一位商人想\zh{行贿} (xínghuì) 小李，送给他一辆豪车和一大笔\zh{贿赂} (huìlù)，希望他帮忙拿到一个工程合同。小李毫不犹豫地\zh{拒绝} (jùjué) 了，并严词\zh{拒绝} (jùjué) 了对方的\zh{利益} (lìyì) 诱惑。他说：「做\zh{官员} (guānyuán)，就要\zh{廉洁} (liánjié)、\zh{诚实} (chéngshí)，对得起\zh{法律} (fǎlǜ) 和人民的\zh{责任} (zérèn)。」后来，这名商人因\zh{行贿} (xínghuì) 罪被\zh{惩罚} (chéngfá)，有关部门也加大了\zh{打击} (dǎjī) \zh{腐败} (fǔbài) 的力度。小李的事迹传开了，大家都说他是\zh{公平} (gōngpíng) 和\zh{廉洁} (liánjié) 的好干部。
\end{textebox}

\begin{exercice}[Questions de compréhension et réinvestissement lexical]
\begin{enumerate}
    \item \textbf{Vocabulaire en contexte :} Relevez dans le texte les 15 mots du glossaire utilisés et notez leur nature grammaticale \emph{dans cette phrase précise} (ex: \zh{腐败} — Nom, sujet de \zh{是}).
    \item \textbf{Grammaire :} Pourquoi \zh{拒绝} est-il utilisé deux fois ? La structure est-elle identique ? (Indice : regardez ce qui suit \zh{了} la première fois et ce qui suit \zh{了} la deuxième fois).
    \item \textbf{Traduction :} Traduisez la phrase : « 做官员，就要廉洁、诚实，对得起法律和人民的责任。 »
    \item \textbf{Expression orale (préparation) :} Imaginez la réponse du commerçant quand Xiao Li refuse. Jouez la scène en binôme (30 secondes). Utilisez au moins 3 mots du glossaire.
\end{enumerate}
\end{exercice}

\begin{corrige}[Exercice — Glossaire et prononciation]
\begin{enumerate}
    \item \textbf{Mots relevés :}
    \zh{腐败} (Nom, sujet) ; \zh{贪污} (Verbe, coordonné à \zh{受贿}) ; \zh{受贿} (Verbe) ; \zh{行贿} (Verbe, but de \zh{想}) ; \zh{贿赂} (Nom, objet de \zh{一大笔}) ; \zh{拒绝} (Verbe 1, + \zh{了} aspect accompli, objet implicite l'action de corrompre) ; \zh{拒绝} (Verbe 2, + \zh{了} aspect accompli, objet explicite \zh{诱惑}) ; \zh{利益} (Nom, attribut de \zh{诱惑}) ; \zh{官员} (Nom, sujet de \zh{就要}) ; \zh{廉洁} (Adj, prédicat, 1ʳᵉ occ.) ; \zh{诚实} (Adj, prédicat) ; \zh{法律} (Nom, objet de \zh{对得起}) ; \zh{责任} (Nom, objet de \zh{对得起}) ; \zh{惩罚} (Verbe, passif \zh{被}) ; \zh{打击} (Verbe, modifie \zh{力度}) ; \zh{腐败} (Nom, objet de \zh{打击}, 2ᵉ occ.) ; \zh{公平} (Adj, coordonné à \zh{廉洁}) ; \zh{廉洁} (Adj, attribut du sujet \zh{好干部}, 2ᵉ occ.).
    \item \textbf{Double \zh{拒绝} :} 1. \zh{拒绝了} (verbe transitif direct, objet \zh{商人} sous-entendu ou l'action de corrompre). 2. \zh{拒绝了对方的利益诱惑} (verbe transitif direct, objet explicite \zh{诱惑}). Structure identique : Suj + \zh{拒绝} + Obj + \zh{了}.
    \item \textbf{Traduction :} « Être fonctionnaire, c'est devoir être intègre et honnête, pour être digne de la loi et de la responsabilité envers le peuple. »
    \item \textbf{Oral :} Production libre attendue. Ex: « 你为什么拒绝贿赂？这是我的利益！/ 不，我要廉洁，不能违法。 »
\end{enumerate}
\end{corrige}

\coursec{Clés (radicaux) et ordre des traits}

Dans la section précédente, nous avons appris à prononcer et à employer le vocabulaire de la corruption : \zh{腐败} (fǔbài), \zh{贪污} (tānwū), \zh{贿赂} (huìlù), \zh{廉洁} (liánjié). Mais savoir prononcer un mot ne suffit pas : un bon élève de chinois doit aussi savoir \cle{l'écrire correctement} et le \cle{retrouver dans un dictionnaire}. Pour cela, il faut comprendre comment les caractères sont construits. C'est l'objet de cette section : les \cle{clés} (radicaux) et l'\cle{ordre des traits}.

\courssub{Qu'est-ce qu'une clé (radical) ?}

Observez d'abord ces quatre caractères de notre leçon : \zh{败} (bài), \zh{贪} (tān), \zh{贿} (huì), \zh{财} (cái). Que remarquez-vous ? Tous les quatre contiennent le même élément : \zh{贝}. Et tous les quatre ont un rapport avec l'argent : \zh{败} (perdre, corrompre), \zh{贪} (convoiter, être avide), \zh{贿} (soudoyer), \zh{财} (richesse). Ce n'est pas un hasard : cet élément commun est une \cle{clé}.

\begin{definition}[La clé (radical), 部首 bùshǒu]
Une \cle{clé} (en chinois \zh{部首}, bùshǒu) est un élément graphique qui sert à \textbf{classer les caractères dans le dictionnaire} et qui donne souvent un \textbf{indice de sens}. La clé \zh{贝} (bèi), qui représentait autrefois un coquillage utilisé comme monnaie, apparaît dans les caractères liés à l'argent et à la richesse. De même, la clé \zh{广} (guǎng, abri, bâtiment) apparaît dans des caractères liés aux constructions ou à des notions étendues, comme \zh{腐} (fǔ) et \zh{廉} (lián).
\end{definition}

\begin{savaistu}[Pourquoi un coquillage pour parler d'argent ?]
Dans la Chine ancienne, bien avant l'invention des pièces de monnaie, on utilisait des coquillages (cauris) comme moyen d'échange. Le caractère \zh{贝} est le dessin stylisé d'un de ces coquillages. C'est pourquoi on le retrouve dans presque tous les mots liés à l'argent : \zh{财} (cái, richesse), \zh{贵} (guì, cher /貴), \zh{买} (mǎi, acheter /買), \zh{贫} (pín, pauvre /貧), \zh{贪} (tān, avide /貪). Apprendre les clés, c'est donc aussi apprendre l'histoire de la Chine !
\end{savaistu}

\courssub{Les règles générales d'écriture des caractères}

Avant d'étudier chaque caractère un par un, il faut connaître les règles qui gouvernent l'ordre des traits. Contrairement au français, où l'on écrit les lettres de gauche à droite sans règle interne, chaque caractère chinois obéit à une \cle{chorégraphie précise} : les traits se tracent dans un ordre fixe, que les Chinois apprennent dès l'école primaire.

\begin{propriete}[Règles générales de l'ordre des traits (笔顺 bǐshùn)]
\begin{enumerate}
\item \textbf{De haut en bas} (从上到下) : on trace d'abord les traits du haut, puis ceux du bas. Exemple : \zh{三} (sān, trois) se trace du trait supérieur vers le trait inférieur.
\item \textbf{De gauche à droite} (从左到右) : on écrit d'abord la partie gauche, puis la partie droite. Exemple : \zh{贿} = \zh{贝} (gauche) puis \zh{有} (droite).
\item \textbf{L'horizontal avant le vertical} (先横后竖) quand ils se croisent : dans \zh{十} (shí, dix), on trace d'abord le trait horizontal, puis le vertical.
\item \textbf{L'extérieur avant l'intérieur} (先外后内) : pour les caractères avec enceinte, on trace d'abord l'enveloppe. Exemple : \zh{腐} = \zh{广} (l'abri, extérieur) puis l'intérieur.
\item \textbf{On ferme en dernier} (后封口) : quand un caractère est entièrement fermé (comme \zh{国} guó, pays), on trace le trait de fermeture du bas en tout dernier.
\end{enumerate}
\end{propriete}

\begin{attention}[Pourquoi respecter l'ordre des traits ?]
Beaucoup d'élèves francophones pensent : « Seul le résultat final compte, peu importe l'ordre. » C'est une erreur ! Respecter l'ordre des traits permet (1) d'écrire plus vite et plus joliment, (2) de mémoriser le caractère durablement par la mémoire motrice, (3) de compter correctement les traits pour utiliser un dictionnaire papier ou électronique, (4) d'être lisible quand on écrit vite, car l'écriture cursive (草书 cǎoshū) suit l'ordre standard. Un caractère tracé dans le désordre est souvent mal équilibré et difficile à lire pour un Chinois.
\end{attention}

\courssub{Étude des caractères du thème, un par un}

\textbf{1. \zh{腐} (fǔ) — pourri, corrompu}\par
Clé : \zh{广} (guǎng, abri). Nombre total de traits : \textbf{11} (caractère simplifié). Structure : caractère à enveloppe supérieure gauche (左上包围). On écrit d'abord \zh{广} (3 traits : point, horizontal, trait tombant gauche), puis l'intérieur \zh{府} (fǔ, 8 traits). \textbf{Étymologie} : la forme traditionnelle \zh{腐} associait \zh{肉} (viande) à \zh{府} (phonétique) ; la forme simplifiée a remplacé \zh{肉} par \zh{广} (bâtiment). L'image reste parlante : de la viande conservée dans un bâtiment finit par pourrir — d'où le sens « pourri », puis au figuré « corrompu ». On le retrouve dans \zh{腐败} (fǔbài, corruption).

\textbf{2. \zh{败} (bài) — perdre, corrompre}\par
Clé : \zh{贝} (bèi, coquillage/argent). Nombre total de traits : \textbf{8}. Structure gauche-droite : on écrit \zh{贝} à gauche (4 traits : vertical, horizontal-rejeté, trait tombant, point), puis \zh{攵} (bǔ, main tenant un bâton) à droite (4 traits : horizontal, trait tombant, vertical/droit, point). La partie droite évoque un coup porté : frapper l'argent, c'est le perdre — d'où « défaite, perte, corruption ».

\textbf{3. \zh{贪} (tān) — avide, corrompu}\par
Clé : \zh{贝} (bèi, argent). Nombre total de traits : \textbf{8}. Structure haut-bas : on écrit \zh{今} (jīn, aujourd'hui) en haut (4 traits), puis \zh{贝} en bas (4 traits), de haut en bas. \textbf{Mnémotechnique} : vouloir l'argent (\zh{贝}) \emph{aujourd'hui} (\zh{今}), tout de suite, sans attendre — c'est la définition même de la cupidité ! On le retrouve dans \zh{贪污} (tānwū, détournement de fonds).

\textbf{4. \zh{贿} (huì) — soudoyer}\par
Clé : \zh{贝} (bèi, argent). Nombre total de traits : \textbf{10}. Structure gauche-droite : \zh{贝} à gauche (4 traits), puis \zh{有} (yǒu, avoir) à droite (6 traits). \textbf{Mnémotechnique} : « avoir (\zh{有}) de l'argent (\zh{贝}) » à offrir pour corrompre quelqu'un. On le retrouve dans \zh{贿赂} (huìlù, pot-de-vin, corruption active).

\textbf{5. \zh{廉} (lián) — honnête, intègre}\par
Clé : \zh{广} (guǎng, abri). Nombre total de traits : \textbf{13}. Comme \zh{腐}, c'est un caractère à enveloppe supérieure gauche : on écrit d'abord \zh{广} (3 traits), puis \zh{兼} (jiān, 10 traits) à l'intérieur. C'est le mot de l'honnêteté : \zh{廉洁} (liánjié, intègre, incorruptible), \zh{廉政} (liánzhèng, gouvernement intègre). Remarquez que \zh{腐} et \zh{廉}, les deux contraires, partagent la même clé \zh{广} : la morale se joue sous le même toit !

\begin{center}
\begin{tabularx}{\linewidth}{|X|X|X|X|}
\hline
\rowcolor{popblueL} Caractère & Clé & Nombre de traits & Ordre d'écriture principal \\
\hline
\zh{腐} fǔ (corrompu) & \zh{广} (abri) & 11 & \zh{广} d'abord, puis \zh{府} à l'intérieur \\
\hline
\zh{败} bài (perdre) & \zh{贝} (argent) & 8 & \zh{贝} à gauche, puis \zh{攵} à droite \\
\hline
\zh{贪} tān (avide) & \zh{贝} (argent) & 8 & \zh{今} en haut, puis \zh{贝} en bas \\
\hline
\zh{贿} huì (soudoyer) & \zh{贝} (argent) & 10 & \zh{贝} à gauche, puis \zh{有} à droite \\
\hline
\zh{廉} lián (intègre) & \zh{广} (abri) & 13 & \zh{广} d'abord, puis \zh{兼} à l'intérieur \\
\hline
\zh{财} cái (richesse) & \zh{贝} (argent) & 7 & \zh{贝} à gauche, puis \zh{才} à droite \\
\hline
\end{tabularx}
\end{center}

\begin{aretenir}[L'indice sémantique de la clé \zh{贝}]
La clé \zh{贝} (bèi, coquillage/argent) revient dans \zh{败}, \zh{贪}, \zh{贿} et \zh{财}. Ce n'est pas une coïncidence : tous ces mots sont liés à \textbf{l'argent}. Quand vous rencontrez un caractère inconnu contenant \zh{贝}, vous pouvez deviner qu'il a un rapport avec l'argent, la richesse, l'achat, la valeur ou la perte. La clé est votre première aide à la devinette du sens — un outil que le français ne possède pas, car nos lettres n'ont pas de signification propre.
\end{aretenir}

\begin{popfigure}
\begin{tikzpicture}[scale=0.95, every node/.style={font=\small}]
% Grille 田字格
\draw[popline, thick] (0,0) rectangle (4,4);
\draw[popline, dashed] (2,0) -- (2,4);
\draw[popline, dashed] (0,2) -- (4,2);
\node[popmuted, below] at (2,-0.2) {Case 田字格 (tiánzìgé)};
% Partie gauche 贝 (4 traits)
\draw[popblue, very thick] (0.6,3.5) -- (0.6,1.2); % 1 Vertical
\node[popblue, font=\footnotesize\bfseries] at (0.3,3.6) {1};
\draw[popblue, very thick] (0.6,3.5) -- (1.6,3.5) -- (1.6,1.2); % 2 Horizontal + turn down
\node[popblue, font=\footnotesize\bfseries] at (1.8,3.6) {2};
\draw[popblue, very thick] (1.1,1.2) -- (0.7,0.5); % 3 Pie left-falling
\node[popblue, font=\footnotesize\bfseries] at (1.25,0.7) {3};
\draw[popblue, very thick] (1.1,1.2) -- (1.5,0.5); % 4 Dian right-falling (dot)
\node[popblue, font=\footnotesize\bfseries] at (1.65,0.7) {4};
\node[popblue, below] at (1.1,0.1) {贝};
% Partie droite 有 (6 traits)
\draw[poporange, very thick] (2.4,3.5) -- (3.7,3.5); % 5 Top horizontal
\node[poporange, font=\footnotesize\bfseries] at (2.2,3.7) {5};
\draw[poporange, very thick] (3.5,3.5) -- (2.6,2.4); % 6 Pie from top right
\node[poporange, font=\footnotesize\bfseries] at (3.7,3.7) {6};
\draw[poporange, very thick] (2.7,2.4) -- (2.7,0.6); % 7 Vertical (left of moon)
\node[poporange, font=\footnotesize\bfseries] at (2.45,2.6) {7};
\draw[poporange, very thick] (2.7,2.4) -- (3.7,2.4) -- (3.7,0.6); % 8 Horizontal + Vertical (fold) - top & right of moon
\node[poporange, font=\footnotesize\bfseries] at (3.9,2.5) {8};
\draw[poporange, very thick] (2.7,1.7) -- (3.6,1.7); % 9 Middle horizontal
\node[poporange, font=\footnotesize\bfseries] at (3.85,1.8) {9};
\draw[poporange, very thick] (2.7,0.6) -- (3.6,0.6); % 10 Bottom horizontal
\node[poporange, font=\footnotesize\bfseries] at (3.85,0.7) {10};
\node[poporange, below] at (3.15,0.1) {有};
% Résultat
\node[font=\Large] at (5.2,2) {= \zh{贿}};
\node[popdark, font=\small] at (5.2,1.3) {huì (4\textsuperscript{e} ton)};
\node[popdark, font=\small] at (5.2,0.9) {« soudoyer »};
\end{tikzpicture}
\legende{Décomposition trait par trait du caractère \zh{贿} (10 traits) : d'abord \zh{贝} à gauche (traits 1 à 4), puis \zh{有} à droite (traits 5 à 10), respectant la règle « gauche avant droite » et « haut avant bas » à l'intérieur de chaque composant.}
\end{popfigure}

\courssub{Attention aux caractères proches}

\begin{attention}[Pièges fréquents pour les francophones]
\begin{itemize}
\item Ne confondez pas \zh{贝} (bèi, coquillage/argent) et \zh{见} (jiàn, voir). Ils se ressemblent beaucoup, mais \zh{贝} se termine par un point final (丶) qui part vers la droite en bas à droite, tandis que \zh{见} se termine par un crochet arrondi (乙) en bas à droite. \zh{贿} (soudoyer) contient \zh{贝}, pas \zh{见}.
\item \zh{贪} (tān, avide) s'écrit avec \zh{今} (jīn, aujourd'hui) en haut, \textbf{pas} avec \zh{令} (lìng, ordonner). La différence : \zh{今} n'a pas de point en bas (dernier trait = horizontal), \zh{令} en a un (dernier trait = point). Un point de trop, et votre caractère n'existe plus !
\item Dans \zh{败} et \zh{贿}, la partie \zh{贝} à gauche est \textbf{plus étroite} que le \zh{贝} écrit seul : les composants se compressent horizontalement pour laisser de la place à la partie droite. C'est une règle d'esthétique et d'équilibre de l'écriture chinoise (结构 jiégòu).
\end{itemize}
\end{attention}

\courssub{Méthode : retrouver un mot dans le dictionnaire}

\begin{methode}[Chercher un caractère inconnu grâce à sa clé (查字典 cházìdiǎn)]
Vous lisez un texte et rencontrez \zh{贿}, que vous ne connaissez pas. Voici la démarche en quatre étapes :
\begin{enumerate}
\item \textbf{Identifier la clé} : cherchez l'élément qui donne le sens, souvent à gauche (structure gauche-droite) ou en haut (structure haut-bas). Dans \zh{贿}, c'est \zh{贝} (argent).
\item \textbf{Compter les traits de la clé} : \zh{贝} compte 4 traits. Repérez la clé \zh{贝} dans l'index des radicaux (souvent classés par nombre de traits).
\item \textbf{Compter les traits restants} : après avoir retiré la clé, comptez les traits de la partie restante (le « corps » du caractère). Ici, \zh{有} compte 6 traits.
\item \textbf{Retrouver la page et vérifier} : dans l'index, sous la clé 4 traits \zh{贝}, cherchez la liste des caractères à 6 traits restants. Le dictionnaire vous donne la page où \zh{贿} est défini : huì, « soudoyer, corrompre ». Vérifiez le sens dans le contexte de votre phrase.
\end{enumerate}
Cette méthode, utilisée par tous les écoliers chinois, vous rend autonome : plus besoin d'attendre le professeur pour comprendre un caractère nouveau !
\end{methode}

\begin{exoresolu}[Identifier la clé et compter les traits pour le dictionnaire]
Énoncé : pour chaque caractère suivant, indiquez la clé, son nombre de traits, le nombre de traits de la partie restante, et le total, comme pour une recherche au dictionnaire : \zh{腐}, \zh{贪}, \zh{财}.
\tcblower
Solution détaillée :
\begin{itemize}
\item \zh{腐} (fǔ) : la clé est l'enveloppe \zh{广} (3 traits). La partie restante (l'intérieur \zh{府}) compte 8 traits. Total : 3 + 8 = \textbf{11 traits}.
\item \zh{贪} (tān) : la clé est \zh{贝} en bas (4 traits). La partie restante \zh{今} en haut compte 4 traits. Total : 4 + 4 = \textbf{8 traits}.
\item \zh{财} (cái) : la clé est \zh{贝} à gauche (4 traits). La partie restante \zh{才} à droite compte 3 traits. Total : 4 + 3 = \textbf{7 traits}.
\end{itemize}
Dans les trois cas, la clé donne déjà un indice de sens : \zh{广} renvoie à un lieu ou une notion étendue (pour \zh{腐}), \zh{贝} renvoie à l'argent (pour \zh{贪} et \zh{财}).
\end{exoresolu}

\begin{exercice}[À vous d'écrire !]
Dans votre cahier, préparez des cases 田字格 (tiánzìgé, un grand carré divisé en quatre par un trait pointillé horizontal et un trait pointillé vertical). Tracez chacun des cinq caractères suivants \textbf{cinq fois}, en respectant scrupuleusement l'ordre des traits vu dans la leçon : \zh{腐}, \zh{败}, \zh{贪}, \zh{贿}, \zh{廉}. Puis, pour chaque caractère, écrivez à côté : la clé, le nombre total de traits, et un mot de la leçon qui l'emploie (exemple : \zh{贿} → \zh{贿赂}, huìlù, pot-de-vin).
\end{exercice}

\begin{corrige}[Exercice d'écriture]
Vérifiez votre travail avec le tableau récapitulatif de la leçon :
\begin{itemize}
\item \zh{腐} : clé \zh{广}, 11 traits, mot \zh{腐败} (fǔbài).
\item \zh{败} : clé \zh{贝}, 8 traits, mot \zh{腐败} (fǔbài).
\item \zh{贪} : clé \zh{贝}, 8 traits, mot \zh{贪污} (tānwū).
\item \zh{贿} : clé \zh{贝}, 10 traits, mot \zh{贿赂} (huìlù).
\item \zh{廉} : clé \zh{广}, 13 traits, mot \zh{廉洁} (liánjié).
\end{itemize}
\textbf{Points de contrôle} : l'enveloppe \zh{广} est tracée en premier dans \zh{腐} et \zh{廉} (règle « extérieur avant intérieur ») ; \zh{贪} a bien \zh{今} (sans point final) en haut ; dans \zh{败} et \zh{贿}, le \zh{贝} de gauche est étroit et laisse la place à la partie droite. Si un de vos caractères est déséquilibré, recommencez en suivant les pointillés de la case 田字格 : chaque composant doit occuper sa zone (moitié gauche/droite, tiers haut/milieu/bas) sans déborder.
\end{corrige}

\begin{aretenir}[L'essentiel de la section]
\begin{itemize}
\item Une \cle{clé} (radical, 部首 bùshǒu) classe le caractère dans le dictionnaire et donne un indice de sens.
\item La clé \zh{贝} (argent) unit \zh{败}, \zh{贪}, \zh{贿}, \zh{财} : tous ces mots parlent d'argent.
\item L'ordre des traits (笔顺 bǐshùn) obéit à des règles fixes : haut→bas, gauche→droite, horizontal→vertical, extérieur→intérieur, fermeture finale.
\item Pour chercher un caractère inconnu : (1) identifier la clé, (2) compter ses traits, (3) compter les traits restants, (4) consulter l'index des radicaux.
\end{itemize}
\end{aretenir}

\coursec{Questions de compréhension et approfondissement}

Après avoir analysé l'écriture et la structure des caractères clés du thème, nous passons maintenant à l'exploitation approfondie du texte support. Cette section vise à vérifier votre compréhension fine, à vous entraîner à formuler des réponses complètes en chinois (compétence écrite et orale) et à préparer l'argumentation nécessaire pour l'épreuve orale du baccalauréat.

\courssub{Rappel du texte support et stratégie de lecture}

Avant de répondre aux questions, relisons le texte ensemble. La stratégie recommandée : \textbf{1.} Identifier les mots-clés (mots-outils, connecteurs logiques comme \zh{因为、所以、但是、不仅……而且……}). \textbf{2.} Repérer les segments qui répondent directement aux questions \zh{谁、什么、怎么、为什么}. \textbf{3.} Reformuler la réponse avec ses propres mots en réutilisant le vocabulaire du glossaire.

\begin{textebox}[Texte support : 反腐败，人人有责]
腐败是一个严重的社会问题，它破坏公平正义，让老百姓失去对政府的信任。在中国，反腐败斗争已经持续很多年。中央政府采取了非常有力的措施：完善法律法规，加强监督检查，严厉惩治贪污受贿。无论是高官「老虎」，还是基层小吏「苍蝇」，一律依法查处，这就是所谓的「打虎拍蝇」。\\
Rán'ér, fǎnfǔbài bù zhǐshì zhèngfǔ de shìqing, yěshì quán shèhuì de zérèn. Lǎobǎixìng xīwàng shēnghuó zài yí gè qīnglián, gōngzhèng de huánjìng lǐ. Rúguǒ měi ge rén dōu cóng xiǎoshì zuòqǐ, jùjué xínghùi shòuhuì, gǎnyú jǔbào wéifǎ xíngwéi, fēngqì jiù huì hǎozhuǎn. Zuòwéi xuéshēng, wǒmen suīrán niánqīng, dàn yě néng gòngxiàn lìliàng: kǎoshì bù zuòbì, jìngxuǎn bù lāpiào, zuò chéngshí shǒuxìn de xiǎo gōngmín. Cóng xiǎoshì zuòqǐ, cóng shēnbiān zuòqǐ, qīnglián Zhōngguó cái néng zhēnzhèng shíxiàn.\\
\textbf{Pinyin :} Fǔbài shì yí gè yánzhòng de shèhuì wèntí, tā pòhuài gōngpíng zhèngyì, ràng lǎobǎixìng shīqù duì zhèngfǔ de xìnrèn. Zài Zhōngguó, fǎnfǔbài dòuzhēng yǐjīng chìxù le hěnduō nián. Zhōngyāng zhèngfǔ cǎiqǔ le fēicháng yǒulì de cuòshī: wánshàn fǎlǜ fǎguī, jiāqiáng jiāndū jiǎnchá, yánlì chéngzhì tānwū shòuhuì. Wúlùn shì gāoguān ``lǎohǔ'', háishì jīcéng xiǎolì ``cāngying'', yīlǜ yīfǎ cháchǔ, jiù shì suǒwèi de ``dǎhǔ pāiying''.\\
\textbf{Traduction :} La corruption est un problème social grave ; elle détruit l'équité et la justice, faisant perdre au peuple la confiance en le gouvernement. En Chine, la lutte contre la corruption dure depuis de nombreuses années. Le gouvernement central a pris des mesures très fermes : perfectionner les lois et règlements, renforcer la supervision et l'inspection, punir sévèrement la corruption et la prise de pots-de-vin. Qu'il s'agisse de hauts fonctionnaires « tigres » ou de petits agents de base « mouches », tous sont traités selon la loi, c'est ce qu'on appelle « frapper les tigres et chasser les mouches ».\\
Cependant, lutter contre la corruption n'est pas seulement l'affaire du gouvernement, c'est la responsabilité de toute la société. Le peuple espère vivre dans un environnement intègre et juste. Si chacun commence par de petites choses, refuse de donner ou recevoir des pots-de-vin, ose signaler les actes illégaux, les mœurs s'amélioreront. En tant qu'élèves, bien que jeunes, nous pouvons aussi contribuer : ne pas tricher aux examens, ne pas acheter de votes lors d'élections, être de jeunes citoyens honnêtes et dignes de confiance. En commençant par les petites choses, par notre entourage, une Chine intègre pourra vraiment se réaliser.
\end{textebox}

\courssub{Questions de compréhension détaillée (Q1 à Q4)}

Nous traitons ici les quatre questions objectives dont les réponses se trouvent explicitement ou implicitement dans le texte. Pour chaque question, nous suivons la méthode : \textbf{Mot-clé repéré $\rightarrow$ Phrase citée $\rightarrow$ Reformulation personnelle}.

\begin{methode}[Méthode : Répondre à une question de compréhension en chinois]
\begin{enumerate}
    \item \textbf{Analyser la question} : Identifier le mot-interrogatif (\zh{什么、怎么、谁、为什么}) et le sujet.
    \item \textbf{Localiser dans le texte} : Chercher les mots-clés (synonymes, antonymes, champs lexicaux).
    \item \textbf{Citer la phrase source} : Copier le segment exact du texte (guillemets chinois \zh{「」}).
    \item \textbf{Reformuler} : Construire une phrase complète \textbf{Sujet + Prédicat} en réutilisant le vocabulaire du glossaire. Ne pas copier-coller bêtement ; adapter la structure (ex: passer du passif à l'actif, nominaliser).
\end{enumerate}
\end{methode}

\begin{exemplebox}[Q1 : 腐败是一个什么问题？]
\textbf{Analyse} : \zh{什么问题} demande une qualification/nature.
\textbf{Mots-clés repérés} : \zh{严重的社会问题、破坏公平正义、失去信任}.
\textbf{Citation} : \zh{「腐败是一个严重的社会问题，它破坏公平正义，让老百姓失去对政府的信任」.}
\textbf{Réponse modèle} :
\zh{腐败是一个严重的社会问题，因为它破坏了公平正义，导致老百姓失去对政府的信任。}
(Fǔbài shì yí gè yánzhòng de shèhuì wèntí, yīnwèi tā pòhuàile gōngpíng zhèngyì, dǎozhì lǎobǎixìng shīqù duì zhèngfǔ de xìnrèn.)
\par
\textit{La corruption est un problème social grave car elle détruit l'équité et la justice, entraînant la perte de confiance du peuple envers le gouvernement.}
\end{exemplebox}

\begin{exemplebox}[Q2 : 中国政府怎么反腐败？]
\textbf{Analyse} : \zh{怎么} demande les mesures/moyens (\zh{措施、办法}).
\textbf{Mots-clés repérés} : \zh{采取了非常有力的措施、完善法律法规、加强监督检查、严厉惩治、打虎拍蝇}.
\textbf{Citation} : \zh{「中央政府采取了非常有力的措施：完善法律法规，加强监督检查，严厉惩治贪污受贿……一律依法查处，这就是所谓的『打虎拍蝇』」.}
\textbf{Réponse modèle} :
\zh{中国政府采取了有力措施，包括完善法律法规、加强监督检查、严厉惩治贪污受贿，对「老虎」和「苍蝇」一律依法查处，坚持「打虎拍蝇」。}
(Zhōngguó zhèngfǔ cǎiqǔle yǒulì cuòshī, bāokuò wánshàn fǎlǜ fǎguī, jiāqiáng jiāndū jiǎnchá, yánlì chéngzhì tānwū shòuhuì, duì ``lǎohǔ'' hé ``cāngying'' yīlǜ yīfǎ cháchǔ, jiānchí ``dǎhǔ pāiying''.)
\par
\textit{Le gouvernement chinois a pris des mesures fermes, notamment perfectionner les lois, renforcer la supervision, punir sévèrement la corruption, traiter « tigres » et « mouches » selon la loi, en persévérant dans la stratégie « frapper les tigres, chasser les mouches ».}
\end{exemplebox}

\begin{exemplebox}[Q3 : 老百姓希望什么？]
\textbf{Analyse} : \zh{希望} + \zh{什么} (objet de l'espoir).
\textbf{Mots-clés repérés} : \zh{希望、生活在、清廉、公正、环境}.
\textbf{Citation} : \zh{「老百姓希望生活在一个清廉、公正的环境里」.}
\textbf{Réponse modèle} :
\zh{老百姓希望生活在一个清廉、公正的社会环境里。}
(Lǎobǎixìng xīwàng shēnghuó zài yí gè qīnglián, gōngzhèng de shèhuì huánjìng lǐ.)
\par
\textit{Le peuple espère vivre dans un environnement social intègre et juste.}
\end{exemplebox}

\begin{exemplebox}[Q4 : 反腐败只是政府的事情吗？请就文中观点作答。]
\textbf{Analyse} : Question fermée (\zh{吗}) + justification (\zh{就文中观点作答}). Réponse attendue : Non + argument du texte.
\textbf{Mots-clés repérés} : \zh{不只是、也是、全社会的责任、每个人、从小事做起}.
\textbf{Citation} : \zh{「反腐败不只是政府的事情，也是全社会的责任……如果每个人都从小事做起……风气就会好转」.}
\textbf{Réponse modèle} :
\zh{不是。反腐败不只是政府的事情，也是全社会、每一个公民的责任。只有人人从小事做起，拒绝贿赂、敢于举报，社会风气才会好转。}
(Bú shì. Fǎnfǔbài bù zhǐshì zhèngfǔ de shìqing, yěshì quán shèhuì, měi yí gè gōngmín de zérèn. Zhǐyǒu rénrén cóng xiǎoshì zuòqǐ, jùjué huìlù, gǎnyú jǔbào, shèhuì fēngqì cái huì hǎozhuǎn.)
\par
\textit{Non. Lutter contre la corruption n'est pas seulement l'affaire du gouvernement, c'est la responsabilité de toute la société et de chaque citoyen. Seulement si chacun commence par de petites choses, refuse les pots-de-vin et ose signaler, les mœurs sociales s'amélioreront.}
\end{exemplebox}

\begin{propriete}[Structure de la réponse argumentée (Q4)]
En chinois, pour répondre « Non, pas seulement A, mais aussi B », on utilise la structure :
\zh{不是。不只是 A，也是 B。}
\par
Pour exprimer la condition nécessaire : \zh{只有……才……} (Seulement si... alors...).
\par
\begin{center}
\begin{tabularx}{\linewidth}{|X|X|X|}
\hline
\rowcolor{popblueL} \textbf{Structure} & \textbf{Exemple (Caractères + Pinyin)} & \textbf{Traduction} \\ \hline
\zh{不是。不只是 A，也是 B。} & \zh{不是。反腐败不只是政府的事，也是大家的事。} (Bú shì. Fǎnfǔbài bù zhǐshì zhèngfǔ de shì, yěshì dàjiā de shì.) & Non. Ce n'est pas seulement l'affaire du gouv., c'est celle de tous. \\ \hline
\zh{只有……才……} & \zh{只有人人参与，反腐败才能成功。} (Zhǐyǒu rénrén cānyù, fǎnfǔbài cáinéng chénggōng.) & Seule la participation de tous permet la réussite. \\ \hline
\end{tabularx}
\end{center}
\end{propriete}

\courssub{Question personnelle et expression orale (Q5)}

La Q5 est une question ouverte (\zh{你认为……}) qui évalue votre capacité à exprimer un avis personnel structuré en réinvestissant le lexique du thème. C'est un entraînement direct à l'expression orale (Partie 2 de l'oral : exposé / interaction).

\begin{methode}[Méthode : Répondre à \zh{你认为……} (Question personnelle)]
\begin{enumerate}
    \item \textbf{Marqueur d'opinion} : \zh{我认为} (Wǒ rènwéi - Je pense que) ou \zh{我觉得} (Wǒ juéde - Je trouve que / J'ai l'impression que). \zh{我认为} est plus formel, adapté à l'écrit et l'oral soutenu.
    \item \textbf{Verbe modal d'obligation/conseil} : \zh{应该} (yīnggāi - devoir / devrait) ou \zh{可以} (kěyǐ - pouvoir / peut).
    \item \textbf{Action concrète} : Verbe + Objet (vocabulaire du thème : \zh{诚实做人、拒绝作弊、监督、举报、从小事做起}).
    \item \textbf{Justification (optionnel mais valorisant)} : \zh{因为……所以……} / \zh{这样才能……}.
\end{enumerate}
\end{methode}

\begin{exemplebox}[Q5 : 你认为学生能为反腐败做什么？]
\textbf{Modèle 1 (Niveau HSK 3 - Standard)} :
\zh{我认为学生应该从小事做起。考试不作弊，不抄袭，做诚实的人。这样才能养成良好的品德，长大后不贪污。}
(Wǒ rènwéi xuéshēng yīnggāi cóng xiǎoshì zuòqǐ. Kǎoshì bù zuòbì, bù chāoxí, zuò chéngshí de rén. Zhèyàng cáinéng yǎngchéng liánghǎo de pǐndé, zhǎngdà hòu bù tānwū.)
\par
\textit{Je pense que les élèves doivent commencer par de petites choses. Ne pas tricher aux examens, ne pas plagier, être des gens honnêtes. Ainsi seulement on cultive une bonne moralité, et adulte on ne sera pas corrompu.}

\textbf{Modèle 2 (Niveau HSK 4 - Approfondi, pour l'oral)} :
\zh{我觉得学生虽然没有权力直接查处贪官，但我们可以监督身边的不正之风。比如，竞选班干部不拉票、不送礼；发现同学买答案要敢于举报。从小树立廉洁意识，未来才能成为廉洁的公职人员。}
(Wǒ juéde xuéshēng suīrán méiyǒu quánlì zhíjiē cháchǔ tānguān, dàn wǒmen kěyǐ jiāndū shēnbiān de bùzhèng zhī fēng. Bǐrú, jìngxuǎn bàngànbu bù lāpiào, bù sònglǐ; fāxiàn tóngxué mǎi dá'àn yào gǎnyú jǔbào. Cóng xiǎo shùlì liánjié yìshí, wèilái cáinéng chéngwéi liánjié de gōngzhì rényuán.)
\par
\textit{Je trouve que bien que les élèves n'aient pas le pouvoir de sanctionner directement les fonctionnaires corrompus, nous pouvons surveiller les mauvaises pratiques autour de nous. Par ex., ne pas faire campagne/acheter des votes pour être délégué ; oser signaler si un camarade achète des réponses. Établir une conscience d'intégrité dès le jeune âge pour devenir futur fonctionnaire intègre.}
\end{exemplebox}

\begin{attention}[Piège fréquent : Réponses « telegrammiques »]
\emph{Erreur typique de francophone} : Répondre par un mot unique ou un fragment nominal comme en français (\textit{« Honnêteté. »}, \textit{« Ne pas tricher. »}).
\par
\textbf{Règle} : En chinois, une réponse à une question ouverte \textbf{doit être une phrase complète avec un sujet explicite}.
\begin{center}
\begin{tabularx}{\linewidth}{|X|X|}
\hline
\rowcolor{poporangeL} \textbf{Incorrect (Style français)} & \textbf{Correct (Phrase complète chinois)} \\ \hline
\zh{诚实。} (Chéngshí.) & \zh{我认为我们应该诚实做人。} (Wǒ rènwéi wǒmen yīnggāi chéngshí zuòrén.) \\ \hline
\zh{不作弊。} (Bù zuòbì.) & \zh{学生应该拒绝作弊。} (Xuéshēng yīnggāi jùjué zuòbì.) \\ \hline
\zh{监督。} (Jiāndū.) & \zh{我们可以监督社会风气。} (Wǒmen kěyǐ jiāndū shèhuì fēngqì.) \\ \hline
\end{tabularx}
\end{center}
\end{attention}

\courssub{Mini-débat guidé : 反腐败，谁的责任最大？政府还是公民？}

Cette activité prépare l'\textbf{Épreuve Orale (Interaction / Débat)}. L'objectif n'est pas de trancher mais d'argumenter en chinois structuré.

\begin{experience}[Activité de classe : Débat contradictoire (15 min)]
\begin{enumerate}
    \item \textbf{Préparation (5 min)} : En binôme, préparez 3 arguments pour le \textbf{Gouvernement} (Groupe A) et 3 pour les \textbf{Citoyens} (Groupe B). Utilisez le vocabulaire du glossaire (\zh{权力、法律、监督、榜样、从小事做起、举报、制度}).
    \item \textbf{Débat (7 min)} : Alternez les tours de parole. Un argument = une phrase complète avec connecteur (\zh{首先、其次、此外、但是、因此}).
    \item \textbf{Synthèse (3 min)} : Le professeur note au tableau la structure \zh{政府负责……，公民负责……，两者缺一不可。}.
\end{enumerate}
\end{experience}

\begin{aretenir}[Expressions utiles pour le débat]
\begin{center}
\begin{tabularx}{\linewidth}{|X|X|X|}
\hline
\rowcolor{poppurpleL} \textbf{Fonction} & \textbf{Chinois (Caractères + Pinyin)} & \textbf{Français} \\ \hline
Donner son avis & \zh{我认为…… / 依我看……} (Wǒ rènwéi / Yī wǒ kàn) & Je pense que / À mon avis... \\ \hline
Ajouter un argument & \zh{首先……，其次……，而且……} (Shǒuxiān, qícì, érqiě) & Premièrement..., deuxièmement..., de plus... \\ \hline
Concéder / Nuancer & \zh{政府固然重要，但是公民……} (Zhèngfǔ gǔrán zhòngyào, dànshì gōngmín...) & Le gouv. est certes important, mais les citoyens... \\ \hline
Conclure & \zh{总之，只有……才能……} (Zǒngzhī, zhǐyǒu... cáinéng...) & En résumé, seulement si... alors... \\ \hline
\end{tabularx}
\end{center}
\end{aretenir}

\courssub{Organigramme de réponse : Méthode visuelle}

Cet organigramme résume la démarche intellectuelle pour passer du texte à la réponse rédigée. Il est utilisable pour toutes les questions de compréhension écrite.

\begin{popfigure}
\begin{tikzpicture}[
    node distance=1.2cm and 1.5cm,
    every node/.style={font=\small, align=center},
    box/.style={rectangle, draw=popblue, thick, fill=popblueL, rounded corners, minimum width=2.8cm, minimum height=1.1cm},
    arrow/.style={-Stealth, thick, popdark},
    keyword/.style={rectangle, draw=poporange, fill=poporangeL, rounded corners, minimum width=2.5cm, minimum height=0.9cm},
    cite/.style={rectangle, draw=popgreen, fill=popgreenL, rounded corners, minimum width=3cm, minimum height=0.9cm},
    reform/.style={rectangle, draw=poppurple, fill=poppurpleL, rounded corners, minimum width=3cm, minimum height=0.9cm}
]
\node[box, align=center] (start){Question\\ (ex: \zh{怎么反腐败？})};
\node[keyword, below=of start, align=center] (kw){Mots-clés repérés\\ \zh{措施、法律、监督、惩治、打虎拍蝇}};
\node[cite, below=of kw, align=center] (cite){Phrase citée du texte\\ \zh{「采取了有力措施……打虎拍蝇」}};
\node[reform, below=of cite, align=center] (reform){Réponse reformulée\\ \zh{政府采取有力措施……坚持打虎拍蝇。}};
\node[box, below=of reform, align=center] (end){Réponse finale\\ (Phrase complète, vocabulaire réinvesti)};

\draw[arrow] (start) -- (kw) node[midway, right, font=\tiny, text=popdark] {1. Analyser};
\draw[arrow] (kw) -- (cite) node[midway, right, font=\tiny, text=popdark] {2. Localiser};
\draw[arrow] (cite) -- (reform) node[midway, right, font=\tiny, text=popdark] {3. Citer};
\draw[arrow] (reform) -- (end) node[midway, right, font=\tiny, text=popdark] {4. Reformuler};
\end{tikzpicture}
\legende{Organigramme : De la question à la réponse rédigée (Méthode en 4 étapes)}
\end{popfigure}

\courssub{Correction collective et réinvestissement lexical}

Au tableau, le professeur construit collectivement les réponses définitives en exigeant l'utilisation des termes précis du glossaire (voir Section 3). Voici le tableau de révision lexicale pour cette section.

\begin{center}
\begin{tabularx}{\linewidth}{|X|X|X|X|}
\hline
\rowcolor{popblueL} \textbf{Caractères} & \textbf{Pinyin} & \textbf{Nature} & \textbf{Traduction} \\ \hline
破坏 & pòhuài & verbe & détruire, nuire à \\ \hline
公平正义 & gōngpíng zhèngyì & nom & équité et justice \\ \hline
失去 & shīqù & verbe & perdre \\ \hline
信任 & xìnrèn & nom/verbe & confiance / faire confiance \\ \hline
措施 & cuòshī & nom & mesure (administrative) \\ \hline
监督 & jiāndū & verbe/nom & superviser, contrôler / supervision \\ \hline
惩治 & chéngzhì & verbe & punir, sanctionner (sévèrement) \\ \hline
一律 & yīlǜ & adv. & tous sans exception, uniformément \\ \hline
查处 & cháchǔ & verbe & enquêter et traiter (un dossier) \\ \hline
清廉 & qīnglián & adj. & intègre, honnête (administration) \\ \hline
拒绝 & jùjué & verbe & refuser, rejeter \\ \hline
举报 & jǔbào & verbe & signaler, dénoncer \\ \hline
风气 & fēngqì & nom & mœurs, atmosphère sociale \\ \hline
好转 & hǎozhuǎn & verbe & s'améliorer (situation) \\ \hline
贡献 & gòngxiàn & verbe/nom & contribuer / contribution \\ \hline
廉洁 & liánjié & adj. & intègre, incorruptible \\ \hline
\end{tabularx}
\end{center}

\begin{savaistu}[Le proverbe \zh{从小事做起} (Cóng xiǎoshì zuòqǐ)]
Ce \zh{成语} (chéngyǔ, expression idiomatique à 4 caractères) signifie « Commencer par les petites choses ». Il est omniprésent dans l'éducation civique chinoise (affiches dans les écoles, spots TV) tout comme au Cameroun (ex: « L'école pour la vie », campagnes anti-corruption au CONAC). Il reflète une vision confucéenne : la vertu se cultive dans le détail quotidien (\zh{勿以恶小而为之，勿以善小而不为} — Ne pas faire le mal sous prétexte qu'il est petit ; ne pas négliger le bien sous prétexte qu'il est petit). En classe de chinois, on l'applique à l'apprentissage : \zh{学中文，从认一个字做起} (Apprendre le chinois, commencer par reconnaître un caractère).
\end{savaistu}

\begin{exoresolu}[Exercice résolu : Reformulation de Q2]
\textbf{Énoncé} : À partir de la phrase du texte : \zh{「中央政府采取了非常有力的措施：完善法律法规，加强监督检查，严厉惩治贪污受贿」}, écrivez une réponse courte à la question \zh{中国政府如何反腐败？} en utilisant la structure \zh{通过……} (par le biais de...) et en remplaçant \zh{中央政府} par \zh{政府}.
\tcblower
\textbf{Solution détaillée} :
\zh{政府通过完善法律法规、加强监督检查、严厉惩治贪污受贿来反腐败。}
(Zhèngfǔ tōngguò wánshàn fǎlǜ fǎguī, jiāqiáng jiāndū jiǎnchá, yánlì chéngzhì tānwū shòuhuì lái fǎnfǔbài.)
\par
\textbf{Analyse} :
\begin{itemize}
    \item \zh{通过} (tōngguò) introduit le moyen/méthode (préposition).
    \item Les trois VP (Verbes + Objets) sont coordonnés par des virgules chinoises \zh{、}.
    \item \zh{来} (lái) + Verbe (\zh{反腐败}) indique le but final.
    \item Le sujet \zh{政府} est placé en début de phrase (sujet agent).
\end{itemize}
\end{exoresolu}

\begin{exercice}[Entraînement : Questions types Bac]
\begin{enumerate}
    \item Répondez en chinois (phrase complète) : \zh{为什么说「反腐败不只是政府的事情」？} (Indice : \zh{全社会的责任、人人、从小事做起}).
    \item Traduisez en chinois : « Si chaque élève refuse de tricher, les mœurs scolaires s'amélioreront. » (Vocabulaire : \zh{如果……就……、拒绝、作弊、校园风气、好转}).
    \item Complétez avec \zh{只有……才……} : \zh{\_\_\_\_\_\_ 人人 \_\_\_\_\_\_ 举报贪污，腐败现象 \_\_\_\_\_\_ 减少。} (Tous / oser signaler / phénomène de corruption / diminuer).
\end{enumerate}
\end{exercice}

\begin{corrige}[Exercice ci-dessus]
\begin{enumerate}
    \item \zh{因为反腐败是全社会的责任，需要人人从小事做起，拒绝贿赂、敢于举报，风气才会好转。}
    \item \zh{如果每个学生都拒绝作弊，校园风气就会好转。}
    \item \zh{只有 人人 敢于 举报贪污，腐败现象 才能 减少。} (\textbf{Attention} : \zh{才能} exprime la possibilité/réalisation future dans une structure \zh{只有……才……}).
\end{enumerate}
\end{corrige}

\coursec{Production guidée et réinvestissement}

Après avoir analysé le texte et répondu aux questions de compréhension, nous passons maintenant à l'appropriation active du lexique et des structures. Cette section vise à transformer la compréhension passive en compétence de production : vous allez rédiger un paragraphe d'opinion structuré, l'améliorer par la relecture entre pairs, puis créer une affiche de sensibilisation. C'est l'étape où le vocabulaire devient vraiment vôtre.

\courssub{Production écrite guidée : « Mon avis sur la corruption »}

Pour exprimer un point de vue argumenté en chinois, il faut respecter un ordre logique des idées : \textbf{constater} $\rightarrow$ \textbf{illustrer} $\rightarrow$ \textbf{proposer (niveau institutionnel)} $\rightarrow$ \textbf{proposer (niveau personnel)} $\rightarrow$ \textbf{conclure par une condition nécessaire}. La trame imposée ci-dessous suit exactement cette logique rhétorique.

\begin{propriete}[Trame rhétorique de l'argumentaire en 5 étapes]
Chaque numéro correspond à une phrase complète. Utilisez les connecteurs et le vocabulaire du thème pour remplir les pointillés.
\begin{enumerate}
    \item \textbf{Définition / Constat :} 腐败是 \underline{\hspace{4cm}} 。
    \item \textbf{Exemple concret :} 有些官员 \underline{\hspace{4cm}} 。
    \item \textbf{Solution institutionnelle :} 政府应该 \underline{\hspace{4cm}} 。
    \item \textbf{Engagement citoyen (nous étudiants) :} 我们学生应该 \underline{\hspace{4cm}} 。
    \item \textbf{Conclusion conditionnelle :} 只有 \underline{\hspace{3cm}}，\underline{\hspace{3cm}} 才 \underline{\hspace{2cm}} 。
\end{enumerate}
\end{propriete}

\begin{center}
\begin{tabularx}{\linewidth}{|X|X|X|X|}
\hline
\rowcolor{popblueL} \textbf{Connecteurs / Mots-outils} & \textbf{Pinyin} & \textbf{Nature} & \textbf{Traduction / Emploi} \\ \hline
因为 & yīnwèi & conj. & Parce que (introduit la cause, souvent en début de phrase ou après le sujet) \\ \hline
所以 & suǒyǐ & adv. & Donc, par conséquent (introduit la conséquence, en début de proposition) \\ \hline
但是 & dànshì & conj. & Mais, cependant (marque l'opposition) \\ \hline
应该 & yīnggāi & v. mod. & Devoir, il faut (se place \textbf{avant} le verbe principal) \\ \hline
只有…才… & zhǐyǒu…cái… & structure & Seulement si… alors… (condition nécessaire, focus sur la condition) \\ \hline
我认为 & wǒ rènwéi & verbe + obj. & Je pense que / À mon avis (introduit l'opinion) \\ \hline
\end{tabularx}
\end{center}

\begin{textebox}[Modèle de paragraphe d'opinion (niveau HSK 3)]
我认为腐败是一个严重的社会问题。因为有些官员收受贿赂，滥用职权，所以百姓失去信任。政府应该加强监督，严惩贪污。我们学生应该从小事做起，拒绝作弊，诚实做人。只有全社会共同努力，反腐倡廉才会成功。
\\Wǒ rènwéi fǔbài shì yí ge yánzhòng de shèhuì wèntí. Yīnwèi yǒuxiē guānyuán shōushòu huìlù, lànyòng zhíquán, suǒyǐ bǎixìng shīqù xìnrèn. Zhèngfǔ yīnggāi jiāqiáng jiāndū, yánchéng tānwū. Wǒmen xuéshēng yīnggāi cóng xiǎoshì zuòqǐ, jùjué zuòbì, chéngshí zuòrén. Zhǐyǒu quán shèhuì gòngtóng nǔlì, fǎnfū chánglián cái huì chénggōng.
\\
Je pense que la corruption est un problème social grave. Parce que certains fonctionnaires acceptent des pots-de-vin et abusent de leur pouvoir, le peuple perd confiance. Le gouvernement devrait renforcer la surveillance et punir sévèrement la corruption. Nous, étudiants, devrions commencer par de petites choses, refuser la triche, être honnêtes. Seulement si toute la société fait des efforts communs, la lutte contre la corruption et la promotion de l'intégrité réussiront.
\end{textebox}

\begin{propriete}[Structure conditionnelle nécessaire : \cle{只有…才…} (zhǐyǒu…cái…)]
Cette structure indique qu'une condition est \textbf{indispensable} pour obtenir un résultat. « 只有 » introduit la condition (sujet + verbe), « 才 » précède le verbe principal de la conséquence.
\medskip
\noindent \textbf{Schéma :} \zh{只有 + Condition, Sujet + 才 + Verbe (Résultat)}.
\medskip
\noindent \textbf{Attention :} En français, on dit « Seulement si..., ... ». En chinois, le sujet de la conséquence se place après la virgule, avant \zh{才}.
\begin{center}
\begin{tabularx}{\linewidth}{|X|X|X|}
\hline
\rowcolor{poppurpleL} \textbf{Structure} & \textbf{Exemple (Caractères + Pinyin)} & \textbf{Traduction} \\ \hline
Seule condition (sujet implicite) & 只有拒绝贿赂，社会才会公平。\newline (Zhǐyǒu jùjué huìlù, shèhuì cái huì gōngpíng.) & Seulement en refusant les pots-de-vin, la société sera équitable. \\ \hline
Condition avec sujet explicite & 只有政府严厉惩罚，腐败现象才会减少。\newline (Zhǐyǒu zhèngfǔ yánlì chéngfá, fǔbài xiànxiàng cái huì jiǎnshǎo.) & Seulement si le gouvernement punit sévèrement, le phénomène de corruption diminuera. \\ \hline
Engagement personnel & 只有我们从小事做起，风气才会变好。\newline (Zhǐyǒu wǒmen cóng xiǎoshì zuòqǐ, fēngqì cái huì biàn hǎo.) & Seulement si nous commençons par de petites choses, les mœurs s'amélioreront. \\ \hline
\end{tabularx}
\end{center}
\end{propriete}

\begin{attention}[Piège majeur : Position de \cle{应该} (yīnggāi)]
En français, « devoir » peut suivre le sujet (\emph{Je dois partir}) ou précéder l'infinitif (\emph{Il faut partir}). En chinois, \zh{应该} est un verbe modal : il se place \textbf{systématiquement AVANT le verbe d'action}, jamais à la fin de la phrase.
\begin{center}
\begin{tabularx}{\linewidth}{|X|X|X|}
\hline
\rowcolor{poporangeL} \textbf{Phrase correcte} & \textbf{Phrase incorrecte (calque français)} & \textbf{Explication} \\ \hline
我们应该拒绝贿赂。\newline (Wǒmen yīnggāi jùjué huìlù.) & 我们拒绝应该贿赂。 & \zh{应该} précède \zh{拒绝}. \\ \hline
政府应该加强监督。\newline (Zhèngfǔ yīnggāi jiāqiáng jiāndū.) & 政府加强监督应该。 & Le verbe modal précède toujours le verbe lexical. \\ \hline
\end{tabularx}
\end{center}
\emph{Astuce :} Remplacez mentalement \zh{应该} par « il faut que » + subjonctif pour sentir l'antériorité logique.
\end{attention}

\begin{methode}[Réussir sa production écrite en 4 étapes]
\begin{enumerate}
    \item \textbf{Une phrase = une idée :} Ne mélangez pas « exemple » et « solution » dans la même phrase. Respectez la trame 1 à 5.
    \item \textbf{Recyclage intelligent :} Reprenez des segments du texte support (ex: \zh{滥用职权}, \zh{诚实做人}) en changeant le sujet ou le temps.
    \item \textbf{Banque de mots :} Cochez sur votre fiche de vocabulaire au moins 5 mots utilisés (ex: \zh{腐败}, \zh{贿赂}, \zh{监督}, \zh{诚实}, \zh{共同努力}).
    \item \textbf{Vérification finale :} Relisez à voix haute pour les tons ; vérifiez l'ordre des traits des caractères clés (\zh{贿}, \zh{监}, \zh{廉}) ; validez la structure \zh{只有…才…}.
\end{enumerate}
\end{methode}

\begin{exoresolu}[Rédaction guidée corrigée]
\textbf{Énoncé :} Complétez la trame ci-dessous en utilisant le vocabulaire du thème (corruption, pots-de-vin, surveillance, honnêteté, efforts communs) et les connecteurs fournis.

\textbf{Trame à compléter :}
1. 我认为腐败是 \underline{\hspace{3cm}} 。
2. 因为有些官员 \underline{\hspace{3cm}} ，\underline{\hspace{2cm}} 。
3. 政府 \underline{\hspace{1cm}} \underline{\hspace{3cm}} 。
4. 我们学生 \underline{\hspace{1cm}} \underline{\hspace{3cm}} 。
5. \underline{\hspace{1cm}} \underline{\hspace{3cm}} ，\underline{\hspace{2cm}} \underline{\hspace{1cm}} \underline{\hspace{2cm}} 。

\tcblower
\textbf{Solution détaillée et commentée :}

1. \zh{我认为腐败是一个严重的社会问题。} (Wǒ rènwéi fǔbài shì yí ge yánzhòng de shèhuì wèntí.) \\
   \textit{Utilisation de « 我认为 » pour marquer l'opinion personnelle + adjectif « 严重 » vu au glossaire.}

2. \zh{因为有些官员收受贿赂、滥用职权，所以百姓对政府失望。} (Yīnwèi yǒuxiē guānyuán shōushòu huìlù, lànyòng zhíquán, suǒyǐ bǎixìng duì zhèngfǔ shīwàng.) \\
   \textit{Structure Cause (因为) $\rightarrow$ Conséquence (所以). Vocabulaire : \zh{收受贿赂} (accepter pots-de-vin), \zh{滥用职权} (abuser pouvoir), \zh{百姓} (peuple), \zh{失望} (déçu).}

3. \zh{政府应该加强监督，严惩贪污腐败。} (Zhèngfǔ yīnggāi jiāqiáng jiāndū, yánchéng tānwū fǔbài.) \\
   \textit{\zh{应该} + Verbe (\zh{加强}, \zh{严惩}). Vocabulaire : \zh{监督} (surveillance), \zh{贪污} (détournement/corruption).}

4. \zh{我们学生应该从小事做起，拒绝作弊，诚实做人。} (Wǒmen xuéshēng yīnggāi cóng xiǎoshì zuòqǐ, jùjué zuòbì, chéngshí zuòrén.) \\
   \textit{Transfert au contexte scolaire : \zh{作弊} (tricher), \zh{诚实做人} (être honnête dans sa conduite). Structure \zh{从...做起} (commencer par...).}

5. \zh{只有全社会共同努力，反腐倡廉才会取得胜利。} (Zhǐyǒu quán shèhuì gòngtóng nǔlì, fǎnfū chánglián cái huì qǔdé shènglì.) \\
   \textit{Structure \zh{只有…才…}. Sujet condition : \zh{全社会} (toute la société). Verbe condition : \zh{共同努力} (efforts communs). Résultat : \zh{反腐倡廉} (lutte anti-corruption / promotion intégrité - chengyu courant) + \zh{才会} + \zh{取得胜利} (remporter la victoire).}
\end{exoresolu}

\courssub{Relecture par les pairs et auto-évaluation}

L'écriture n'est pas finie tant qu'elle n'a pas été relue. En classe, nous utilisons une grille de relecture croisée (Pair Review) pour développer l'esprit critique et la métacognition linguistique.

\begin{experience}[Atelier « Correcteur d'un jour » (20 min)]
\begin{enumerate}
    \item \textbf{Échange :} Donnez votre brouillon à votre voisin. Vous recevez le sien.
    \item \textbf{Mission 1 - Sons et Tons :} Surlignez en \textcolor{poporange}{orange} les mots dont le pinyin semble douteux (tons 2/3 confondus, \zh{ü} mal prononcé). Lisez la phrase à voix haute à l'auteur.
    \item \textbf{Mission 2 - Caractères et Ordre des traits :} Cochez en \textcolor{popgreen}{vert} les 5 caractères « thème » utilisés. Vérifiez sur la fiche « Clés et ordre des traits » (Section 4) si l'auteur a respecté l'ordre (ex: \zh{贿} : \zh{贝} en bas, \zh{监} : \zh{皿} en bas).
    \item \textbf{Mission 3 - Richesse lexicale :} Comptez les mots du glossaire réinvestis (objectif : au moins 5). Notez le score en bas de page.
    \item \textbf{Mission 4 - Structure :} Entourez \zh{应该} (vérifiez position) et \zh{只有…才…} (vérifiez présence de \zh{才}).
    \item \textbf{Retour :} Rendez la copie avec un commentaire bienveillant : « J'ai aimé ta phrase 3 car... / Attention au ton de \zh{贿} (huì 4e ton) / Il manque \zh{才} à la phrase 5. »
    \item \textbf{Réécriture :} Corrigez votre version finale (propre) pour l'affichage.
\end{enumerate}
\end{experience}

\begin{aretenir}[Check-list de révision finale (à coller dans le cahier)]
\begin{itemize}
    \item [ ] 5 phrases complètes, ponctuation chinoise (。).
    \item [ ] Au moins 5 mots du glossaire thème « Corruption » utilisés correctement.
    \item [ ] \zh{应该} placé AVANT le verbe à chaque occurrence.
    \item [ ] Structure \zh{只有…才…} correcte (condition + \zh{才} + résultat).
    \item [ ] Connecteurs logiques : \zh{因为/所以}, \zh{但是}, \zh{我认为}.
    \item [ ] Tons du pinyin vérifiés à l'oral (auto + pair).
    \item [ ] Ordre des traits respecté pour les caractères nouveaux (\zh{贿}, \zh{监}, \zh{廉}, \zh{惩}, \zh{滥}).
\end{itemize}
\end{aretenir}

\courssub{Extension : Affiche de classe bilingue — Slogan anti-corruption}

La production écrite trouve son aboutissement dans une tâche finale authentique : la communication publique. Une affiche (标语, \zh{biāoyǔ}) doit être \textbf{courte, percutante, rythmée} (souvent 4+4 ou 5+5 caractères) et \textbf{bilingue} pour l'environnement scolaire camerounais.

\begin{exemplebox}[Exemples de slogans (标语) — Structure 4+4 ou 5+5]
\begin{center}
\begin{tabularx}{\linewidth}{|X|X|X|}
\hline
\rowcolor{popgreenL} \textbf{Chinois (Caractères + Pinyin)} & \textbf{Français} & \textbf{Analyse structurelle} \\ \hline
拒绝贿赂，诚实做人！\newline (Jùjué huìlù, chéngshí zuòrén!) & Refuse les pots-de-vin, sois honnête ! & Impératif + Verbe (V+O) // Adverbe + Verbe (Adv+V). Parallélisme parfait. \\ \hline
权为民所用，利为民所谋。\newline (Quán wéi mín suǒ yòng, lì wéi mín suǒ móu.) & Le pouvoir sert le peuple, le profit profite au peuple. & Structure passive classique \zh{为...所...} (être utilisé par...). Style formel, solennel. \\ \hline
反腐倡廉，从我做起。\newline (Fǎnfū chánglián, cóng wǒ zuòqǐ.) & Lutte contre la corruption, promeus l'intégrité, commence par moi. & Chengyu (4 caractères) + Structure \zh{从...做起}. Engagement personnel fort. \\ \hline
清正廉洁，光明磊落。\newline (Qīngzhèng liánjié, guāngmíng lěiluò.) & Intègre et honnête, droit et franc. & Deux chengyu (idiomes 4 caractères) synonymes. Qualité littéraire, vocabulaire HSK 4+. \\ \hline
\end{tabularx}
\end{center}
\end{exemplebox}

\begin{experience}[Création de l'affiche collective (30 min + affichage)]
\begin{enumerate}
    \item \textbf{Groupes de 4 :} Chaque groupe propose 3 slogans originaux (chinois + français) sur papier A3.
    \item \textbf{Contraintes :} Minimum 8 caractères chinois. Utiliser au moins 2 mots du glossaire. Traduction française fidèle.
    \item \textbf{Design :} Caractères grands (feutre noir/rouge), pinyin petit (crayon gris) sous chaque caractère, français en bas.
    \item \textbf{Vote :} Galerie marchande. Chaque élève pose 2 gommettes (\"Meilleur message\", \"Meilleur chinois\"). Les 3 meilleures sont plastifiées pour le couloir du bâtiment des langues.
    \item \textbf{Exemples d'idées pour démarrer :}
    \begin{itemize}
        \item \zh{一笔贿赂，毁了前程。} (Un pot-de-vin, un avenir détruit.) — Focus conséquence.
        \item \zh{清廉是福，贪婪是祸。} (L'intégrité est une bénédiction, la cupidité un malheur.) — Parallélisme philosophique.
        \item \zh{不拿一针一线，还我清清白白。} (Ne pas prendre un fil ni une aiguille, me rendre ma pureté.) — Référence culturelle (Zhu Rongji).
    \end{itemize}
\end{enumerate}
\end{experience}

\begin{popfigure}
\begin{tikzpicture}[
    node distance=1.2cm and 1.5cm,
    box/.style={draw=popblue, thick, rounded corners, fill=popblueL, text width=2.8cm, align=center, minimum height=1.8cm, font=\small},
    arrow/.style={-Stealth, thick, popdark},
    labelbox/.style={font=\tiny, text=popmuted, align=center}
]
% Nœuds principaux
\node[box, align=center] (opinion){\textbf{1. Opinion}\\\zh{我认为腐败是...}\\(Je pense que la corruption est...)};
\node[box, right=of opinion, align=center] (exemple){\textbf{2. Exemple}\\\zh{因为有些官员...所以...}\\(Parce que certains fonctionnaires... donc...)};
\node[box, right=of exemple, align=center] (gov){\textbf{3. Solution Gov}\\\zh{政府应该...}\\(Le gouvernement devrait...)};
\node[box, right=of gov, align=center] (citoyen){\textbf{4. Solution Citoyenne}\\\zh{我们学生应该...}\\(Nous étudiants devrions...)};
\node[box, right=of citoyen, align=center] (conclusion){\textbf{5. Conclusion}\\\zh{只有...才...}\\(Seulement si... alors...)};

% Flèches
\draw[arrow] (opinion) -- (exemple) node[midway, above, labelbox] {Constater $\rightarrow$ Illustrer};
\draw[arrow] (exemple) -- (gov) node[midway, above, labelbox] {Cause $\rightarrow$ Remède institutionnel};
\draw[arrow] (gov) -- (citoyen) node[midway, above, labelbox] {Haut $\rightarrow$ Bas (Citoyen)} ;
\draw[arrow] (citoyen) -- (conclusion) node[midway, above, labelbox] {Action $\rightarrow$ Condition nécessaire} ;

% Annotations clés sous les boîtes
\node[labelbox, below=0.2cm of opinion] {\zh{我认为 / 严重 / 问题}};
\node[labelbox, below=0.2cm of exemple] {\zh{因为 / 所以 / 贿赂 / 滥用职权}};
\node[labelbox, below=0.2cm of gov] {\zh{应该 + V / 监督 / 惩罚}};
\node[labelbox, below=0.2cm of citoyen] {\zh{应该 + V / 诚实 / 拒绝}};
\node[labelbox, below=0.2cm of conclusion] {\zh{只有 / 才 / 共同努力 / 成功}};

% Titre global
\node[above=0.5cm of opinion, font=\bfseries\large, text=popblue] {Schéma rhétorique de l'argumentaire anti-corruption};
\end{tikzpicture}
\legende{Schéma de la trame de production : progression logique de l'opinion à la condition nécessaire (只有…才…).}
\end{popfigure}

\begin{savaistu}[Le vocabulaire de l'intégrité en Chine : \cle{廉} (lián) et \cle{贪} (tān)]
Le caractère \zh{廉} (lián, intègre, honnête, bas prix) compose \zh{廉洁} (liánjié, intégrité morale) et \zh{清廉} (qīnglián, honnête et pur). Son antonyme culturel est \zh{贪} (tān, cupide, avide) : \zh{贪污} (tānwū, corruption/détournement), \zh{贪婪} (tānlán, cupidité).
\medskip
Historiquement, l'idéal du lettré-fonctionnaire (士大夫, shìdàfū) confucéen est d'être \zh{两袖清风} (liǎng xiù qīng fēng) — « deux manches pleines de vent pur », c'est-à-dire ne rien emporter en partant de sa charge. Aujourd'hui, la campagne \zh{反腐倡廉} (fǎnfū chánglián — « lutter contre la corruption, promouvoir l'intégrité ») est une priorité politique majeure (\zh{「打老虎、拍苍蝇」} dǎ lǎohǔ, pāi cāngying — frapper les tigres, chasser les mouches). Au Cameroun, l'opération \zh{Épervier} vise un objectif similaire : la tolérance zéro.
\end{savaistu}

\begin{exercice}[Devoir maison : Version finale et affiche]
\begin{enumerate}
    \item Rédigez la \textbf{version finale propre} de votre paragraphe (5 phrases) sur feuille A4 lignée. Écrivez les caractères soigneusement (ordre des traits). Ajoutez le pinyin en petit au-dessus des mots nouveaux seulement. Traduisez en français en regard.
    \item Finalisez votre \textbf{affiche de slogan} (format A3) choisie par votre groupe. Remettez-la pour l'affichage la semaine prochaine.
    \item \textbf{Défi bonus (oral) :} Enregistrez-vous (1 min max) en lisant votre paragraphe + votre slogan. Envoyez le fichier audio sur la plateforme OnBuch+ (rubrique « Oral 1ère A - Module 3 »). Critères : tons justes, débit fluide, pas de hésitation sur \zh{只有…才…}.
\end{enumerate}
\end{exercice}

\begin{corrige}[Exercice - Attendus pour la version finale]
\textbf{Paragraphe type (niveau visé) :}
\zh{我认为腐败是危害社会公平的毒瘤。因为有些官员利用职权谋取私利，所以破坏了法治精神。政府应该建立透明制度，终身追责。我们学生应该抵制不正之风，做诚信青年。只有全民参与监督，清廉中国才能早日实现。}
\\
\textit{Traduction : Je pense que la corruption est une tumeur maligne nuisant à l'équité sociale. Parce que certains fonctionnaires utilisent leur pouvoir pour chercher des profits privés, ils détruisent l'esprit de l'État de droit. Le gouvernement devrait établir un système transparent, poursuivre à vie. Nous étudiants devrions résister aux mauvaises pratiques, être des jeunes intègres. Seulement si tout le peuple participe à la surveillance, une Chine intègre pourra se réaliser rapidement.}
\\
\textbf{Vocabulaire bonus HSK 4+ utilisé :} \zh{毒瘤} (tumeur maligne/métaphore), \zh{谋取私利} (chercher profit personnel), \zh{法治} (État de droit), \zh{透明制度} (système transparent), \zh{终身追责} (responsabilité à vie), \zh{不正之风} (mauvaises pratiques/vents), \zh{诚信} (intégrité/honnêteté), \zh{清廉中国} (Chine intègre - slogan officiel).
\end{corrige}

\newpage

\coursec{Activité d'intégration}

\courssub{Objectifs de l'activité}
Cette activité de synthèse vise à mobiliser l'ensemble des acquis de la leçon 27 dans une tâche de communication authentique et engageante. À l'issue de cette séance, l'élève doit être capable de :
\begin{itemize}
    \item \textbf{Réinvestir le lexique thématique} : employer correctement au moins six mots du glossaire (noms, verbes, adjectifs) liés à la corruption et à l'intégrité.
    \item \textbf{Maîtriser la structure \cle{只有...才...}} : l'utiliser de manière pertinente pour exprimer une condition nécessaire (seule cette condition permet d'atteindre le résultat).
    \item \textbf{Produire un écrit court et percutant} : rédiger un slogan et un court texte argumentatif en caractères chinois simplifiés, accompagnés de pinyin et de traduction française.
    \item \textbf{Présenter oralement} : exposer le travail du groupe en chinois (1 minute) avec une prononciation soignée (tons, liaisons).
    \item \textbf{Travailler en équipe} : négocier le contenu, répartir les rôles (scribe, illustrateur, orateur, correcteur) et respecter le temps imparti.
\end{itemize}

\courssub{Situation}
\begin{textebox}[Contexte : Le Lycée Bilingue de Yaoundé lance une campagne « Intégrité »]
Le proviseur du \textbf{Lycée Bilingue de Yaoundé} (雅温得双语高中 \emph{Yǎwēndé Shuāngyǔ Gāozhōng}) organise une semaine de la citoyenneté. Le thème de cette année est \zh{「拒绝腐败，从我做起」} \emph{(Jùjué fǔbài, cóng wǒ zuòqǐ)} — « Refuser la corruption, ça commence par moi ».

Votre classe de \textbf{1ère A} participe à un concours d'affiches bilingues (chinois / français). Les meilleures affiches seront laminées et affichées dans les couloirs du bâtiment administratif, à la cantine et à la bibliothèque (图书馆 \emph{túshūguǎn}).

\medskip
\textbf{Document déclencheur : Extrait du règlement intérieur affiché au tableau d'affichage (布告栏 \emph{bùgàolán})}
\begin{center}
\begin{tabularx}{\linewidth}{|X|}
\hline
\zh{校规第12条：严禁贿赂、贪污、徇私舞弊。} \\
\emph{Xiàoguī dì shí'èr tiáo: Yánjìn huìlù, tānwū, xùnsīwǔbì.} \\
« Article 12 du règlement : Sont strictement interdits la corruption, la malversation et le favoritisme. » \\
\hline
\zh{发现违规行为，请向纪检部 (纪律检查部) 举报。} \\
\emph{Fāxiàn wéiguī xíngwéi, qǐng xiàng jìjiébù (jìlǜ jiǎnchá bù) jǔbào.} \\
« Tout comportement contrevenant doit être signalé au Bureau de la discipline (纪检部). » \\
\hline
\zh{只有全体师生共同监督，才能建设廉洁校园。} \\
\emph{Zhǐyǒu quántǐ shīshēng gòngtóng jiāndū, cáinéng jiànshè liánjié xiàoyuán.} \\
« Seule la supervision de toute la communauté éducative permet de bâtir un campus intègre. » \\
\hline
\end{tabularx}
\end{center}
\emph{Votre mission : En groupe de 4, concevoir l'affiche la plus convaincante.}

\end{textebox}

\courssub{Consignes}
\begin{enumerate}
    \item \textbf{Formation des groupes (5 min)} : Constituez des groupes de 4 élèves. Désignez : un \textbf{secrétaire} (note les phrases en caractères), un \textbf{correcteur} (vérifie pinyin, tons, ordre des traits), un \textbf{illustrateur} (dessine le visuel), un \textbf{porte-parole} (présente à l'oral).
    \item \textbf{Création du slogan (10 min)} : Rédigez \textbf{un slogan court (4 à 8 caractères)} en chinois. Il doit être percutant, facile à mémoriser et utiliser au moins un mot du glossaire.
        \begin{itemize}
            \item Écrivez le slogan en \textbf{grands caractères} au centre de l'affiche (format A3 ou A2).
            \item Ajoutez le \textbf{pinyin} (avec tons) au-dessus ou en dessous.
            \item Ajoutez la \textbf{traduction française} en petit.
        \end{itemize}
    \item \textbf{Rédaction de l'argumentaire (20 min)} : Rédigez \textbf{3 phrases complètes en chinois} expliquant \emph{pourquoi} il faut lutter contre la corruption dans l'école.
        \begin{itemize}
            \item \textbf{Contrainte lexicale :} Réutilisez \textbf{au moins 6 mots distincts} du glossaire de la leçon (voir \courssub{Ressources à mobiliser}).
            \item \textbf{Contrainte grammaticale :} L'une des 3 phrases \textbf{doit contenir la structure \cle{只有...才...}}.
            \item Écrivez les phrases en caractères, avec pinyin et traduction française.
        \end{itemize}
    \item \textbf{Illustration et mise en page (15 min)} : L'illustrateur réalise un dessin simple (symbole, scène, métaphore visuelle) qui renforce le message. Soignez la lisibilité (couleurs contrastées, police lisible).
    \item \textbf{Répétition orale (5 min)} : Le porte-parole s'entraîne à lire le slogan et les 3 phrases. Vérifiez les tons difficiles (ex. \zh{贿赂 huìlù} 4-4, \zh{廉洁 liánjié} 2-2, \zh{监督 jiāndū} 1-1).
    \item \textbf{Présentation et vote (15 min)} : Chaque groupe présente son affiche (1 minute max). La classe vote à bulletin secret pour : (1) Meilleur slogan, (2) Meilleur argumentaire, (3) Meilleure illustration. Les affiches sont exposées.
\end{enumerate}

\courssub{Ressources à mobiliser}
\medskip
\noindent\textbf{1. Glossaire thématique (Extrait de la leçon 27) — Mots « à réinvestir »}
\begin{center}
\begin{tabularx}{\linewidth}{|X|X|X|X|}
\hline
\rowcolor{popblueL} \textbf{Caractères} & \textbf{Pinyin} & \textbf{Nature} & \textbf{Traduction} \\ \hline
腐败 & fǔbài & n./v. & corruption ; corrompre \\
贿赂 & huìlù & n./v. & pot-de-vin ; corrompre (offrir de l'argent) \\
贪污 & tānwū & v. & détourner des fonds, malverser \\
徇私舞弊 & xùnsīwǔbì & v. (idiome) & favoritisme et fraude / abuser de sa charge pour des intérêts privés \\
举报 & jǔbào & v. & dénoncer, signaler (aux autorités) \\
监督 & jiāndū & v./n. & superviser, contrôler ; supervision \\
廉洁 & liánjié & adj. & intègre, honnête, incorruptible \\
公正 & gōngzhèng & adj. & juste, équitable, impartial \\
透明 & tòumíng & adj. & transparent \\
责任 & zérèn & n. & responsabilité \\
风气 & fēngqì & n. & atmosphère, mœurs, climat (moral) \\
损害 & sǔnhài & v. & nuire à, porter préjudice \\
\hline
\end{tabularx}
\end{center}

\medskip
\noindent\textbf{2. Point de grammaire clé : \cle{只有...才...} (Seulement si... alors...)}
\begin{propriete}{Structure conditionnelle nécessaire : 只有...才...}
\emph{Seule la réalisation de la condition (只有) permet l'aboutissement du résultat (才). L'ordre est fixe.}
\begin{center}
\begin{tabularx}{\linewidth}{|X|X|X|}
\hline
\rowcolor{popblueL} \textbf{Schéma} & \textbf{Exemple (Caractères + Pinyin)} & \textbf{Traduction} \\ \hline
\zh{只有 + Condition + 才 + Résultat} & \zh{只有拒绝贿赂，才能保持廉洁。} \emph{Zhǐyǒu jùjué huìlù, cáinéng bǎochí liánjié.} & Ce n'est qu'en refusant les pots-de-vin qu'on peut rester intègre. \\ \hline
\zh{只有 + Suj. + Verbe + 才 + Verbe} & \zh{只有全体师生监督，风气才会变好。} \emph{Zhǐyǒu quántǐ shīshēng jiāndū, fēngqì cái huì biàn hǎo.} & Ce n'est que si tout le personnel et les élèves surveillent que les mœurs s'amélioreront. \\ \hline
\end{tabularx}
\end{center}
\medskip
\end{propriete}

\begin{attention}{Pièges fréquents pour francophones}
\begin{itemize}
    \item \textbf{Ne pas inverser 只有 et 才 :} \zh{*才只有...} est faux.
    \item \textbf{Ne pas oublier 才 :} En français « Seulement si..., (alors)... », le « alors » est souvent implicite. En chinois, \cle{才} est \textbf{obligatoire} dans la proposition principale pour marquer la restriction.
    \item \textbf{Négation :} On nie généralement le résultat : \zh{如果不举报，就不能制止腐败。} (Si on ne dénonce pas, on ne peut pas arrêter la corruption). \zh{只有} ne s'utilise pas directement avec \zh{不} dans la proposition conditionnelle standard.
    \item \textbf{Confusion avec 只要 :} \zh{只要...就...} = « Aussitôt que / Pourvu que... alors... » (condition suffisante). \zh{只有...才...} = « Seulement si... » (condition \textbf{nécessaire}).
\end{itemize}
\end{attention}

\medskip
\noindent\textbf{3. Rappel : Ordre des traits pour 3 caractères « affiche »}
\begin{itemize}
    \item \zh{\cle{腐} fǔ} (radical \zh{月} « viande/lune » + \zh{府} « bureau/gouvernement ») : 1. 月 (gauche : trait horizontal, trait vertical, 2 horizontaux) → 2. 府 (haut → bas : 广 + 付). \emph{Idée : la chair pourrit dans les bureaux.}
    \item \zh{\cle{廉} lián} (radical \zh{广} « toit/abri » + \zh{兼} « simultanément ») : 1. 广 (point, horizontal, vertical crochet) → 2. 兼 (haut → bas : 八 + 丷 + 一 + 八). \emph{Idée : sous le toit, on partage équitablement.}
    \item \zh{\cle{举} jǔ} (radical \zh{廾} « deux mains ») : 1. 兴 (haut : 八 + 同 sans le trait final) → 2. 廾 (bas : deux mains levées). \emph{Idée : lever les deux mains pour signaler/dénoncer.}
\end{itemize}

\courssub{Proposition de production et corrigé commenté}
\begin{exoresolu}{Affiche modèle — Groupe « Intégrité » (Intégrité 组)}
\textbf{Production élève (à reproduire sur l'affiche A3) :}

\medskip
\noindent\textbf{SLOGAN (Grand format, centre de l'affiche) :}
\begin{center}
\zh{\LARGE 拒绝贿赂，人人有责！} \\
\emph{Jùjué huìlù, rénrén yǒu zé!} \\
« Refuser la corruption, c'est l'affaire de tous ! »
\end{center}

\medskip
\noindent\textbf{ARGUMENTAIRE (3 phrases, alignées à gauche) :}

1. \zh{腐败损害学校公正，破坏良好风气。} \\
\emph{Fǔbài sǔnhài xuéxiào gōngzhèng, pòhuài liánghǎo fēngqì.} \\
« La corruption nuit à l'équité de l'école et détruit un bon climat moral. »

2. \zh{只有全体师生敢于举报，贪污现象才会减少。} \\
\emph{Zhǐyǒu quántǐ shīshēng gǎnyú jǔbào, tānwū xiànxiàng cái huì jiǎnshǎo.} \\
« Ce n'est que si tout le personnel et les élèves osent dénoncer que les phénomènes de malversation diminueront. »

3. \zh{建设廉洁校园，需要你我共同监督。} \\
\emph{Jiànshè liánjié xiàoyuán, xūyào nǐ wǒ gòngtóng jiāndū.} \\
« Construire un campus intègre nécessite une supervision commune de ta part et de la mienne. »

\medskip
\noindent\textbf{ILLUSTRATION (Description pour l'illustrateur) :}
Un grand arbre (校园树 \emph{xiàoyuán shù}) dont les racines sont \zh{廉洁} (intégrité), \zh{公正} (équité), \zh{透明} (transparence). Un élève arrose l'arbre avec un arrosoir marqué \zh{监督} (supervision). Un ciel bleu avec le soleil \zh{阳光} (soleil/lumière). Au pied de l'arbre, un serpent \zh{腐败} est coupé en deux par une hache \zh{举报}.

\tcblower
\textbf{Commentaire pédagogique du professeur (Corrigé commenté) :}

\begin{itemize}
    \item \textbf{Choix lexical (6+ mots validés) :} \zh{贿赂 huìlù} (slogan), \zh{腐败 fǔbài}, \zh{公正 gōngzhèng}, \zh{风气 fēngqì}, \zh{举报 jǔbào}, \zh{贪污 tānwū}, \zh{廉洁 liánjié}, \zh{监督 jiāndū}. \textbf{Objectif atteint (8 mots).}
    \item \textbf{Structure \cle{只有...才...} (Phrase 2) :} \zh{只有全体师生敢于举报 (Condition) ，贪污现象才会减少 (Résultat)}. La structure est respectée. L'adverbe \zh{敢于 gǎnyú} (oser) enrichit le sujet. \zh{才会} marque le futur hypothétique nécessaire.
    \item \textbf{Grammaire / Syntaxe :}
        \begin{itemize}
            \item Phrase 1 : Structure \zh{Sujet + Verbe + Objet 1, Verbe + Objet 2} (série verbale). \zh{损害} et \zh{破坏} sont bien des verbes transitifs.
            \item Phrase 3 : Structure \zh{Sujet (thème : 建设廉洁校园) + 需要 + Objet (你我共同监督)}. L'usage de \zh{你我 nǐ wǒ} (toi et moi / nous) est idiomatique pour « tout le monde / chacun ».
        \end{itemize}
    \item \textbf{Pinyin \& Tons :} Attention aux tons 4-4 sur \zh{huìlù}, 2-2 sur \zh{liánjié}, 1-1 sur \zh{jiāndū}. Le ton neutre sur \zh{xiànxiàng} (phénomène) n'existe pas, c'est 4-4. Le \zh{de} de \zh{良好} n'est pas nécessaire devant \zh{风气} (adjectif monosyllabique ou bisyllabique qualificatif direct souvent sans \zh{de} dans l'écrit formel/slogan).
    \item \textbf{Registre :} Langage formel, style affichage / slogan (四字成语 style pour le slogan : \zh{拒绝贿赂，人人有责} — 4+4+4 caractères), phrases complètes pour l'argumentaire. Adapté au contexte scolaire camerounais (référence implicite aux valeurs civiques du MINESEC).
\end{itemize}
\end{exoresolu}

\courssub{Bilan de l'activité}
\begin{aretenir}{Ce que vous avez démontré}
\begin{itemize}
    \item \textbf{Compétence lexicale :} Vous savez piocher dans un glossaire thématique précis (\zh{腐败, 贿赂, 贪污, 举报, 监督, 廉洁...}) pour construire un discours cohérent.
    \item \textbf{Compétence grammaticale :} Vous maîtrisez l'articulation logique \cle{只有...才...} pour exprimer une condition \textbf{nécessaire et exclusive}, structure clé de l'argumentation en chinois (HSK 3/4).
    \item \textbf{Compétence scripturale :} Vous écrivez des caractères complexes (\zh{腐, 廉, 举, 贿, 污}) en respectant l'ordre des traits et les radicaux sémantiques (月, 广, 廾).
    \item \textbf{Compétence communicative :} Vous avez produit un texte \textbf{authentique} (affiche publique), \textbf{bilingue} (chinois/français), \textbf{persuasif} (slogan + arguments) et \textbf{contextualisé} (Lycée de Yaoundé, règlement intérieur).
    \item \textbf{Compétence interculturelle :} Vous reliez la valeur universelle d'intégrité (\zh{廉洁}) au contexte camerounais (lutte contre la corruption, CONAC, opération \emph{Épervier}) et chinois (campagne \zh{反腐倡廉 fǎnfū chánglián} « Lutte contre la corruption et promotion de l'intégrité »).
\end{itemize}

\end{aretenir}

\begin{methode}{Pour aller plus loin (Prolongements possibles)}
\begin{enumerate}
    \item \textbf{Expression orale approfondie :} Organisez un débat en chinois : \zh{「学生应该举报老师的腐败行为吗？」} (Les élèves doivent-ils dénoncer la corruption des professeurs ?). Utilisez \zh{虽然...但是...}, \zh{不仅...而且...}.
    \item \textbf{Traduction Thème/Version :} Traduisez en chinois : « Au Cameroun, la Commission Nationale Anti-Corruption (CONAC) encourage les citoyens à dénoncer les actes de corruption. Seule la transparence peut garantir la justice sociale. » (Mots-clés : \zh{国家反腐败委员会}, \zh{鼓励}, \zh{公民}, \zh{透明度}, \zh{社会正义}).
    \item \textbf{Projet numérique :} Créez une version numérique de l'affiche (Canva, PowerPoint) avec enregistrement audio du slogan par le porte-parole. Publiez sur la plateforme de l'établissement (OnBuch+, Pronote).
\end{enumerate}
\end{methode}

\begin{savaistu}{Regard croisé Cameroun — Chine : Lutte contre la corruption}
\begin{itemize}
    \item \textbf{Cameroun :} La \textbf{CONAC} (Commission Nationale Anti-Corruption), créée en 2006, pilote la stratégie nationale. L'opération \textbf{Épervier} (depuis 2006) vise la répression des grands dossiers. Le numéro vert \textbf{1517} permet les signalements anonymes. Le Code pénal (art. 134 et suivants) sanctionne la corruption passive et active.
    \item \textbf{Chine :} La campagne \zh{反腐倡廉 fǎnfū chánglián} (litt. « S'opposer à la corruption, promouvoir l'intégrité ») est une priorité politique majeure depuis 2012. Le \zh{中央纪委国家监委} (Commission centrale de discipline / Commission nationale de supervision) enquête sur les cadres du Parti et de l'État. Les peines vont jusqu'à la peine de mort avec sursis pour les montants très élevés (\zh{巨额贪污}).
    \item \textbf{Point commun :} Dans les deux pays, l'école est vue comme le lieu premier de l'éducation à l'intégrité (\zh{廉洁教育 liánjié jiàoyù}). Les affiches \zh{拒绝腐败，从我做起} sont omniprésentes dans les établissements chinois (écoles primaires, universités, administrations).
\end{itemize}
\end{savaistu}

\newpage

\coursec{Méthodes et exercices résolus}

\coursec{Leçon 27 — Vocabulaire et compréhension de texte : corruption}

\courssub{Mise en route et découverte du thème}

\begin{experience}[Entrée en matière : Corruption et citoyenneté]
\begin{enumerate}
\item \textbf{Observation d'une affiche.} Regardez l'affiche ci-dessous (imaginée pour une campagne au Cameroun) et décrivez-la en français : lieu, personnages, message visuel, slogan.
\item \textbf{Brainstorming lexical.} En binôme, listez en 2 minutes tous les mots français qui vous viennent à l'esprit sur le thème de la \cle{corruption} (noms, verbes, adjectifs). Mettez en commun au tableau ; le professeur écrit les équivalents chinois au fur et à mesure.
\item \textbf{Question déclencheur.} \zh{你认为腐败对国家有什么坏处？} \emph{(Nǐ rènwéi fǔbài duì guójiā yǒu shénme huàichù ?)} « Selon toi, quels sont les méfaits de la corruption pour le pays ? » Réponse courte attendue : \zh{腐败让国家变穷，人民不信任政府。} \emph{(Fǔbài ràng guójiā biàn qióng, rénmín bù xìnrèn zhèngfǔ.)}
\end{enumerate}
\end{experience}

\begin{popfigure}
\begin{tikzpicture}[mindmap, grow cyclic, every node/.style=concept, concept color=popblue!60,
    level 1/.style={level distance=4.5cm,sibling angle=90},
    level 2/.style={level distance=3cm,sibling angle=45}]
\node {腐败 \emph{fǔbài}}
    child[concept color=poporange!70] {node {贪污 \emph{tānwū}}}
    child[concept color=popgreen!70] {node {受贿 \emph{shòuhuì}}}
    child[concept color=poppurple!70] {node {行贿 \emph{xínghuì}}}
    child[concept color=poppink!70] {node {贿赂 \emph{huìlù}}};
\end{tikzpicture}
\legende{Carte mentale : champ lexical de la corruption (niveau HSK 3-4)}
\end{popfigure}

\courssub{Texte support : lecture et compréhension globale}

\begin{textebox}[Texte support : « Un fonctionnaire honnête »]
王叔叔是一个公务员。有一天，一个商人想给他很多钱，请他帮忙做一件不对的事情。王叔叔说：« 不行！这样做是腐败，是犯法的。» 他把这件事告诉了警察。后来，那个商人受到了法律的惩罚。大家都说王叔叔是一个诚实的人。 \\
Wáng shūshu shì yī gè gōngwùyuán. Yǒu yī tiān, yī gè shāngrén xiǎng gěi tā hěn duō qián, qǐng tā bāngmáng zuò yī jiàn bù duì de shìqing. Wáng shūshu shuō: « Bù xíng! Zhèyàng zuò shì fǔbài, shì fànfǎ de. » Tā bǎ zhè jiàn shì gàosu le jǐngchá. Hòulái, nà gè shāngrén shòudào le fǎlǜ de chéngfá. Dàjiā dōu shuō Wáng shūshu shì yī gè chéngshí de rén. \\
« L'oncle Wang est un fonctionnaire. Un jour, un homme d'affaires voulut lui donner beaucoup d'argent et lui demanda de l'aider à faire quelque chose de mal. L'oncle Wang dit : "Non ! Agir ainsi, c'est de la corruption, c'est illégal." Il raconta cette affaire à la police. Plus tard, cet homme d'affaires subit la punition de la loi. Tout le monde dit que l'oncle Wang est un homme honnête. »
\end{textebox}

\begin{methode}[Comprendre un texte en trois lectures]
Pour comprendre un texte chinois sur un thème de société comme la corruption, procède toujours en trois lectures méthodiques :
\begin{enumerate}
\item \textbf{Lecture globale (silencieuse)} : lis le texte entier sans t'arrêter sur les mots inconnus. Repère le \cle{thème} (ici : \zh{腐败} \emph{fǔbài}, « la corruption »), le \cle{lieu} (implicite : lieu de travail), les \cle{personnages} (\zh{王叔叔} \emph{Wáng shūshu}, \zh{商人} \emph{shāngrén}, \zh{警察} \emph{jǐngchá}) et la \cle{conclusion} (dernière phrase). Le titre et la dernière phrase donnent souvent l'idée principale.
\item \textbf{Lecture analytique} : souligne les mots du champ lexical du thème (\zh{腐败， 贪污， 受贿， 法律， 惩罚， 诚实}…). Pour chaque mot inconnu, devine son sens grâce au \cle{contexte} et aux \cle{clés (radicaux)} : la clé \zh{讠} (parole) indique souvent un verbe de communication (\zh{告诉} \emph{gàosu}), la clé \zh{心} (cœur) un sentiment ou qualité morale (\zh{诚实} \emph{chéngshí}), la clé \zh{钅/贝} (métal/coquillage $\rightarrow$ argent/biens) une notion liée à l'argent (\zh{贪污} \emph{tānwū}, \zh{受贿} \emph{shòuhuì}, \zh{贿赂} \emph{huìlù}).
\item \textbf{Lecture de vérification} : relis en traduisant mentalement phrase par phrase, en respectant l'ordre chinois \cle{Sujet + Temps/Lieu + Verbe + Complément}. Vérifie que ta compréhension est cohérente avec les questions posées.
\end{enumerate}
\textbf{Réflexe d'examen :} avant de répondre, \textbf{localise dans le texte la phrase qui contient la réponse} et recopie-en les mots-clés dans ta réponse en phrase complète.
\end{methode}

\courssub{Glossaire du thème, prononciation, clés et ordre des traits}

\begin{aretenir}[Tableau de vocabulaire : La corruption et la loi]
\begin{center}
\begin{tabularx}{\linewidth}{|X|X|X|X|}
\hline
\rowcolor{popblueL} \textbf{Caractères} & \textbf{Pinyin} & \textbf{Nature} & \textbf{Traduction} \\ \hline
腐败 & fǔbài & nom / verbe & corruption ; corrompre \\ \hline
贪污 & tānwū & verbe & détourner des fonds publics, piller \\ \hline
受贿 & shòuhuì & verbe & accepter des pots-de-vin \\ \hline
行贿 & xínghuì & verbe & offrir des pots-de-vin \\ \hline
贿赂 & huìlù & nom & pot-de-vin, corruption (acte) \\ \hline
公务员 & gōngwùyuán & nom & fonctionnaire \\ \hline
商人 & shāngrén & nom & commerçant, homme d'affaires \\ \hline
警察 & jǐngchá & nom & policier \\ \hline
法律 & fǎlǜ & nom & loi \\ \hline
犯法 & fànfǎ & verbe & violer la loi, commettre un délit \\ \hline
惩罚 & chéngfá & nom / verbe & punition ; punir \\ \hline
受到 & shòudào & verbe & recevoir (subir : punition, éducation, accueil) \\ \hline
诚实 & chéngshí & adj. & honnête, sincère \\ \hline
公平 & gōngpíng & adj. & juste, équitable \\ \hline
反对 & fǎnduì & verbe & s'opposer à, combattre \\ \hline
遵守 & zūnshǒu & verbe & respecter (loi, règle, discipline) \\ \hline
告诉 & gàosu & verbe & dire à, informer, raconter \\ \hline
行为 & xíngwéi & nom & comportement, acte, agissement \\ \hline
影响 & yǐngxiǎng & nom / verbe & influence ; influencer \\ \hline
相信 & xiāngxìn & verbe & croire, faire confiance à \\ \hline
政府 & zhèngfǔ & nom & gouvernement \\ \hline
\end{tabularx}
\end{center}
\end{aretenir}

\begin{propriete}[Prononciation : tons et pièges sonores]
\begin{itemize}
\item \zh{腐败} \emph{fǔbài} : 3\textsuperscript{e} ton (descendant-montant) + 4\textsuperscript{e} ton (descendant). Ne pas confondre avec \emph{fùbái}.
\item \zh{诚实} \emph{chéngshí} : deux 2\textsuperscript{e} tons (montants). Attention : \emph{shí} (réel/honnête) $\neq$ \emph{shì} (être/affaire).
\item \zh{贪污} \emph{tānwū} : 1\textsuperscript{er} ton (haut) + 1\textsuperscript{er} ton. \zh{污} \emph{wū} (sale) $\neq$ \zh{乌} \emph{wū} (noir/corbeau).
\item \zh{受贿} \emph{shòuhuì} : 4\textsuperscript{e} ton + 4\textsuperscript{e} ton. \zh{受} \emph{shòu} (recevoir) $\neq$ \zh{授} \emph{shòu} (enseigner/transmettre).
\item \zh{法律} \emph{fǎlǜ} : 3\textsuperscript{e} ton + 4\textsuperscript{e} ton. \zh{律} \emph{lǜ} (règle/loi) contient \zh{聿} \emph{yù} (pinceau), pas \zh{刷} \emph{shuā} (brosse).
\end{itemize}
\textbf{Exercice oral :} Le professeur dicte les mots dans le désordre ; l'élève écrit le pinyin avec tons puis les caractères.
\end{propriete}

\begin{definition}[Clés (radicaux) et ordre des traits : 6 caractères clés]
\begin{center}
\begin{tabularx}{\linewidth}{|c|X|X|X|}
\hline
\rowcolor{popblueL} \textbf{Caractère} & \textbf{Clé (Radical)} & \textbf{Sens de la clé} & \textbf{Ordre des traits (résumé)} \\ \hline
\zh{腐} \emph{fǔ} & \zh{氵} (trois points d'eau) & Liquide, pourriture & Haut $\rightarrow$ bas : \zh{氵} puis \zh{府} (haut $\rightarrow$ bas, gauche $\rightarrow$ droite) \\ \hline
\zh{败} \emph{bài} & \zh{攴} (pū, frapper) & Action de frapper, briser & Gauche \zh{啇} (haut $\rightarrow$ bas) puis droite \zh{攴} (horizontal, point, crochet) \\ \hline
\zh{贪} \emph{tān} & \zh{贝} (bèi, coquillage) & Argent, biens, valeur & Haut \zh{今} (horizontal, vertical, horizontal) puis bas \zh{贝} \\ \hline
\zh{诚} \emph{chéng} & \zh{讠} (yán, parole) & Parole, langage & Gauche \zh{讠} (point, courbe) puis droite \zh{成} (haut $\rightarrow$ bas) \\ \hline
\zh{法} \emph{fǎ} & \zh{氵} (eau) & Norme, standard (l'eau qui va droit) & \zh{氵} puis \zh{去} (haut $\rightarrow$ bas : \zh{土} puis \zh{厶}) \\ \hline
\zh{律} \emph{lǜ} & \zh{彳} (chì, pas) & Marche, règle, ordre & Gauche \zh{彳} (trois traits) puis droite \zh{聿} (point, horizontaux, vertical, crochet) \\ \hline
\end{tabularx}
\end{center}
\textbf{Rappel règles d'ordre :} 1. Horizontal avant vertical (\zh{十}) 2. Gauche avant droite (\zh{人}) 3. Haut avant bas (\zh{三}) 4. Extérieur avant intérieur (\zh{口}) 5. Fermer en dernier (\zh{田}) 6. Trait central avant ailes (\zh{小}).
\end{definition}

\courssub{Questions de compréhension et approfondissement}

\begin{exoresolu}[Compréhension détaillée du texte support]
\textbf{Consigne.} Réponds en chinois (caractères + pinyin) par des phrases complètes en te basant \emph{uniquement} sur le texte.
\begin{enumerate}
\item 王叔叔是做什么工作的？
\item 商人想做什么？
\item 王叔叔同意了吗？为什么？
\item 后来商人怎么样了？
\item 为什么大家都说王叔叔是诚实的人？
\end{enumerate}
\tcblower
\textbf{Solution détaillée.}
\begin{enumerate}
\item \textbf{王叔叔是一个公务员。} \emph{Wáng shūshu shì yī gè gōngwùyuán.} (Phrase 1 : \zh{是} + classif. \zh{个} + nom).
\item \textbf{商人想给他很多钱，请他帮忙做一件不对的事情。} \emph{Shāngrén xiǎng gěi tā hěn duō qián, qǐng tā bāngmáng zuò yī jiàn bù duì de shìqing.} (Phrase 2 : \zh{想} + V1, \zh{请} + V2). \zh{一件} classif. pour \zh{事情}.
\item \textbf{王叔叔没有同意，因为这样做是腐败，是犯法的。} \emph{Wáng shūshu méiyǒu tóngyì, yīnwèi zhèyàng zuò shì fǔbài, shì fànfǎ de.}
\zh{没有} + V (négation passé accompli) ; \zh{因为}… introduit la cause ; \zh{的} final marque l'affirmation/insistance.
\item \textbf{后来，那个商人受到了法律的惩罚。} \emph{Hòulái, nà gè shāngrén shòudào le fǎlǜ de chéngfá.} Structure \zh{受到} + \zh{惩罚} (verbe + objet) ; \zh{了} marque l'accompli.
\item \textbf{因为他不收贿赂，把坏人的事告诉警察，所以大家说他诚实。} \emph{Yīnwèi tā bù shōu huìlù, bǎ huàirén de shì gàosu jǐngchá, suǒyǐ dàjiā shuō tā chéngshí.} Inférence logique : \zh{不收贿赂} (n'accepte pas de pots-de-vin) = \zh{诚实}.
\end{enumerate}
\end{exoresolu}

\begin{exoresolu}[Vocabulaire : complétion et collocations]
\textbf{Énoncé.} Complète les phrases avec les mots du tableau (forme correcte) : \zh{腐败、诚实、法律、惩罚、反对、遵守、受贿、贪污、受到、告诉}。
\begin{enumerate}
\item 王叔叔是一个 \_\_\_\_ 的人，他从来不\_\_\_\_别人的钱。
\item \_\_\_\_ 和 \_\_\_\_ 都是 \_\_\_\_ 的行为。
\item 犯法的人应该 \_\_\_\_ \_\_\_\_ 。
\item 每个公民都应该 \_\_\_\_ \_\_\_\_ 。
\item 我们喀麦隆学生 \_\_\_\_ \_\_\_\_ 。
\item 王叔叔把这件事 \_\_\_\_ 了警察。
\end{enumerate}
\tcblower
\textbf{Solution détaillée.}
\begin{enumerate}
\item \zh{诚实} (chéngshí) ; \zh{收} (shōu) — \zh{收钱} « accepter de l'argent » (contexte corruption). \textit{王叔叔是一个诚实的人，他从来不收别人的钱。}
\item \zh{贪污} (tānwū) ; \zh{受贿} (shòuhuì) ; \zh{腐败} (fǔbài). \textit{贪污和受贿都是腐败的行为。} (\zh{腐败的} attribut).
\item \zh{受到} (shòudào) ; \zh{惩罚} (chéngfá). Collocation figée \zh{受到惩罚}. \textit{犯法的人应该受到惩罚。}
\item \zh{遵守} (zūnshǒu) ; \zh{法律} (fǎlǜ). Collocation figée \zh{遵守法律}. \textit{每个公民都应该遵守法律。}
\item \zh{反对} (fǎnduì) ; \zh{腐败} (fǔbài). Verbe transitif direct. \textit{我们喀麦隆学生反对腐败。}
\item \zh{告诉} (gàosu). Structure \zh{把} + objet + \zh{告诉} + destinataire + \zh{了}. \textit{王叔叔把这件事告诉了警察。}
\end{enumerate}
\textbf{Point méthode :} Au Bac, ces questions testent les \cle{collocations} (\zh{受到惩罚}， \zh{遵守法律}， \zh{反对腐败}， \zh{收贿赂/受贿}). Apprends les verbes \textbf{avec} leurs compléments habituels.
\end{exoresolu}

\courssub{Production guidée et réinvestissement}

\begin{methode}[Construire une phrase correcte avec le nouveau vocabulaire]
\begin{enumerate}
\item \textbf{Squelette} : Sujet + (Temps) + (Lieu) + Verbe + Complément. \textbf{Jamais d'inversion.}
\item \textbf{Compléments Temps/Lieu} : Placés \textbf{avant} le verbe (ou avant le sujet). Ex: \zh{在喀麦隆，腐败是大问题。} / \zh{腐败在喀麦隆是大问题。}
\item \textbf{Mots-outils indispensables} :
\begin{itemize}
\item \zh{了} : action accomplie / changement d'état (\zh{受到了惩罚了}).
\item \zh{的} : attribut (adj + \zh{的} + nom : \zh{诚实的人}) / possession (\zh{法律的惩罚}).
\item \zh{得} : complément de manière (V + \zh{得} + adj : \zh{做得好}).
\item \zh{过} : expérience passée (\zh{去过中国}).
\end{itemize}
\item \textbf{Classificateurs} : Obligatoires après nombre/démonstratif. \zh{一个公务员}， \zh{一件事}， \zh{一个商人}， \zh{这件事}.
\item \textbf{Négation} : \zh{不} (habitude, futur, volonté) vs \zh{没有} (passé accompli : \zh{没有同意}).
\end{enumerate}
\end{methode}

\begin{exoresolu}[Expression écrite : Traduction thème (Français $\rightarrow$ Chinois)]
\textbf{Énoncé.} Traduis en chinois (caractères + pinyin avec tons).
\begin{enumerate}
\item La corruption est un grand problème dans beaucoup de pays.
\item Ce fonctionnaire a accepté de l'argent (pot-de-vin), donc il a été puni.
\item À l'école, nos professeurs nous apprennent à être honnêtes.
\item Nous devons signaler les actes de corruption à la police.
\end{enumerate}
\tcblower
\textbf{Solution détaillée.}
\begin{enumerate}
\item \zh{腐败是很多国家的一个大问题。} \emph{Fǔbài shì hěn duō guójiā de yī gè dà wèntí.}
Structure : Sujet \zh{腐败} + \zh{是} + [Modificateur \zh{很多国家的}] + Tête \zh{一个大问题}. \zh{的} relie le modificateur.
\item \zh{这个公务员受贿了，所以受到了惩罚。} \emph{Zhè gè gōngwùyuán shòuhuì le, suǒyǐ shòudào le chéngfá.}
\zh{受贿} verbe composé (verbe+objet) $\rightarrow$ \zh{了} après le verbe. \zh{所以} en tête de 2\textsuperscript{e} proposition. \zh{受到惩罚} passif/Subi.
\item \zh{在学校里，我们的老师教我们做诚实的人。} \emph{Zài xuéxiào lǐ, wǒmen de lǎoshī jiāo wǒmen zuò chéngshí de rén.}
Complément lieu \zh{在学校里} en tête. \zh{教} qqn \zh{做} qqch. \zh{诚实的} attribut.
\item \zh{我们应该向警察举报腐败行为。} \emph{Wǒmen yīnggāi xiàng jǐngchá jǔbào fǔbài xíngwéi.}
\zh{向} + destinataire + \zh{举报} (signaler/dénoncer, vocab. HSK 4 utile ici) + objet. \zh{应该} (devoir moral).
\end{enumerate}
\end{exoresolu}

\begin{methode}[Rédiger une réponse d'approfondissement / opinion]
\begin{enumerate}
\item \textbf{Reprends le sujet de la question} : \zh{你认为……} $\rightarrow$ \zh{我认为……} / \zh{我觉得……}.
\item \textbf{Structure argumentaire} : \zh{我认为} + Opinion + \zh{因为} / \zh{所以} + Arguments + \zh{所以/因此} + Proposition d'action (\zh{应该} / \zh{要} / \zh{必须}) + \zh{这样} + Conséquence positive (\zh{越来越好}).
\item \textbf{Formules utiles} :
\begin{itemize}
\item \zh{A 对 B 有坏影响} (A a une mauvaise influence sur B).
\item \zh{让/使} + [qqn] + [verbe/adj] « faire que » (\zh{让国家变穷}).
\item \zh{越来越} + adj « de plus en plus » (\zh{越来越好}).
\end{itemize}
\end{enumerate}
\end{methode}

\begin{exoresolu}[Expression orale/écrite : Débat citoyen]
\textbf{Énoncé.} Réponds en chinois (5-6 phrases) : \zh{你认为腐败对社会有什么影响？我们年轻人应该怎么做？}
\tcblower
\textbf{Modèle de réponse (niveau HSK 3/4 visé).}

我认为腐败对社会有很坏的影响。第一，腐败让国家变穷，老百姓的生活变难。第二，腐败让人民不相信政府，社会不公平。所以，我们年轻人应该从自己做起，做诚实的人，不收贿赂，不办坏事。我们还要监督政府，发现腐败就向警察举报。这样，我们的国家才会越来越好，越来越强。

\emph{Wǒ rènwéi fǔbài duì shèhuì yǒu hěn huài de yǐngxiǎng. Dì yī, fǔbài ràng guójiā biàn qióng, lǎobǎixìng de shēnghuó biàn nán. Dì èr, fǔbài ràng rénmín bù xiāngxìn zhèngfǔ, shèhuì bù gōngpíng. Suǒyǐ, wǒmen niánqīngrén yīnggāi cóng zìjǐ zuòqǐ, zuò chéngshí de rén, bù shōu huìlù, bù bàn huàishì. Wǒmen hái yào jiāndū zhèngfǔ, fāxiàn fǔbài jiù xiàng jǐngchá jǔbào. Zhèyàng, wǒmen de guójiā cái huì yuèláiyuè hǎo, yuèláiyuè qiáng.}

« Je pense que la corruption a une très mauvaise influence sur la société. Premièrement, la corruption appauvrit le pays et rend la vie dure au peuple. Deuxièmement, elle fait que le peuple ne fait plus confiance au gouvernement, la société n'est pas juste. Donc, nous les jeunes, nous devons commencer par nous-mêmes, être des gens honnêtes, ne pas accepter de pots-de-vin, ne pas faire de mauvaises actions. Nous devons aussi surveiller le gouvernement, dès qu'on découvre de la corruption, la signaler à la police. Ainsi, notre pays deviendra de mieux en mieux, de plus en plus fort. »

\textbf{Justifications grammaticales} :
\begin{itemize}
\item \zh{第一… 第二…} (Premièrement… Deuxièmement…) : structuration du discours.
\item \zh{让/使} + objet + complément de résultat/état : \zh{让国家变穷}, \zh{让人民不相信}.
\item \zh{从自己做起} (commencer par soi-même) : expression idiomatique.
\item \zh{监督} (jiāndū, surveiller/contrôler) + \zh{发现} (découvrir) + \zh{就} (alors/aussitôt) + \zh{举报} (jǔbào, dénoncer).
\item \zh{才} (cái) : « seulement ainsi / ce n'est qu'alors que » (condition nécessaire), plus fort que \zh{就}.
\item \zh{越来越} + adj : progression.
\end{itemize}
\end{exoresolu}

\begin{attention}[Les 5 erreurs qui coûtent des points au Bac (Première A)]
\begin{enumerate}
\item \textbf{Tons négligés ou faux en pinyin} : Écrire \emph{fubai} sans tons ou \emph{fùbái} est éliminatoire. \zh{腐败} = \emph{fǔbài} (3+4). \zh{诚实} = \emph{chéngshí} (2+2), pas \emph{chéngshì}. \zh{法律} = \emph{fǎlǜ} (3+4). \textbf{Règle :} Note le ton sur la voyelle principale (a, o, e, i, u, ü), jamais sur la consonne.
\item \textbf{Caractères confondus (clés différentes)} :
\begin{itemize}
\item \zh{受} \emph{shòu} (clé \zh{又} « main droite », recevoir) $\neq$ \zh{授} \emph{shòu} (clé \zh{扌} « main », enseigner).
\item \zh{诚} \emph{chéng} (clé \zh{讠} « parole », sincère) $\neq$ \zh{城} \emph{chéng} (clé \zh{土} « terre », ville).
\item \zh{法} \emph{fǎ} (clé \zh{氵} « eau », loi) $\neq$ \zh{去} \emph{qù} (aller) / \zh{灭} \emph{miè} (éteindre).
\item \zh{贿} \emph{huì} (clé \zh{贝} « argent », pot-de-vin) $\neq$ \zh{惠} \emph{huì} (clé \zh{心} « cœur », bienveillance).
\end{itemize}
\item \textbf{Ordre des mots calqué sur le français} :
\begin{itemize}
\item \textbf{Faux :} \zh{腐败是问题大在很多国家} / \zh{老师教我们在学校里}.
\item \textbf{Vrai :} \zh{腐败是很多国家的一个大问题。} \zh{在学校里，老师教我们。} (Lieu/Temps $\rightarrow$ Sujet $\rightarrow$ Verbe).
\end{itemize}
\item \textbf{Mots-outils oubliés ou mal placés} :
\begin{itemize}
\item Passé accompli : \zh{他受到了惩罚\cle{了}} (pas \zh{他受到惩罚}).
\item Négation passé : \zh{王叔叔\cle{没有}同意} (pas \zh{不同意} qui = « refuse maintenant »).
\item Attribut adjectival : \zh{诚实\cle{的}人}， \zh{腐败\cle{的}行为} (pas \zh{诚实人}).
\item Structure \zh{把} : \zh{把} + objet + \zh{告诉} + qqn + \zh{了} (objet déplacé avant verbe).
\end{itemize}
\item \textbf{Classificateur omis ou erroné} :
\begin{itemize}
\item \zh{一\cle{个}公务员}， \zh{一\cle{个}商人}， \zh{一\cle{件}事}， \zh{一\cle{条}法律} (loi = texte long).
\item Jamais \zh{一公务员}， \zh{一商人}， \zh{一事}.
\end{itemize}
\end{enumerate}
\end{attention}

\begin{savaistu}[Regard croisé Cameroun – Chine : Lutte contre la corruption]
Au Cameroun, la \cle{CONAC} (Commission Nationale Anticorruption) créée en 2006 mène des campagnes « \textit{Corruption : Non !} » dans les administrations et écoles. En Chine, la lutte anticorruption (\zh{反腐倡廉} \emph{fǎnfū chánglián} « combattre la corruption, promouvoir l'intégrité ») est une priorité politique majeure depuis 2012 (\zh{打虎拍蝇} \emph{dǎhǔ pāiyíng} « frapper les tigres et écraser les mouches » : hauts cadres et petits fonctionnaires). Les deux pays partagent le même objectif : \zh{权力关进制度的笼子里} « enfermer le pouvoir dans la cage des institutions » (Xi Jinping). En classe, on peut comparer les slogans : Cameroun « \textit{La corruption tue le développement} » vs Chine \zh{腐败是党执政最大的威胁} « La corruption est la plus grande menace pour le pouvoir du Parti ».
\end{savaistu}

\begin{exercice}[Devoirs / Entraînement personnel]
\begin{enumerate}
\item \textbf{Écriture :} Écris 5 fois chaque caractère du tableau « Clés et ordre des traits » en respectant l'ordre des traits. Récite le pinyin à voix haute.
\item \textbf{Lexique :} Apprends par cœur le tableau de vocabulaire (chinois $\rightarrow$ français et français $\rightarrow$ chinois). Teste-toi sur les collocations : \zh{遵守法律}， \zh{受到惩罚}， \zh{反对腐败}， \zh{受贿/行贿}， \zh{举报腐败}.
\item \textbf{Rédaction :} Écris un court paragraphe (60-80 caractères) sur le thème : \zh{《拒绝腐败，从我做起》} (« Refuser la corruption, commencer par moi »). Utilise au moins 6 mots du vocabulaire de la leçon. Soigne la ponctuation chinoise (， 。).
\item \textbf{Oral :} Prépare une présentation orale de 1 minute : « Pourquoi la corruption est-elle dangereuse pour le Cameroun ? » (en chinois, notes autorisées sur mots-clés).
\end{enumerate}
\end{exercice}

\begin{corrige}[Exercice n° Devoirs - Éléments de correction]
\begin{enumerate}
\item \textbf{Écriture :} Vérification par le professeur (ordre traits, proportions, clés).
\item \textbf{Collocations types :}
\zh{遵守法律} (zūnshǒu fǎlǜ), \zh{受到惩罚} (shòudào chéngfá), \zh{反对腐败} (fǎnduì fǔbài), \zh{受贿/行贿} (shòuhuì/xínghuì), \zh{举报腐败} (jǔbào fǔbài), \zh{诚实守信} (chéngshí shǒuxìn - honnête et digne de confiance).
\item \textbf{Modèle de rédaction :}
\zh{腐败害人害己。作为学生，我们要诚实守信，不作弊，不收贿赂。发现腐败要举报。遵守法律，反对腐败，让社会更公平。} \emph{(Fǔbài hàirén hàijǐ. Zuòwéi xuésheng, wǒmen yào chéngshí shǒuxìn, bù zuòbì, bù shōu huìlù. Fāxiàn fǔbài yào jǔbào. Zūnshǒu fǎlǜ, fǎnduì fǔbài, ràng shèhuì gèng gōngpíng.)}
\item \textbf{Oral :} Évaluation sur : prononciation (tons), fluidité, vocabulaire du thème, structure \zh{因为…所以…}, \zh{应该…}.
\end{enumerate}
\end{corrige}

\newpage

\coursec{Exercices}

\courssub{Je vérifie mes connaissances}

\begin{exercice}[QCM : le bon sens du mot]
Pour chaque phrase, choisis la bonne réponse (a, b ou c).

1. \zh{腐败} (fǔbài) signifie :\\
a) la justice \quad b) la corruption \quad c) l'honnêteté

2. \zh{受贿} (shòuhuì) signifie :\\
a) accepter des pots-de-vin \quad b) donner un cadeau d'anniversaire \quad c) payer ses impôts

3. \zh{廉洁} (liánjié) signifie :\\
a) riche \quad b) intègre, honnête \quad c) célèbre

4. \zh{惩罚} (chéngfá) signifie :\\
a) récompenser \quad b) inviter \quad c) punir, sanctionner

5. \zh{贪污} (tānwū) signifie :\\
a) détourner de l'argent public \quad b) économiser de l'argent \quad c) prêter de l'argent
\end{exercice}

\begin{exercice}[Vrai ou faux ? Justifie ta réponse]
Dis si chaque affirmation est vraie (V) ou fausse (F), puis justifie en une phrase en français.

1. Le caractère \zh{贿} contient la clé \zh{贝} (bèi), liée à l'argent et aux richesses.

2. \zh{腐败} est un mot positif qui décrit un bon fonctionnaire.

3. La clé de \zh{贪} est aussi liée à l'argent (\zh{贝} en bas du caractère).

4. En chinois, « lutter contre la corruption » se dit \zh{反腐败} (fǎn fǔbài).

5. \zh{责任} (zérèn) signifie « liberté ».
\end{exercice}

\begin{exercice}[Vocabulaire : trouve le mot]
Trouve dans le vocabulaire de la leçon le mot chinois (caractères + pinyin) qui correspond à chaque définition française.

1. Un fonctionnaire, un agent de l'État.
2. L'argent donné illégalement pour obtenir un service.
3. Les intérêts, les droits du peuple.
4. La loi, la législation.
5. Porter atteinte à, nuire à (quelqu'un ou quelque chose).
\end{exercice}

\begin{exercice}[Associe caractères et pinyin]
Relie chaque caractère ou mot à son pinyin.

\begin{center}
\begin{tabularx}{\linewidth}{|X|X|}
\hline
\rowcolor{popblueL} Caractères & Pinyin \\
\hline
1. 腐败 & a) liánjié \\
\hline
2. 廉洁 & b) chéngfá \\
\hline
3. 贿赂 & c) fǔbài \\
\hline
4. 惩罚 & d) zérèn \\
\hline
5. 责任 & e) huìlù \\
\hline
\end{tabularx}
\end{center}
\end{exercice}

\courssub{Je m'entraîne}

\begin{exercice}[Traduis en français]
Traduis les phrases suivantes en français.

1. \zh{腐败是一个严重的问题。} (Fǔbài shì yí ge yánzhòng de wèntí.)

2. \zh{那个官员受贿了。} (Nà ge guānyuán shòuhuì le.)

3. \zh{政府应该惩罚贪污的人。} (Zhèngfǔ yīnggāi chéngfá tānwū de rén.)

4. \zh{我们要做廉洁的人。} (Wǒmen yào zuò liánjié de rén.)
\end{exercice}

\begin{exercice}[Texte à trous : banque de mots]
Complète le texte avec les mots de la banque : \zh{腐败}、\zh{贪污}、\zh{受贿}、\zh{廉洁}、\zh{惩罚}、\zh{责任}.

\begin{textebox}[Texte à compléter]
在一些国家，\_\_\_\_\_\_\_\_ 是一个大问题。有的官员 \_\_\_\_\_\_\_\_ 公款，有的官员 \_\_\_\_\_\_\_\_，收了别人的钱。这些行为损害了人民的利益。政府应该 \_\_\_\_\_\_\_\_ 他们。我们每个人都有 \_\_\_\_\_\_\_\_ 反对腐败，做一个 \_\_\_\_\_\_\_\_ 的人。\\
\emph{Zài yìxiē guójiā, \_\_\_\_ shì yí ge dà wèntí. Yǒude guānyuán \_\_\_\_ gōngkuǎn, yǒude guānyuán \_\_\_\_, shōu le biérén de qián. Zhèxiē xíngwéi sǔnhài le rénmín de lìyì. Zhèngfǔ yīnggāi \_\_\_\_ tāmen. Wǒmen měi ge rén dōu yǒu \_\_\_\_ fǎnduì fǔbài, zuò yí ge \_\_\_\_ de rén.}\\
« Dans certains pays, \_\_\_\_ est un grand problème. Certains fonctionnaires \_\_\_\_ l'argent public, d'autres \_\_\_\_ (des pots-de-vin). Ces comportements nuisent aux intérêts du peuple. Le gouvernement doit les \_\_\_\_. Chacun de nous a la \_\_\_\_ de lutter contre la corruption et d'être une personne \_\_\_\_. »
\end{textebox}
\end{exercice}

\begin{exercice}[Remets les mots dans l'ordre]
Remets les mots dans le bon ordre pour former une phrase correcte, puis traduis-la en français.

1. 问题 / 腐败 / 是 / 严重的 / 一个

2. 应该 / 政府 / 官员 / 惩罚 / 贪污的

3. 人民 / 的 / 损害 / 利益 / 了 / 腐败

4. 拒绝 / 我们 / 贿赂 / 应该
\end{exercice}

\begin{exercice}[Appariement : début et fin de phrase]
Relie chaque début de phrase (1-5) à sa fin (a-e), puis traduis les phrases complètes en français.

\begin{center}
\begin{tabularx}{\linewidth}{|X|X|}
\hline
\rowcolor{popblueL} Début & Fin \\
\hline
1. 贪污公款的人 & a) 是我们的责任。 \\
\hline
2. 反腐败 & b) 应该受到惩罚。 \\
\hline
3. 那个官员因为受贿 & c) 损害了国家的利益。 \\
\hline
4. 腐败行为 & d) 被警察抓了。 \\
\hline
5. 做一个廉洁的人 & e) 很重要。 \\
\hline
\end{tabularx}
\end{center}
\end{exercice}

\courssub{Je me prépare au Bac}

\begin{exercice}[Type Bac — Compréhension écrite]
Lis le texte suivant, puis réponds aux questions.

\begin{textebox}[中国的反腐败斗争]
近年来，中国政府非常重视反腐败工作。腐败损害人民的利益，也损害国家的发展。有些官员贪污公款，有些官员受贿，收了不该收的钱。政府用法律惩罚这些人，很多腐败的官员被抓了。政府还教育年轻人要诚实、廉洁，拒绝贿赂。老百姓支持政府的反腐败工作，因为他们希望社会更公平。反腐败不是一件容易的事，需要每个人的努力。只有大家一起努力，社会才能变得更好。\\
\emph{Jìnnián lái, Zhōngguó zhèngfǔ fēicháng zhòngshì fǎn fǔbài gōngzuò. Fǔbài sǔnhài rénmín de lìyì, yě sǔnhài guójiā de fāzhǎn. Yǒuxiē guānyuán tānwū gōngkuǎn, yǒuxiē guānyuán shòuhuì, shōu le bù gāi shōu de qián. Zhèngfǔ yòng fǎlǜ chéngfá zhèxiē rén, hěn duō fǔbài de guānyuán bèi zhuā le. Zhèngfǔ hái jiàoyù niánqīngrén yào chéngshí, liánjié, jùjué huìlù. Lǎobǎixìng zhīchí zhèngfǔ de fǎn fǔbài gōngzuò, yīnwèi tāmen xīwàng shèhuì gèng gōngpíng. Fǎn fǔbài bú shì yí jiàn róngyì de shì, xūyào měi ge rén de nǔlì. Zhǐyǒu dàjiā yìqǐ nǔlì, shèhuì cái néng biàn de gèng hǎo.}\\
« Ces dernières années, le gouvernement chinois accorde une grande importance à la lutte contre la corruption. La corruption nuit aux intérêts du peuple et nuit aussi au développement du pays. Certains fonctionnaires détournent l'argent public, d'autres acceptent des pots-de-vin, de l'argent qu'ils ne devraient pas recevoir. Le gouvernement punit ces personnes par la loi, et beaucoup de fonctionnaires corrompus ont été arrêtés. Le gouvernement éduque aussi les jeunes à être honnêtes et intègres, et à refuser les pots-de-vin. Les citoyens soutiennent la lutte anticorruption du gouvernement, car ils espèrent une société plus juste. Lutter contre la corruption n'est pas une chose facile : cela demande les efforts de chacun. Ce n'est qu'en travaillant tous ensemble que la société peut devenir meilleure. »
\end{textebox}

\textbf{A. Vrai ou faux ?} Recopie V (vrai) ou F (faux) et justifie en citant une phrase du texte en chinois.

1. 中国政府不重视反腐败工作。

2. 腐败只损害人民的利益，不损害国家的发展。

3. 老百姓支持政府的反腐败工作。

\textbf{B. Questions ouvertes.} Réponds en chinois, avec des phrases complètes.

4. 政府用什么惩罚腐败的官员？

5. 根据课文，反腐败需要什么？
\end{exercice}

\begin{exercice}[Type Bac — Traduction : thème et version]
\textbf{A. Version.} Traduis en français la phrase suivante :

\zh{有些官员贪污受贿，损害了人民的利益。政府应该惩罚他们。}\\
\emph{Yǒuxiē guānyuán tānwū shòuhuì, sǔnhài le rénmín de lìyì. Zhèngfǔ yīnggāi chéngfá tāmen.}

\textbf{B. Thème.} Traduis en chinois (caractères + pinyin) les phrases suivantes :

1. Nous devons refuser les pots-de-vin et être honnêtes.

2. La corruption est un problème grave dans beaucoup de pays.

3. Les jeunes ont la responsabilité de lutter contre la corruption.
\end{exercice}

\begin{exercice}[Type Bac — Expression écrite]
\textbf{Sujet :} \zh{年轻人在反腐败中的作用} (Le rôle des jeunes dans la lutte contre la corruption au Cameroun).

Rédige un texte de \textbf{5 à 8 phrases en chinois} (environ 80 à 120 caractères) en suivant la trame vue en classe :

\begin{enumerate}
\item Phrase d'introduction : la corruption est un problème (utilise \zh{腐败是……的问题}).
\item Ce que les jeunes doivent faire : refuser les pots-de-vin (\zh{拒绝贿赂}), être honnêtes (\zh{诚实}、\zh{廉洁}).
\item Ce que les jeunes peuvent faire à l'école et dans la société.
\item Phrase de conclusion : l'avenir du pays (\zh{国家的未来}).
\end{enumerate}

\textbf{Écriture des caractères (partie notée) :} recopie soigneusement les caractères \zh{腐}、\zh{败}、\zh{贪}、\zh{贿}、\zh{廉} en respectant l'ordre des traits, puis indique pour chacun sa clé (radical).
\end{exercice}

\newpage

\coursec{Corrigés des exercices}

\coursec{Corrigés des exercices de la leçon 27 — Vocabulaire et compréhension : corruption}

\begin{corrige}[Exercice 1 — QCM : le bon sens du mot]
1. \textbf{b)} \zh{腐败} (fǔbài) = la corruption. \zh{腐} évoque la pourriture, \zh{败} la défaite : image de la « pourriture » morale.

2. \textbf{a)} \zh{受贿} (shòuhuì) = accepter des pots-de-vin. \zh{受} = recevoir, subir ; \zh{贿} = pot-de-vin.

3. \textbf{b)} \zh{廉洁} (liánjié) = intègre, honnête. C'est l'antonyme moral de \zh{腐败}.

4. \textbf{c)} \zh{惩罚} (chéngfá) = punir, sanctionner. Attention à ne pas confondre avec \zh{奖励} (jiǎnglì, récompenser).

5. \textbf{a)} \zh{贪污} (tānwū) = détourner de l'argent public. \zh{贪} = avide, cupide ; \zh{污} = souillé, malhonnête.
\end{corrige}

\begin{corrige}[Exercice 2 — Vrai ou faux ? Justifie ta réponse]
1. \textbf{V.} \zh{贿} contient bien la clé \zh{贝} (bèi, le coquillage, ancienne monnaie), placée à gauche ; cette clé se retrouve dans de nombreux caractères liés à l'argent : \zh{财} (cái, richesse), \zh{贵} (guì, cher), \zh{贫} (pín, pauvre).

2. \textbf{F.} \zh{腐败} est un mot très négatif : il désigne la corruption, un comportement condamnable ; un bon fonctionnaire est \zh{廉洁} (liánjié, intègre).

3. \textbf{V.} \zh{贪} (tān, cupide) contient \zh{贝} en bas du caractère ; la partie haute \zh{今} (jīn) donne le son. La cupidité est donc liée à l'argent jusque dans l'écriture.

4. \textbf{V.} \zh{反} (fǎn) = s'opposer à, anti- ; \zh{反腐败} (fǎn fǔbài) = lutter contre la corruption. On trouve aussi \zh{反对腐败} (fǎnduì fǔbài).

5. \textbf{F.} \zh{责任} (zérèn) signifie « responsabilité », non « liberté » ; « liberté » se dit \zh{自由} (zìyóu).
\end{corrige}

\begin{corrige}[Exercice 3 — Vocabulaire : trouve le mot]
1. Un fonctionnaire : \zh{官员} (guānyuán).

2. L'argent donné illégalement : \zh{贿赂} (huìlù) ; pour le verbe « accepter des pots-de-vin » : \zh{受贿} (shòuhuì).

3. Les intérêts du peuple : \zh{人民的利益} (rénmín de lìyì) ; le mot cherché est \zh{利益} (lìyì).

4. La loi : \zh{法律} (fǎlǜ).

5. Porter atteinte à, nuire à : \zh{损害} (sǔnhài).
\end{corrige}

\begin{corrige}[Exercice 4 — Associe caractères et pinyin]
\begin{center}
\begin{tabularx}{\linewidth}{|X|X|X|}\hline
\rowcolor{popblueL} Caractères & Pinyin & Tons \\ \hline
\zh{腐败} & fǔbài & 3\textsuperscript{e} + 4\textsuperscript{e} \\ \hline
\zh{廉洁} & liánjié & 2\textsuperscript{e} + 2\textsuperscript{e} \\ \hline
\zh{贿赂} & huìlù & 4\textsuperscript{e} + 4\textsuperscript{e} \\ \hline
\zh{惩罚} & chéngfá & 2\textsuperscript{e} + 2\textsuperscript{e} \\ \hline
\zh{责任} & zérèn & 2\textsuperscript{e} + 4\textsuperscript{e} \\ \hline
\end{tabularx}
\end{center}

\textbf{Point de vigilance :} ne pas confondre \zh{贿赂} (huìlù, le pot-de-vin, nom) et \zh{受贿} (shòuhuì, accepter un pot-de-vin, verbe). Le caractère \zh{赂} contient aussi la clé \zh{贝} de l'argent.
\end{corrige}

\begin{corrige}[Exercice 5 — Traduis en français]
1. \zh{腐败是一个严重的问题。} (Fǔbài shì yí ge yánzhòng de wèntí.)\\
« La corruption est un problème grave. »\\
Structure : Sujet + 是 + (un + classificateur 个 + adjectif + 的 + nom). \zh{严重} (yánzhòng) = grave, sérieux.

2. \zh{那个官员受贿了。} (Nà ge guānyuán shòuhuì le.)\\
« Ce fonctionnaire a accepté des pots-de-vin. »\\
Le \zh{了} final marque l'accomplissement de l'action (passé). \zh{那个} (nà ge) = ce/cette + classificateur.

3. \zh{政府应该惩罚贪污的人。} (Zhèngfǔ yīnggāi chéngfá tānwū de rén.)\\
« Le gouvernement doit punir les personnes qui détournent (de l'argent public). »\\
\zh{应该} (yīnggāi) = devoir (conseil moral) ; \zh{贪污的人} = « les gens qui détournent » : le verbe qualifie le nom grâce à \zh{的}.

4. \zh{我们要做廉洁的人。} (Wǒmen yào zuò liánjié de rén.)\\
« Nous devons être des personnes intègres. »\\
\zh{要} (yào) = devoir, vouloir ; \zh{做……的人} = « être quelqu'un de … ».
\end{corrige}

\begin{corrige}[Exercice 6 — Texte à trous : banque de mots]
Texte complété :

\begin{textebox}[Texte à trous complété]
在一些国家，\textbf{腐败} 是一个大问题。有的官员 \textbf{贪污} 公款，有的官员 \textbf{受贿}，收了别人的钱。这些行为损害了人民的利益。政府应该 \textbf{惩罚} 他们。我们每个人都有 \textbf{责任} 反对腐败，做一个 \textbf{廉洁} 的人。
\end{textebox}

\emph{Zài yìxiē guójiā, fǔbài shì yí ge dà wèntí. Yǒude guānyuán tānwū gōngkuǎn, yǒude guānyuán shòuhuì, shōu le biérén de qián. Zhèxiē xíngwéi sǔnhài le rénmín de lìyì. Zhèngfǔ yīnggāi chéngfá tāmen. Wǒmen měi ge rén dōu yǒu zérèn fǎnduì fǔbài, zuò yí ge liánjié de rén.}

« Dans certains pays, la corruption est un grand problème. Certains fonctionnaires détournent l'argent public, d'autres acceptent des pots-de-vin. Ces comportements nuisent aux intérêts du peuple. Le gouvernement doit les punir. Chacun de nous a la responsabilité de lutter contre la corruption et d'être une personne intègre. »

\textbf{Justifications :}
\begin{itemize}
\item Trou 1 : il faut un nom-sujet général du texte → \zh{腐败}.
\item Trou 2 : devant \zh{公款} (gōngkuǎn, fonds publics), le verbe est \zh{贪污} (détourner).
\item Trou 3 : la suite « 收了别人的钱 » (a reçu l'argent des autres) impose \zh{受贿}.
\item Trou 4 : après \zh{政府应该}, il faut un verbe → \zh{惩罚}.
\item Trou 5 : \zh{有责任 + verbe} = « avoir la responsabilité de » → \zh{责任}.
\item Trou 6 : adjectif avant \zh{的人} → \zh{廉洁}.
\end{itemize}
\end{corrige}

\begin{corrige}[Exercice 7 — Remets les mots dans l'ordre]
1. \zh{腐败是一个严重的问题。} (Fǔbài shì yí ge yánzhòng de wèntí.)\\
« La corruption est un problème grave. »\\
Ordre : Sujet + 是 + 一个 + adjectif + 的 + nom. L'adjectif \zh{严重的} se place avant le nom \zh{问题}.

2. \zh{政府应该惩罚贪污的官员。} (Zhèngfǔ yīnggāi chéngfá tānwū de guānyuán.)\\
« Le gouvernement doit punir les fonctionnaires corrompus (qui détournent). »\\
Le complément d'objet \zh{贪污的官员} se place après le verbe \zh{惩罚}.

3. \zh{腐败损害了人民的利益。} (Fǔbài sǔnhài le rénmín de lìyì.)\\
« La corruption a nui aux intérêts du peuple. »\\
\zh{了} se place juste après le verbe \zh{损害} ; \zh{人民的利益} = « les intérêts du peuple » (possession avec 的).

4. \zh{我们应该拒绝贿赂。} (Wǒmen yīnggāi jùjué huìlù.)\\
« Nous devons refuser les pots-de-vin. »\\
Ordre : Sujet + modal 应该 + verbe 拒绝 + objet 贿赂.

\textbf{Rappel de méthode :} en chinois, l'ordre de base est Sujet – (modal) – Verbe – Objet ; le déterminant (的) relie le qualifiant au nom.
\end{corrige}

\begin{corrige}[Exercice 8 — Appariement : début et fin de phrase]
1 → \textbf{b)} \zh{贪污公款的人应该受到惩罚。} (Tānwū gōngkuǎn de rén yīnggāi shòudào chéngfá.)\\
« Les personnes qui détournent les fonds publics doivent être punies. » (\zh{受到} = recevoir, subir.)

2 → \textbf{a)} \zh{反腐败是我们的责任。} (Fǎn fǔbài shì wǒmen de zérèn.)\\
« Lutter contre la corruption est notre responsabilité. »

3 → \textbf{d)} \zh{那个官员因为受贿被警察抓了。} (Nà ge guānyuán yīnwèi shòuhuì bèi jǐngchá zhuā le.)\\
« Ce fonctionnaire a été arrêté par la police parce qu'il avait accepté des pots-de-vin. » (Passif avec \zh{被} + agent.)

4 → \textbf{c)} \zh{腐败行为损害了国家的利益。} (Fǔbài xíngwéi sǔnhài le guójiā de lìyì.)\\
« Les actes de corruption ont nui aux intérêts du pays. »

5 → \textbf{e)} \zh{做一个廉洁的人很重要。} (Zuò yí ge liánjié de rén hěn zhòngyào.)\\
« Être une personne intègre est très important. »
\end{corrige}

\begin{corrige}[Exercice 9 — Type Bac — Compréhension écrite]
\textbf{Barème indicatif :} A : 3 × 2 pts (1 pt réponse + 1 pt justification citée) ; B : 2 × 2 pts. Total : 10 pts.

\textbf{A. Vrai ou faux ?}

1. \textbf{F.} Le texte dit le contraire : \zh{中国政府非常重视反腐败工作。} (Zhōngguó zhèngfǔ fēicháng zhòngshì fǎn fǔbài gōngzuò.) — « Le gouvernement chinois accorde une grande importance à la lutte anticorruption. »

2. \textbf{F.} Le texte dit : \zh{腐败损害人民的利益，也损害国家的发展。} (Fǔbài sǔnhài rénmín de lìyì, yě sǔnhài guójiā de fāzhǎn.) — « La corruption nuit aux intérêts du peuple et nuit aussi au développement du pays. » Le mot \zh{也} (aussi) contredit le « seulement » de l'affirmation.

3. \textbf{V.} Le texte dit : \zh{老百姓支持政府的反腐败工作。} (Lǎobǎixìng zhīchí zhèngfǔ de fǎn fǔbài gōngzuò.) — « Les citoyens soutiennent la lutte anticorruption du gouvernement. »

\textbf{B. Questions ouvertes.}

4. \zh{政府用法律惩罚腐败的官员。} (Zhèngfǔ yòng fǎlǜ chéngfá fǔbài de guānyuán.)\\
« Le gouvernement punit les fonctionnaires corrompus par la loi. »\\
Structure : \zh{用 + outil/moyen + verbe} = « au moyen de ».

5. \zh{反腐败需要每个人的努力。} (Fǎn fǔbài xūyào měi ge rén de nǔlì.)\\
« Lutter contre la corruption demande les efforts de chacun. »\\
On peut ajouter : \zh{只有大家一起努力，社会才能变得更好。} (Zhǐyǒu dàjiā yìqǐ nǔlì, shèhuì cái néng biàn de gèng hǎo.) — « Ce n'est qu'en travaillant ensemble que la société peut devenir meilleure. »

\textbf{Points de vigilance :} citer la phrase du texte \emph{en chinois} pour justifier (exigence de l'épreuve) ; en B, répondre par une phrase complète avec sujet + verbe, pas par un mot isolé.
\end{corrige}

\begin{corrige}[Exercice 10 — Type Bac — Traduction : thème et version]
\textbf{Barème indicatif :} version 4 pts ; thème 3 × 2 pts. Total : 10 pts.

\textbf{A. Version.}\\
\zh{有些官员贪污受贿，损害了人民的利益。政府应该惩罚他们。}\\
« Certains fonctionnaires détournent des fonds et acceptent des pots-de-vin, ce qui nuit aux intérêts du peuple. Le gouvernement doit les punir. »\\
Remarque : \zh{贪污受贿} est une combinaison fréquente de deux verbes (« détourner et toucher des pots-de-vin ») ; \zh{他们} (tāmen) reprend « les fonctionnaires ».

\textbf{B. Thème.}

1. \zh{我们应该拒绝贿赂，做诚实的人。}\\
\emph{Wǒmen yīnggāi jùjué huìlù, zuò chéngshí de rén.}\\
« Nous devons refuser les pots-de-vin et être honnêtes. »

2. \zh{腐败在很多国家是一个严重的问题。}\\
\emph{Fǔbài zài hěn duō guójiā shì yí ge yánzhòng de wèntí.}\\
« La corruption est un problème grave dans beaucoup de pays. »\\
Attention : le complément de lieu \zh{在很多国家} se place avant le verbe \zh{是}.

3. \zh{年轻人有责任反对腐败。}\\
\emph{Niánqīngrén yǒu zérèn fǎnduì fǔbài.}\\
« Les jeunes ont la responsabilité de lutter contre la corruption. »

\textbf{Points de vigilance :} ne pas écrire « \zh{是严重的问题一个} » (le classificateur \zh{一个} précède l'adjectif) ; \zh{有责任 + verbe} sans \zh{的} ; bien placer \zh{在…} avant le verbe.
\end{corrige}

\begin{corrige}[Exercice 11 — Type Bac — Expression écrite]
\textbf{Barème indicatif :} respect du sujet et de la trame 4 pts ; correction grammaticale (ordre des mots, 的, 了) 4 pts ; richesse du vocabulaire du thème 2 pts ; longueur (5-8 phrases, 80-120 caractères) 2 pts ; écriture des caractères et clés 3 pts. Total : 15 pts.

\textbf{Proposition de rédaction modèle :}

\begin{textebox}[Rédaction modèle]
腐败是喀麦隆的一个严重的问题。它损害人民的利益，也损害国家的发展。我们年轻人应该拒绝贿赂，做诚实、廉洁的人。在学校里，我们不应该作弊；在社会上，我们应该批评腐败的行为。反腐败是每个人的责任。只有年轻人一起努力，国家的未来才能更美好。
\end{textebox}

\emph{Fǔbài shì Kāmàilóng de yí ge yánzhòng de wèntí. Tā sǔnhài rénmín de lìyì, yě sǔnhài guójiā de fāzhǎn. Wǒmen niánqīngrén yīnggāi jùjué huìlù, zuò chéngshí, liánjié de rén. Zài xuéxiào lǐ, wǒmen bù yīnggāi zuòbì; zài shèhuì shàng, wǒmen yīnggāi pīpíng fǔbài de xíngwéi. Fǎn fǔbài shì měi ge rén de zérèn. Zhǐyǒu niánqīngrén yìqǐ nǔlì, guójiā de wèilái cái néng gèng měihǎo.}

« La corruption est un problème grave au Cameroun. Elle nuit aux intérêts du peuple et au développement du pays. Nous, les jeunes, devons refuser les pots-de-vin et être des personnes honnêtes et intègres. À l'école, nous ne devons pas tricher ; dans la société, nous devons dénoncer les comportements corrompus. Lutter contre la corruption est la responsabilité de chacun. Ce n'est que si les jeunes travaillent ensemble que l'avenir du pays pourra être plus beau. »

\textbf{Écriture des caractères — clés et ordre des traits :}
\begin{itemize}
\item \zh{腐} (fǔ) : clé \zh{月} (yuè, radical « viande/charcuterie », écrit à gauche). Structure gauche-droite. On écrit d'abord la clé \zh{月} (traits : oblique, horizontale, horizontale, verticale), puis la partie droite \zh{府} (fǔ) : enveloppe \zh{广} (point, horizontale, crochet), puis \zh{付} à l'intérieur.
\item \zh{败} (bài) : clé \zh{贝} (bèi, l'argent). Structure gauche-droite. D'abord \zh{贝} à gauche, puis \zh{攵} (la clé « frapper ») à droite — de gauche à droite.
\item \zh{贪} (tān) : clé \zh{贝} (bèi, en bas). Structure haut-bas. D'abord \zh{今} (jīn) en haut, puis \zh{贝} en bas — de haut en bas.
\item \zh{贿} (huì) : clé \zh{贝} (bèi, à gauche). Structure gauche-droite. D'abord \zh{贝} à gauche, puis \zh{有} (yǒu) à droite — de gauche à droite.
\item \zh{廉} (lián) : clé \zh{广} (guǎng, le toit/abri, en haut). Structure semi-encadrée. D'abord l'enveloppe \zh{广} (point, horizontale, crochet), puis \zh{兼} (jiān) à l'intérieur — de l'extérieur vers l'intérieur.
\end{itemize}

\textbf{Points de vigilance :}
\begin{itemize}
\item Bien employer les mots du thème : \zh{腐败}、\zh{拒绝贿赂}、\zh{廉洁}、\zh{责任}、\zh{国家的未来}.
\item Structure de conclusion vue en classe : \zh{只有……才……} (zhǐyǒu … cái …) = « ce n'est que … que … ».
\item Ne pas oublier \zh{的} entre le qualifiant et le nom : \zh{严重的问题}, \zh{腐败的行为}.
\item Complément de lieu avant le verbe : \zh{在学校里} + verbe, jamais après.
\end{itemize}
\end{corrige}

\newpage

\coursec{Fiche bilan}

\begin{popfigure}
\begin{tikzpicture}[
  mindmap,
  grow cyclic,
  every node/.style={concept, execute at begin node=\hskip0pt},
  concept color=popblue,
  level 1/.append style={level distance=4.5cm, sibling angle=360/7},
  level 2/.append style={level distance=3cm, sibling angle=360/4},
  text=white,
  font=\small\bfseries,
  popblue/.style={concept color=popblue},
  poporange/.style={concept color=poporange},
  popgreen/.style={concept color=popgreen},
  poppurple/.style={concept color=poppurple},
  poppink/.style={concept color=poppink},
  popgold/.style={concept color=popgold},
  popdark/.style={concept color=popdark},
]
\node[concept, scale=1.2, align=center]{Leçon 27\\Corruption}
  child[poporange] { node[concept, align=center]{Vocabulaire\\clé} }
  child[popgreen] { node[concept, align=center]{Texte support\\« Lutte anticorruption »} }
  child[poppurple] { node[concept, align=center]{Radicaux\\Écriture} }
  child[poppink] { node[concept, align=center]{Compréhension\\QCM / Ouvertes} }
  child[popgold] { node[concept, align=center]{Structures\\Grammaire} }
  child[popdark] { node[concept, align=center]{Contexte\\Cameroun / Chine} }
  child[popblue] { node[concept, align=center]{Production\\Débat / Rédaction} };
\end{tikzpicture}
\legende{Carte mentale de la leçon 27 : Vocabulaire et compréhension — Corruption}
\end{popfigure}

\begin{bilan}

\end{bilan}

\courssub{Lexique clé (HSK 3--4) — Thème : Corruption \& Citoyenneté}

\begin{center}
\begin{tabularx}{\linewidth}{|X|X|X|X|}
\hline
\rowcolor{popblueL} \textbf{Caractères} & \textbf{Pinyin} & \textbf{Nature} & \textbf{Traduction} \\ \hline
腐败 & fǔbài & n./adj. & corruption ; corrompu \\ \hline
贪污 & tānwū & v. & détourner des fonds, péculation \\ \hline
受贿 & shòuhuì & v. & accepter des pots-de-vin \\ \hline
行贿 & xínghuì & v. & verser des pots-de-vin \\ \hline
贿赂 & huìlù & n. & pot-de-vin, corruption (acte) \\ \hline
反腐 & fǎnfǔ & v./n. & lutte anticorruption (abrév. de 反腐败) \\ \hline
廉政 & liánzhèng & n. & intégrité administrative, probité \\ \hline
监察 & jiānchá & v./n. & inspecter, contrôler ; inspection \\ \hline
纪律 & jìlǜ & n. & discipline (règles du Parti/État) \\ \hline
违纪 & wéijì & v. & violer la discipline \\ \hline
违法 & wéifǎ & v. & enfreindre la loi \\ \hline
查处 & cháchǔ & v. & enquêter et sanctionner \\ \hline
零容忍 & líng róngrěn & expr. & tolérance zéro \\ \hline
打虎拍蝇 & dǎ hǔ pāi yíng & idiome & « frapper les tigres et écraser les mouches » (grands et petits corruptos) \\ \hline
公款 & gōngkuǎn & n. & fonds publics \\ \hline
挪用 & nuóyòng & v. & détourner (fonds) à d'autres fins \\ \hline
举报 & jǔbào & v. & dénoncer, signaler \\ \hline
保护伞 & bǎohùsǎn & n. & parapluie protecteur (complicité hiérarchique) \\ \hline
阳光 & yángguāng & n./adj. & transparence (ex. 阳光政务) \\ \hline
诚信 & chéngxìn & n. & intégrité, honneur, crédit \\ \hline
\end{tabularx}
\end{center}

\courssub{Points de grammaire essentiels à maîtriser}

\begin{center}
\begin{tabularx}{\linewidth}{|X|X|X|}
\hline
\rowcolor{popblueL} \textbf{Structure / Règle} & \textbf{Exemple (caractères + pinyin)} & \textbf{Traduction} \\ \hline
\emph{被} (bèi) passif formel (agent souvent présent) & 该官员\textbf{被}纪委\textbf{调查}。\newline (Gāi guānyuán \textbf{bèi} jìwěi \textbf{diàochá}.) & Cet officiel \textbf{a été enquêté} par la Commission de discipline. \\ \hline
\emph{遭} (zāo) passif (nuance négative, malheur) & 群众\textbf{遭}贪官\textbf{欺压}。\newline (Qúnzhòng \textbf{zāo} tānguān \textbf{qīyā}.) & Les masses \textbf{subissent} l'oppression des officiels corrompus. \\ \hline
\emph{由于}... \emph{所以}... (Cause $\rightarrow$ Conséquence) & \textbf{由于}监管不力，\textbf{所以}腐败滋生。\newline (\textbf{Yóuyú} jiānguǎn bùlì, \textbf{suǒyǐ} fǔbài zīshēng.) & \textbf{En raison d'une} surveillance faible, \textbf{donc} la corruption prolifère. \\ \hline
\emph{导致} (dǎozhì) « entraîner, causer » (souvent négatif) & 贪污\textbf{导致}国家财产损失。\newline (Tānwū \textbf{dǎozhì} guójiā cáichǎn sǔnshī.) & La corruption \textbf{a entraîné} des pertes aux biens de l'État. \\ \hline
\emph{不仅}... \emph{而且}... (Progression) & 反腐\textbf{不仅}靠法律，\textbf{而且}靠舆论。\newline (Fǎnfǔ \textbf{bùjǐn} kào fǎlǜ, \textbf{érqiě} kào yúlùn.) & La lutte anticorruption \textbf{ne repose pas seulement} sur la loi, \textbf{mais aussi} sur l'opinion publique. \\ \hline
\emph{除非}... \emph{否则}... (Condition nécessaire) & \textbf{除非}制度透明，\textbf{否则}难根治腐败。\newline (\textbf{Chúfēi} zhìdù tòumíng, \textbf{fǒuzé} nán gēnzhì fǔbài.) & \textbf{À moins que} le système soit transparent, \textbf{sinon} difficile d'éradiquer la corruption. \\ \hline
\emph{把} (bǎ) + objet + V + 结果/方向 (Disposition) & 纪委\textbf{把}贪官\textbf{查处了}。\newline (Jìwěi \textbf{bǎ} tānguān \textbf{cháchǔ le}.) & La Commission \textbf{a enquêté et sanctionné} l'officiel corrompu. \\ \hline
\end{tabularx}
\end{center}

\courssub{Méthodes express}

\begin{enumerate}
\item \textbf{Lecture active :} Repérer les mots-clés (腐败, 贪污, 反腐, 纪委) $\rightarrow$ Identifier la structure Cause/Effet (由于/所以, 导致) $\rightarrow$ Noter les passifs (被, 遭) pour trouver l'agent et la victime.
\item \textbf{Questions de compréhension :} 
\begin{itemize}
\item \emph{Global :} Quel est le thème ? (Lutte anticorruption / Cas précis).
\item \emph{Détail :} Qui ? Quoi ? Quand ? Combien ? (Chiffres, noms).
\item \emph{Inférence :} Quelle mesure propose l'auteur ? (制度, 监察, 阳光政务).
\end{itemize}
\item \textbf{Traduction thème/version :} Respecter l'ordre chinois (Sujet + Temps + Lieu + Manière + Verbe). Attention aux passifs : « 被 + agent » $\rightarrow$ « par + agent » en français.
\item \textbf{Production orale/écrite :} Utiliser la structure « 不仅…而且… » pour argumenter ; citer l'idiome « 打虎拍蝇 » pour montrer la maîtrise culturelle.
\end{enumerate}



\begin{aretenir}[Checklist avant le Bac -- Leçon 27]
$\square$ Je connais par cœur les 20 mots du tableau de vocabulaire (caractères, pinyin, sens).\\
$\square$ J'écris correctement les caractères clés : 腐, 贪, 贿, 廉, 监, 纪, 挪, 举。\\
$\square$ J'identifie les radicaux : \zh{氵} (eau) dans 贪/腐/挪, \zh{贝} (coquillage/argent) dans 贿/贪/贪, \zh{广} (maison) dans 底/廊 (ex. 廉).\\
$\square$ Je maîtrise l'ordre des traits : haut $\rightarrow$ bas, gauche $\rightarrow$ droite, horizontal avant vertical, trait central avant ailes.\\
$\square$ Je distingue \textbf{被} (neutre/formel) et \textbf{遭} (négatif) à l'oral et à l'écrit.\\
$\square$ Je construis des phrases avec \textbf{由于}... \textbf{所以}... et \textbf{导致} sans confusion.\\
$\square$ J'utilise \textbf{不仅}... \textbf{而且}... pour structurer un argumentaire écrit.\\
$\square$ Je sais expliquer l'idiome \zh{打虎拍蝇} (dǎ hǔ pāi yíng) et son sens politique.\\
$\square$ Je cite au moins 2 différences / similitudes Cameroun--Chine (ex. CONAC / Commission nationale anti-corruption vs 纪委/监委).\\
$\square$ Je suis capable de rédiger un paragraphe de 80--100 caractères sur « Comment lutter contre la corruption ? » avec connecteurs logiques.\\
$\square$ Je réponds aux questions de compréhension (choix multiples + questions ouvertes) en justifiant par le texte.
\end{aretenir}

\findecours

\end{document}
